\documentclass[11pt]{article}

\usepackage[margin=1in]{geometry}
\usepackage{amsmath,amssymb,amsthm,mathtools}
\usepackage{mathrsfs}
\usepackage{enumitem}
\usepackage{booktabs}
\usepackage{array}
\usepackage[dvipsnames]{xcolor}
\usepackage[most]{tcolorbox}
\usepackage{microtype}
\usepackage[hidelinks]{hyperref}
\usepackage[nameinlink]{cleveref}

\allowdisplaybreaks
\setlength{\parskip}{0.4em}

%==================== macros ====================
\newcommand{\E}{\mathbb{E}}
\newcommand{\C}{\mathbb{C}}
\newcommand{\R}{\mathbb{R}}
\newcommand{\Z}{\mathbb{Z}}
\newcommand{\Lcal}{\mathcal{L}}
\newcommand{\Om}{\Omega}
\newcommand{\UQ}{U_{3^r}}
\newcommand{\Dcoll}{\mathfrak{D}_{\mathrm{Coll}}}
\DeclareMathOperator{\tr}{tr}
\newcommand{\KL}[2]{D\!\left(#1\,\middle\|\,#2\right)}

\newtcolorbox{principlebox}[1]{%
  enhanced, sharp corners, boxrule=0.8pt,
  colback=black!3, colframe=black!75,
  coltitle=black, fonttitle=\bfseries, title={#1},
  left=8pt, right=8pt, top=6pt, bottom=6pt,
  before skip=10pt, after skip=10pt}

\newtcolorbox{keystatement}{%
  enhanced, sharp corners, boxrule=0.6pt,
  colback=black!4, colframe=black!70,
  left=8pt, right=8pt, top=6pt, bottom=6pt,
  before skip=8pt, after skip=8pt}

\theoremstyle{plain}
\newtheorem{theorem}{Theorem}
\newtheorem{proposition}{Proposition}
\newtheorem{lemma}{Lemma}
\newtheorem{corollary}{Corollary}
\newtheorem{conjecture}{Conjecture}
\newtheorem{openproblem}{Open Problem}
\theoremstyle{definition}
\newtheorem{definition}{Definition}
\theoremstyle{remark}
\newtheorem*{remark}{Remark}
\newtheorem*{convention}{Convention}

\crefname{conjecture}{Conjecture}{Conjectures}
\Crefname{conjecture}{Conjecture}{Conjectures}
\crefname{definition}{Definition}{Definitions}
\Crefname{definition}{Definition}{Definitions}

\title{Finite Spectral Shadows for the Collatz Valuation Cocycle}
\author{John A. Janik\\\texttt{john.janik@gmail.com}}
\date{\today}

\begin{document}
\maketitle

\begin{abstract}
We study residue equidistribution for the accelerated Collatz/Syracuse map
$S(x)=(3x+1)/2^{v_2(3x+1)}$ through a pathwise twisted Birkhoff functional---the
\emph{deterministic defect cocycle} $W_n(x;\chi,s)$---and the finite
residue--height transfer operators $\Lcal_{s,Q,L,\chi}$ that govern it. Our aim
is not to prove the Collatz conjecture but to identify, rigorously, the correct
finite spectral object and the correct quotient on which its gap should be
sought. We prove two clean finite identities: a \emph{sliding-window congruence}
expressing $x_n\bmod 3^r$ as a finite window of the valuation word, and the fact
that the quadratic character $\chi_2$ of $(\Z/3^r)^\times$ satisfies
$\chi_2(S(x))=(-1)^{v_2(3x+1)}$, i.e.\ $\chi_2$ is \emph{valuation parity in
disguise} and is an exact coboundary of the faithful residue--valuation skew
product. Consequently the right form of the finite spectral hypothesis (``H23a'')
is a gap for all \emph{non-coboundary} characters of the \emph{faithful}
(synchronized) operator, not for all nontrivial characters. We then prove three
theorems about that operator: a \emph{classification of coboundaries}
($\operatorname{Cob}=\langle\chi_2\rangle$, for arbitrary positive branch
weights), whence by Wielandt's theorem a \emph{strict} spectral gap for every
non-coboundary character at every finite level and tilt; a \emph{conductor
collapse} (the nonzero spectrum of a conductor-$3^c$ sector is computed at level
$3^c$), which reduces uniformity of the gap in $Q$ to the boundedness of a
single sequence of computable constants $\gamma_c(s)$; and a \emph{rotation
identity} showing the quadratic sector of the killed height-strip operator is a
rigid rotation of the principal one. The uniformity question is then
\emph{resolved}: a threshold ($m=c$) Schur analysis of the adjoint kernel
yields bounded Plancherel shells for the renewal constants
($\tilde Q_L\le3C_1<5.16$), cascade row contraction at rate $(\sqrt3/2)^c$,
a level-cost law with the exact constant $\tfrac47<3^{-1/2}$, and a
root-ball analysis closing the degenerate budget---together the finite
H23a theorem $\sup_c\gamma_c(0)\le(\sqrt3/2)^{1/2}+o(1)<1$ at $s=0$. On
the global side we prove the \emph{keystone arrow} (every bad orbit casts
a finite non-coboundary two-point shadow), reduce the remaining quenched
exclusion to a joint $(2,3)$-adic max-plus rigidity, and prove two
unconditional exclusion theorems: periodic traps via linear forms in
logarithms, and---by a new \emph{repetition rigidity} argument in which
repeated valuation blocks are exact cycles---all bounded-discrepancy
valuation words of subexponential complexity, in particular every
Sturmian word. A counterexample must therefore mimic randomness at every
banded scale. The elementary core (the $2$-adic coding lemma and
repetition rigidity) is machine-verified in Lean~4. A companion
computational report~\cite{companion} documents the numerical suite.
\end{abstract}

\tableofcontents

\begin{convention}
From \cref{sec:identities} onward the residue space is specialised to
$\UQ=(\Z/3^r\Z)^\times$, $Q=3^r$, and all characters are multiplicative unless
stated otherwise. The additive-character experiments reported in the companion
are diagnostic only and are not the final operator vehicle. The character twist
in the transfer operator is always placed on the \emph{source} vertex $a$ (see
\cref{conv:twist}); with endpoint twisting the coboundary formulae are conjugate
but take a different form.
\end{convention}

\part{The object and the open inverse principle}

\section{The obstacle that fixes the shape of the object}

A single odd integer $x_0$ produces one deterministic valuation word
$w(x_0)=(b_0,b_1,\dots)$, $b_i=v_2(3x_i+1)\ge1$, with $x_{i+1}=S(x_i)$. The
annealed law is the pushforward of Haar measure on $\Z_2$ under
$x\mapsto v_2(3x+1)$, i.e.\ $\Pr(b)=2^{-b}$. ``The orbit obeys the annealed law''
asserts that the one prescribed word $w(x_0)$ equidistributes for the
residue--height skew product. No ergodic theorem applies to a single prescribed
non-generic point, so the missing object cannot be a measure-theoretic average.

It must instead be a \emph{dichotomy with a closure}: a pathwise quantity forced
either to decay (the orbit equidistributes) or to localise as a finite,
excludable obstruction. It must reduce, at the trivial character and after
ensemble averaging, to the annealed pressure and the Cram\'er tilt; and its
localisation must land in a fixed finite character sector.

\section{The deterministic defect cocycle}
\label{sec:cocycle}

For each modulus $Q$, height truncation $L$, tilt $s$, and character $\chi$ of
the residue space, define the pathwise twisted Birkhoff functional
\[
W_n(x;\chi,s)=\frac1n\sum_{i=0}^{n-1}\chi(x_i\bmod Q)\,2^{\,s\Delta_{b_i}},
\qquad \Delta_b=\log_2 3-b .
\]
The family $\{W_n(x;\chi,s)\}$ fibered over the finite levels is the
\emph{deterministic defect cocycle} $\Dcoll$. At $s=0$ it is the residue Weyl
sum, so $x_i\bmod Q$ equidistributes iff $W_n(x;\chi,0)\to0$ for all $\chi\ne1$.
At the trivial character the ensemble average reproduces the annealed pressure
$M(s)=\sum_{b\ge1}2^{-b}2^{s\Delta_b}=3^s/(2^{1+s}-1)$ and its Cram\'er root
$M(1)=1$ with tilted law $p^{(1)}_b=3\cdot4^{-b}$. Thus $\Dcoll$ is the quenched
lift of the thermodynamic formalism.

At finite level the intended content is, informally, an equivalence between
equidistribution failure, an approximate coboundary along the orbit, and a
near-peripheral eigenstate of $\Lcal_{s,Q,L,\chi}$. We use only the safe
direction.

\begin{remark}
The full three-way equivalence is too strong. A non-decaying twisted sum signals
correlation with $\chi$ but need not be an exact coboundary (it may reflect
concentration, subsequential trapping, or nonstationarity); and producing a
near-\emph{top} eigenstate is precisely the open keystone. The load-bearing
implication is: a persistent $W_n$ bias $\Rightarrow$ a nontrivial finite
character shadow, and (conjecturally) a near-peripheral state of
$\Lcal_{s,Q,L,\chi}$.
\end{remark}

\section{The inverse principle (conjectural)}

\begin{principlebox}{Quenched Inverse Principle (conjectural)}
Let $x_0$ have an infinite forward orbit and suppose, for some
\emph{non-coboundary} character $\chi$ on some $\UQ$, some tilt $s$, and some
$\delta>0$, that $\lvert W_n(x_0;\chi,s)\rvert\ge\delta$ along blocks of length
$\ell_j\to\infty$. Then there exist a finite pair $(Q,L)$ and a $\chi$-sector
eigenstate of $\Lcal_{s,Q,L,\chi}$ with $\lvert\lambda\rvert\ge
\rho(\Lcal_{s,Q,L,0})-o(1)$, exhibited by the orbit's empirical measure.
Conjoined with the non-coboundary gap
$\rho(\Lcal_{s,Q,L,\chi})\le(1-\kappa_Q)\rho(\Lcal_{s,Q,L,0})$, this is a
contradiction. Hence every orbit satisfying the hypotheses has \emph{no
persistent non-coboundary finite spectral shadow}.
\end{principlebox}

\begin{remark}[Scope of the conclusion]
This does not by itself give full equidistribution. It excludes persistent
non-coboundary finite spectral shadows. If, in addition, every bounded-deficit or
large-deviation failure of the annealed law produces such a shadow---the keystone
inverse arrow---then the relevant quenched obstruction is excluded. The keystone
is open and is morally a normality/disjointness statement for the $2$-adic word,
of Sarnak-hard flavour.
\end{remark}

\paragraph{The two leaks.} Closing the principle quantitatively requires (i) a
\emph{finiteness} input: a uniform gap $\kappa_Q$ that does not decay faster than
the entropy production of the bias (the subject of \cref{sec:target}); and (ii) a
\emph{circularity} input: reading a top eigenfunction off the orbit's own
empirical measure presumes near-stationarity, which is what one is proving. The
companion report shows the second leak is real at the single-orbit level.

\part{Rigorous finite identities}
\label{sec:identities}

These identities are the rigorous core: they bridge the valuation word and the
residue characters, and they pin down the unique structural degeneracy of the
operator.

\begin{convention}\label{conv:twist}
The transfer operator $\Lcal_{s,Q,L,\chi}$ has matrix element, on an edge
$(a,z)\xrightarrow{b}(a',z')$ with $a'\equiv(3a+1)2^{-b}$ and
$z'=z+\lfloor R(\log_2 3-b)\rceil$ (killed off $[0,L]$),
\[
\Lcal_{s,Q,L,\chi}[(a',z'),(a,z)]=2^{-b}2^{s\Delta_b}\,\chi(a),
\]
i.e.\ the character twist sits on the \emph{source} residue $a$.
\end{convention}

\section{The sliding-window congruence}

\begin{lemma}[Sliding-window congruence]\label{lem:cong}
With $B_n=\sum_{i<n}b_i$ and the exact affine iterate
$x_n=(3^n x_0+A_n)/2^{B_n}$, $A_n=\sum_{j=0}^{n-1}3^{\,n-1-j}2^{B_j}$, one has for
$n\ge r$
\[
x_n\equiv\sum_{k=1}^{r}3^{\,k-1}\,2^{-(b_{n-1}+\cdots+b_{n-k})}\pmod{3^r}.
\]
\end{lemma}
\begin{proof}
Modulo $3^r$ the term $3^n x_0$ vanishes for $n\ge r$. In $A_n$ the summand
indexed by $j$ carries $3^{\,n-1-j}$, which is nonzero mod $3^r$ only when
$n-1-j<r$, i.e.\ $j\in\{n-r,\dots,n-1\}$. Writing $j=n-k$ ($k=1,\dots,r$) gives
$3^{\,n-1-j}=3^{\,k-1}$, so $A_n\equiv\sum_{k=1}^r 3^{k-1}2^{B_{n-k}}\pmod{3^r}$.
Since $2$ is a unit mod $3^r$, multiply by $2^{-B_n}$:
$x_n\equiv\sum_{k=1}^r3^{k-1}2^{B_{n-k}-B_n}=\sum_{k=1}^r
3^{k-1}2^{-(b_{n-k}+\cdots+b_{n-1})}\pmod{3^r}$.
\end{proof}

\noindent Thus a finite $3^r$-residue character is a finite sliding-window
observable of the valuation process---one of the few rigorous bridges from
valuation word to residue shadow. Since $2$ is a primitive root mod $3^r$ (order
$\phi=2\cdot3^{r-1}$), the resonant frequencies attached to nontrivial characters
are $\xi=k/(2\cdot3^{r-1})$, and the single-symbol Fourier transform
$\hat p(\xi)=\sum_{b\ge1}2^{-b}e(b\xi)=e(\xi)/(2-e(\xi))$ obeys
$\lvert\hat p(\xi)\rvert=(5-4\cos2\pi\xi)^{-1/2}<1$ for $\xi\notin\Z$---the
annealed decorrelation mechanism. (The companion confirms the \emph{exact}
operator departs from this i.i.d.\ value.)

\section{The quadratic character is valuation parity}

\begin{lemma}[$\chi_2$ factors through reduction mod $3$]\label{lem:fac3}
$\UQ$ is cyclic of order $\phi=2\cdot3^{r-1}$, and its unique order-$2$ character
$\chi_2$ has $\ker\chi_2=\{u:\,u\equiv1\!\!\pmod3\}$. Equivalently
$\chi_2(a)=+1$ if $a\equiv1$ and $-1$ if $a\equiv2\pmod3$.
\end{lemma}
\begin{proof}
$2$ is a primitive root mod $3^r$, so the group is cyclic and $\ker\chi_2$ is its
unique index-$2$ subgroup (the squares), of order $3^{r-1}$. Reduction
$\UQ\to(\Z/3)^\times\cong\Z/2$ has kernel $\{u\equiv1\!\pmod3\}$, also of order
$3^{r-1}$; in a cyclic group the subgroup of a given order is unique, so the two
coincide.
\end{proof}

\begin{theorem}[Key identity]\label{thm:key}
For every odd $x$, with $b=v_2(3x+1)$,
\[
\chi_2\!\bigl(S(x)\bmod 3^r\bigr)=(-1)^{b}.
\]
Hence $\chi_2(x_{i+1})=(-1)^{b_i}$ along every orbit.
\end{theorem}
\begin{proof}
Work mod $3$: $2\equiv-1$ and $3x+1\equiv1$, so
$S(x)=(3x+1)2^{-b}\equiv(-1)^{b}\pmod3$. By \cref{lem:fac3},
$\chi_2(S(x))=(-1)^b$. (Since $3x+1\equiv1\pmod3$, $S(x)$ is a unit.)
\end{proof}

\begin{corollary}\label{cor:run}
$\prod_{i=1}^n\chi_2(x_i)=(-1)^{B_n}$: the accumulated quadratic character is
exactly the parity of the cumulative $2$-adic valuation. In particular the
residues are \emph{not} uniform mod $3$; they occur in ratio
$\Pr(x\equiv1)=\sum_{b\text{ even}}2^{-b}=\tfrac13$, $\Pr(x\equiv2)=\tfrac23$.
\end{corollary}

\section{The quadratic coboundary and the operator dichotomy}

Consider the synchronized residue--valuation skew product with state
$(a,\tau)\in\UQ\times(\Z/2)$ and dynamics
$(a,\tau)\xrightarrow{b}\bigl((3a+1)2^{-b},\,\tau+b\bigr)$ (one valuation $b$
drives both coordinates), and let $\Lcal_{\chi}$ be the corresponding twisted
annealed operator (\cref{conv:twist}), $\Lcal_{\mathbf1}$ the principal one.

\begin{theorem}[Quadratic coboundary]\label{thm:cobound}
The phase $d(a,\tau)=\chi_2(a)(-1)^{\tau}\in\{\pm1\}$ satisfies
$d(a',\tau')/d(a,\tau)=\chi_2(a)$ on every edge. Hence
$\Lcal_{\chi_2}=D\,\Lcal_{\mathbf1}\,D^{-1}$ with $D=\mathrm{diag}(d)$, so
$\rho(\Lcal_{\chi_2})=\rho(\Lcal_{\mathbf1})$: the $\chi_2$ sector has \emph{no
gap}, for every tilt $s$ and any height discretisation.
\end{theorem}
\begin{proof}
By \cref{thm:key} $\chi_2(a')=(-1)^b$, and $\tau'=\tau+b$ gives
$(-1)^{\tau'}=(-1)^\tau\chi_2(a')$. Thus
$d(a',\tau')=\chi_2(a')(-1)^{\tau'}=(-1)^\tau$ and
$d(a',\tau')/d(a,\tau)=(-1)^\tau/(\chi_2(a)(-1)^\tau)=\chi_2(a)$. With the source
twist, the conjugated matrix element is $d(a',\tau')^{-1}2^{-b}2^{s\Delta_b}
\chi_2(a)d(a,\tau)=2^{-b}2^{s\Delta_b}$, the principal element; $D$ is real
orthogonal, so the operators are isospectral.
\end{proof}

So $\chi_2$ is not a defect---it belongs to the height/valuation reference. This
forces the operator question.

\begin{proposition}[Residue-only gap; the $\Lcal$ dichotomy]\label{prop:dich}
Without the valuation-parity coordinate, $\chi_2$ is not a coboundary: on the
$\{\pm1\}$ (mod $3$) reduction the twisted $2\times2$ operator has eigenvalues $0$
and $\sum_b(-1)^bw_s(b)$, giving the closed-form residue-only gap ratio
\[
\frac{\rho(\Lcal^{\mathrm{res}}_{\chi_2})}{\rho(\Lcal^{\mathrm{res}}_{\mathbf1})}
=\frac{\bigl|\sum_b(-1)^bw_s(b)\bigr|}{\sum_bw_s(b)}
=\frac{2^{1+s}-1}{2^{1+s}+1}
\qquad\Bigl(\tfrac13\text{ at }s{=}0,\ \tfrac35\text{ at }s{=}1\Bigr),
\quad w_s(b)=2^{-b}2^{s\Delta_b}.
\]
\end{proposition}

\begin{remark}[Faithful \texttt{sync} vs.\ \texttt{factored}; why faithfulness is forced]
Two finite operators couple residue and height. The \emph{faithful}
(\texttt{sync}) operator uses one valuation $b$ for both, matching the orbit; the
\texttt{factored} operator $\Lcal_{\mathrm{res}}(\chi)\otimes T_{\mathrm{height}}$
draws the height's valuation independently. The factored operator reports a gap
at $\chi_2$, but \cref{cor:run} shows $W_n(x;\chi_2,0)\to-\tfrac13\ne0$: the
$\chi_2$ component does not decay. The factored operator therefore predicts decay
that does not occur. Because the inverse principle runs \emph{backwards} (a
persistent $W_n$ must force a peripheral eigenstate), the operator's spectral
radius must equal the true growth rate of $W_n$; only \texttt{sync} qualifies. We
adopt the faithful \texttt{sync} operator and use \texttt{factored}/residue-only
operators only as diagnostics.
\end{remark}

\section{Classification of coboundaries and the strict gap at every level}
\label{sec:class}

We now prove that $\chi_2$ is the \emph{only} nontrivial coboundary of the
synchronized skew product---the central finite theorem of this paper. Throughout
this section the operator is the synchronized residue--parity operator of
\cref{thm:cobound}: states $(a,\tau)\in\UQ\times(\Z/2)$, edges
$(a,\tau)\xrightarrow{b}(S_b(a),\tau+b)$ with $S_b(a)=(3a+1)2^{-b}$, source twist,
and \emph{arbitrary} positive summable branch weights $w_b>0$ (in particular
$w_b=2^{-b}2^{s\Delta_b}$ for any tilt $s$). We allow the twisted-conjugacy
(Wielandt) form with a global phase.

\begin{theorem}[Coboundary classification]\label{thm:class}
Let $\chi$ be a character of $\UQ$, and suppose there exist $\lambda\in S^1$ and
$d:\UQ\times(\Z/2)\to S^1$ with
\[
d\bigl(S_b(a),\tau+b\bigr)=\lambda\,\chi(a)\,d(a,\tau)
\qquad\text{for all }a\in\UQ,\ b\ge1,\ \tau\in\Z/2 .
\]
Then $\chi\in\{\mathbf1,\chi_2\}$. Conversely both occur ($\chi=\mathbf1$ with
$d\equiv1$; $\chi=\chi_2$ with $d(a,\tau)=\chi_2(a)(-1)^\tau$, $\lambda=1$, by
\cref{thm:cobound}).
\end{theorem}
\begin{proof}
\emph{Step 1 ($d$ factors through $(\chi_2,\tau)$).} Fix $(a,\tau)$ and compare
the edges with valuations $b$ and $b+2$: they have the same parity increment and
$S_{b+2}(a)=4^{-1}S_b(a)$, so the displayed identity gives
\[
d\bigl(4^{-1}S_b(a),\,\tau+b\bigr)=\lambda\chi(a)d(a,\tau)=d\bigl(S_b(a),\,\tau+b\bigr).
\]
For fixed $a$, as $b$ ranges over $\Z_{\ge1}$ the residue $S_b(a)$ ranges over the
coset $(3a+1)\langle2\rangle=\UQ$ (since $2$ is a primitive root mod $3^r$), and
for each fixed value $v=S_b(a)$ the residue class of $b$ mod $\phi$ is determined,
so letting $\tau$ vary gives both parities $\sigma=\tau+b$. Hence
$d(4^{-1}v,\sigma)=d(v,\sigma)$ for \emph{every} $(v,\sigma)$: $d(\cdot,\sigma)$
is constant on cosets of $\langle4\rangle$. Since $\UQ$ is cyclic,
$\langle4\rangle=\langle2^2\rangle$ is its unique index-$2$ subgroup, the squares
$=\ker\chi_2$ (\cref{lem:fac3}). Therefore $d(a,\tau)=D(\chi_2(a),\tau)$ for some
$D:\{\pm1\}\times\Z/2\to S^1$.

\emph{Step 2 ($\chi$ is trivial on the squares).} Substituting into the identity
and using $\chi_2(S_b(a))=(-1)^b$ (\cref{thm:key}),
\[
D\bigl((-1)^b,\,\tau+b\bigr)=\lambda\,\chi(a)\,D(\chi_2(a),\tau)
\qquad\text{for all }a,b,\tau .
\]
Fix $b$ and $\tau$: the left side does not depend on $a$, so
$a\mapsto\chi(a)D(\chi_2(a),\tau)$ is constant on $\UQ$. In particular $\chi$ is
constant on the subgroup $\ker\chi_2$ of squares; a character constant on a
subgroup equals $1$ there. Hence $\chi$ factors through
$\UQ/\ker\chi_2\cong\Z/2$, i.e.\ $\chi\in\{\mathbf1,\chi_2\}$.
\end{proof}

\begin{lemma}[Irreducibility]\label{lem:irred}
The principal synchronized operator $\Lcal_{\mathbf1}$ on $\UQ\times(\Z/2)$ is
irreducible (indeed primitive): every state is reachable from every state in
exactly two steps.
\end{lemma}
\begin{proof}
From any $a$, one step reaches any prescribed residue $v$ (choose $b$ with
$2^{-b}=v/(3a+1)$, possible since $2$ generates $\UQ$), and the parity increment
of that step is determined by $v$ alone: $b\equiv[\chi_2(v)=-1]\pmod 2$, because
$\chi_2(v)=(-1)^b$. Given a target $(u,\sigma)$, choose the intermediate residue
$v$ with $\chi_2(v)$ set so that the two forced parity increments sum to
$\sigma-\tau$; both square classes are nonempty, so this is always possible.
\end{proof}

\begin{corollary}[Strict spectral gap at every finite level]\label{cor:strict}
For every $r\ge2$, every tilt $s$, and every character
$\chi\notin\{\mathbf1,\chi_2\}$ of $\UQ$,
\[
\rho\bigl(\Lcal^{\mathrm{par}}_{s,Q,\chi}\bigr)<\rho\bigl(\Lcal^{\mathrm{par}}_{s,Q,\mathbf1}\bigr),
\]
where $\Lcal^{\mathrm{par}}$ denotes the synchronized residue--parity operator. In
particular $\kappa^{\mathrm{par}}_Q(s)>0$ for every fixed $Q=3^r$ and $s$.
\end{corollary}
\begin{proof}
$\lvert\Lcal_{\chi}\rvert=\Lcal_{\mathbf1}$ entrywise (the twist has unit
modulus), and $\Lcal_{\mathbf1}$ is irreducible (\cref{lem:irred}). By Wielandt's
equality theorem, $\rho(\Lcal_\chi)=\rho(\Lcal_{\mathbf1})$ would force
$\Lcal_\chi=e^{i\varphi}D\,\Lcal_{\mathbf1}D^{-1}$ with $D$ diagonal
unimodular---which is precisely the identity hypothesised in \cref{thm:class}
with $\lambda=e^{-i\varphi}$. By \cref{thm:class} this requires
$\chi\in\{\mathbf1,\chi_2\}$, a contradiction.
\end{proof}

\begin{remark}
The proof of \cref{thm:class} is purely structural: it uses only which edges
exist (the richness of the branch set), not the values of the weights. It
therefore holds simultaneously for all tilts and for any summable positive
weighting of the valuations, and it bypasses the principal-unit functional
equation entirely. This resolves the coboundary-classification question for the
synchronized skew product; see \cref{sec:strip} for the height-strip refinement.
\end{remark}

\section{Conductor autonomy and collapse}

\begin{proposition}[Conductor reduction is autonomous]\label{prop:cond}
A character of conductor $3^c$ depends only on $a\bmod3^c$. Since
$3a\equiv3(a\bmod3^{c-1})\pmod{3^c}$, the quantity $(3a+1)2^{-b}\bmod3^c$ is a
function of $(a\bmod3^c,b)$ alone; hence the mod-$3^c$ reduction of the dynamics
is autonomous, and the gap ratio of a conductor-$3^c$ sector is \emph{independent
of the ambient $r$} once $r\ge c$.
\end{proposition}
\begin{proof}
Immediate from $3a\equiv3(a\bmod3^{c-1})\pmod{3^c}$: the induced twisted operator
on $(\Z/3^c)^\times\times(\text{height})$ makes no reference to $r$.
\end{proof}

Autonomy gives an embedding of the level-$3^c$ spectrum into the level-$3^r$
operator. The following sharper statement says nothing else survives: the
\emph{entire} nonzero spectrum of a conductor-$3^c$ sector is computed at level
$3^c$.

\begin{theorem}[Conductor collapse]\label{thm:collapse}
Let $\chi$ have conductor $3^c$, $2\le c\le r$. Then for the synchronized
operator at level $Q=3^r$ (residue--parity, or the height-strip operator with or
without killing),
\[
\operatorname{spec}\bigl(\Lcal^{(r)}_{\chi}\bigr)\setminus\{0\}
=\operatorname{spec}\bigl(\Lcal^{(c)}_{\chi}\bigr)\setminus\{0\}.
\]
Consequently every gap ratio is exactly conductor-determined, and
\[
\kappa_{3^r}(s)=1-\max_{2\le c\le r}\gamma_c(s),
\qquad
\gamma_c(s)=\max_{\operatorname{cond}\chi=3^c}
\frac{\rho(\Lcal^{(c)}_{s,\chi})}{\rho(\Lcal^{(c)}_{s,\mathbf1})},
\]
a maximum of finitely many $r$-independent constants. In particular
$\kappa_{3^r}$ is non-increasing in $r$, and (with \cref{cor:strict})
$\kappa^{\mathrm{par}}_{3^r}(s)>0$ for every fixed $r$; uniform positivity as
$r\to\infty$ is equivalent to $\sup_c\gamma_c(s)<1$.
\end{theorem}
\begin{proof}
Work with the Koopman form $M_\chi g(a,\cdot)=\chi(a)\sum_b w_b\,g(S_b(a),\cdot)$,
where $(\cdot)$ stands for the auxiliary coordinate (parity $\tau$, or the strip
coordinate $z$ with or without killing); the argument below acts only on the
$a$-dependence and carries the auxiliary coordinate along unchanged. Here
$A(a)=3a+1$ and $S_b(a)=A(a)2^{-b}$. Every function of $(a,\cdot)$ expands
uniquely as $g=\sum_{\eta\in\widehat{U_{3^r}}}\eta\otimes g_\eta(\cdot)$; the
conductor of $\eta$ is the least $3^j$ such that $\eta$ factors through
$U_{3^j}$, and we set
\[
V_j=\Bigl\{\textstyle\sum_\eta\eta\otimes g_\eta:\ g_\eta=0\text{ unless }
\operatorname{cond}\eta\le3^j\Bigr\},
\qquad V_0\subset V_1\subset\cdots\subset V_r .
\]
For a single character $\eta$,
\[
M_\chi(\eta\otimes u)=\chi\cdot(\eta\circ A)\otimes(\text{transported }u),
\qquad A(a)=3a+1 .
\]
The function $\eta\circ A$ depends only on $a\bmod3^{e-1}$ where
$3^e=\operatorname{cond}\eta$ (since $3a\bmod3^{e}$ depends only on
$a\bmod3^{e-1}$), so its Fourier support lies in conductors $\le3^{e-1}$;
multiplying by $\chi$ moves support into conductors $\le3^{\max(c,e-1)}$. Hence
\[
M_\chi V_j\subseteq V_{\max(c,\,j-1)} .
\]
Thus $V_c$ is invariant, and on the quotient filtration
$V_c\subset V_{c+1}\subset\cdots\subset V_r$ the induced maps are strictly
lowering, so $M_\chi$ is nilpotent on $V_r/V_c$ (of index $\le r-c$). Therefore
$\operatorname{spec}(M_\chi)=\operatorname{spec}(M_\chi|_{V_c})\cup\{0\}$, and
$M_\chi|_{V_c}$ is the level-$3^c$ operator by \cref{prop:cond}. The principal
radii agree at all levels ($M_{\mathbf1}\mathbf1=\bigl(\sum_bw_b\bigr)\mathbf1$),
giving the ratio statement.
\end{proof}

\begin{remark}
\cref{thm:collapse} is confirmed exactly by computation: on the residue--parity
operator each $\gamma_c$ is identical across all ambient levels tested
($Q\le3^6$), e.g.\ $\gamma_2(0)=0.654654$ and $\gamma_3(0)=0.763817$ at every
$Q$. The same collapse holds for the killed height-strip operator (the proof
never touches the height index), which is why the strip conductor table of
\cref{sec:numerics} is $r$-independent.
\end{remark}

\section{Height-strip rigidity of the quadratic sector}
\label{sec:strip}

On the discretised height strip (integer scale $R\ge1$, increment
$z'=z+c^\ast-Rb$ with $c^\ast=\lfloor R\log_23\rceil$, truncation to $z\in[0,L]$
with killing) the quadratic sector is not merely peripheral---its spectrum is a
rigid rotation of the principal one.

\begin{proposition}[Rotation identity]\label{prop:rot}
Let $D=\operatorname{diag}\bigl(\chi_2(a)\,e^{i\pi z/R}\bigr)$. Then on the
height-strip operator, truncated or not, and for every tilt $s$,
\[
\Lcal_{\chi_2}=e^{-i\pi c^\ast/R}\,D\,\Lcal_{\mathbf1}\,D^{-1},
\qquad\text{hence}\qquad
\operatorname{spec}\bigl(\Lcal_{s,Q,L,\chi_2}\bigr)
=e^{-i\pi c^\ast/R}\operatorname{spec}\bigl(\Lcal_{s,Q,L,\mathbf1}\bigr).
\]
\end{proposition}
\begin{proof}
On an edge $(a,z)\to(a',z+c^\ast-Rb)$, using $\chi_2(a')=(-1)^b$
(\cref{thm:key}) and $e^{i\pi(c^\ast-Rb)/R}=e^{i\pi c^\ast/R}(-1)^b$,
\[
\frac{d(a',z')}{d(a,z)}
=\frac{\chi_2(a')e^{i\pi z'/R}}{\chi_2(a)e^{i\pi z/R}}
=\chi_2(a')\,(-1)^b\,e^{i\pi c^\ast/R}\,\chi_2(a)
=e^{i\pi c^\ast/R}\,\chi_2(a),
\]
since $\chi_2(a')(-1)^b=1$ and $\chi_2(a)^{-1}=\chi_2(a)$. This is the twisted
coboundary identity with $\lambda=e^{i\pi c^\ast/R}$; as $D$ is diagonal in
$(a,z)$ it commutes with the truncation to $z\in[0,L]$.
\end{proof}

\begin{keystatement}
$\chi_2$ is part of the baseline dynamics, not an obstruction to
equidistribution: its sector is a rigid rotation of the principal one, its Weyl
sum is slaved to the valuation statistics, and it is quotiented into the
reference law.
\end{keystatement}

\noindent This explains the machine-precision equality
$\rho(\Lcal_{\chi_2})=\rho(\Lcal_{\mathbf1})$ observed on the killed strip for
every $(R,L)$ tested. For characters \emph{outside} $\langle\chi_2\rangle$ no
product-form strip potential exists: a potential
$d(a,z)=\chi'(a)e^{i\omega z}$ satisfies the strip coboundary identity iff
$\chi'(2)e^{i\omega R}=1$ and $\chi'(3a+1)=\lambda'\chi(a)\chi'(a)$ on $\UQ$; an
exhaustive search over all character pairs for $r\le4$ finds only
$(\chi',\chi)=(\mathbf1,\mathbf1)$ and $(\chi_2,\chi_2)$. (The general---%
non-product-form---strip classification reduces to this principal-unit
functional equation and is left open; \cref{cor:strict} settles the
residue--parity case unconditionally.)

\part{Numerical evidence (summary)}
\label{sec:numerics}

The full suite \texttt{collatz\_tests} (ISO~C11; \textsf{FAST128}/\textsf{GMP}
backends; Arnoldi spectral solver) and Experiments A--I are documented in the
companion~\cite{companion}; every figure there is reproducible. We summarise only
what bears on the statements above.

\begin{enumerate}[leftmargin=1.6em]
\item \textbf{Congruence (\cref{lem:cong}).} Verified in
$23{,}905{,}968/23{,}905{,}968$ checks (seeds $3$--$10^6$, $r\le8$, native) and
$144/144$ in exact \textsc{gmp} arithmetic for a $16$-digit seed at $r=12$ with
$n$ up to $155$ (values $\gg2^{128}$).
\item \textbf{Key identity (\cref{thm:key}) and \cref{cor:run}.}
$\chi_2(S(x))=(-1)^b$ holds in $10^5/10^5$ tested seeds;
$W_n(x;\chi_2,0)$ has empirical mean $-0.353\approx-\tfrac13$ over
$3.2\times10^4$ orbits, confirming the residues are biased $\tfrac13:\tfrac23$.
\item \textbf{Annealed laws.} Height tails fit $\log_2\Pr(h\ge H)\approx-H$
(slope $\approx-0.97$); high-ascent excursion bulk leans to the Cram\'er tilt
$p^{(1)}$ ($\KL{\cdot}{p^{(1)}}=0.093\ll\KL{\cdot}{p}=0.191$, boundary-excluded).
\item \textbf{Operator gaps.} On the residue-only operator every nontrivial
sector is gapped; \cref{prop:dich} is reproduced to the digit.
\cref{thm:class,cor:strict} are confirmed: on the residue--parity operator, a
full-group scan finds exactly one peripheral sector ($k=\phi/2=\chi_2$) at every
level and tilt. \cref{thm:collapse} is confirmed \emph{exactly}: each $\gamma_c$
is identical across all ambient levels. The computed sequences:
\[
\begin{array}{c|ccccc}
\gamma_c\ (\text{residue--parity}) & c{=}2 & c{=}3 & c{=}4 & c{=}5 & c{=}6\\\hline
s=0 & 0.6547 & \mathbf{0.7638} & 0.7629 & 0.7202 & 0.7147\\
s=1 & 0.9078 & 0.8691 & \mathbf{0.9264} & 0.9210 & 0.9021
\end{array}
\]
\[
\begin{array}{c|cccc}
\gamma_c\ (\text{killed strip},\,s{=}0) & c{=}2 & c{=}3 & c{=}4 & c{=}5\\\hline
 & \mathbf{0.9069} & 0.8530 & 0.8261 & 0.7306
\end{array}
\]
Both sequences are bounded away from $1$ and peak at a \emph{fixed small}
conductor (boldface), decaying beyond it. Note the residue--parity sequence is
\emph{not} monotone in $c$ (e.g.\ $\gamma_4>\gamma_2$ at $s{=}1$), while the
killed strip is monotone over the tested range: the extremal conductor depends on
the operator, but bounded-conductor extremality is seen in both. See
\cref{sec:target}.
\item \textbf{Shadow $\Rightarrow$ gap, only by tilt.} In a $1.25\times10^6$-window
search, single-orbit ``shadows'' concentrate at the weakest-gap tilt $s=1$ ($82\%$
of flagged windows; $100\%$ of long ones) and in weaker-gap sectors, but the
per-sector signal is dominated by the smallest modulus---the circularity leak made
explicit.
\end{enumerate}

\begin{remark}[Provisional labels]
The conclusions involving the labels \texttt{critical\_C\_A}, \texttt{high\_ascent},
and \texttt{post\_exit} (Experiments B, D, I) are \emph{exploratory} until the
formal $C_A$-critical certificate is fixed: the literal criterion $V^2\le2^A P$ is
degenerate and the suite uses a working \texttt{height} definition. These results
should not be read as theorem evidence in their current form.
\end{remark}

\part{The refined H23a target}
\label{sec:target}

The finite-operator side is now largely a theorem. \cref{thm:class} classifies
the coboundaries; \cref{cor:strict} gives a strict gap at every finite level;
\cref{thm:collapse} reduces uniformity in $Q$ to a single sequence of computable
constants. What remains open is one statement about that sequence.

\begin{keystatement}
This work identifies the correct operator (the faithful \texttt{sync}
synchronized operator), \emph{proves} the correct quotient
($\operatorname{Cob}=\langle\chi_2\rangle$, \cref{thm:class}), proves a strict
non-coboundary gap at every finite level (\cref{cor:strict}), and proves that the
uniform gap is governed by the $r$-independent conductor sequence
$\gamma_c(s)$ (\cref{thm:collapse}):
\[
\kappa_{3^r}(s)=1-\max_{2\le c\le r}\gamma_c(s),\qquad
\gamma_c(s)<1\ \text{for every }c .
\]
The single remaining finite-operator statement is that
$\sup_c\gamma_c(s)<1$.
\end{keystatement}

\noindent\textbf{Proof status.}
\begin{itemize}[leftmargin=1.6em,itemsep=1pt]
\item \emph{Proved:} $\chi_2$ is an exact coboundary (\cref{thm:cobound});
$\operatorname{Cob}=\langle\chi_2\rangle$ on the residue--parity skew product
(\cref{thm:class}); strict gap $\rho(\Lcal_\chi)<\rho(\Lcal_{\mathbf1})$ at every
finite level for every non-coboundary $\chi$ and tilt (\cref{cor:strict});
conductor autonomy and collapse (\cref{prop:cond,thm:collapse}); the
height-strip rotation identity (\cref{prop:rot}); the residue-only $\chi_2$
ratio $\tfrac{2^{1+s}-1}{2^{1+s}+1}$ (\cref{prop:dich}); one-step mode damping
with the coboundary modes as the only undamped ones (\cref{lem:damp}); the
stability extraction step---near-peripheral spectrum yields a near-coboundary
(\cref{prop:stab}); \emph{exact descent}---the primitive sector at level $3^c$
is isospectral to a multiplier-twisted operator one level down, so deep-digit
comparisons are lossless (\cref{thm:descent,cor:descentcomp}).
\item \emph{Proved (canonical witness):} the explicit cycle pair of
\cref{thm:witness} has $W\equiv4\bmod9$, giving full principal-unit generation
and the uniform $\sqrt3$ cube-root separation for every primitive character---%
on the two-descent range unconditionally, and for all $c$ granted carry depth.
\item \emph{Proved (pair-correlation engine):} the phase formula
(\cref{prop:paircorr}); the \emph{exact nonstationary vanishing} lemma
(\cref{lem:vanish}); the collision classification and mass identity
(\cref{lem:collision}); the endpoint obstruction---fixed-$n$ operator-norm
bounds are structurally false, so the assembly must run through spectral/trace
quantities (\cref{prop:endpoint}); and the trace identity
$\operatorname{tr}((T^m)^*T^m)=\dim\Delta_s^m+\Theta_{m,c}$ with the thin
remainder verified $\le6\times10^{-4}$ and decreasing in $c$
(\cref{lem:trace}).
\item \emph{Proved (assembly bookkeeping):} the thin-sector bound
(\cref{lem:thin}), and---two further structural negative results---fixed-$m$
chain powers bound the norm rather than the spectral radius (non-normal
inflation), while the single trace saturates \emph{exactly} at $\rho^2\le1$ at
$s=0$, with zero margin (\cref{prop:marginal}).
\item \emph{Proved (wrap-regime collision flattening):} the walk--window
factorization $C_w=2^{B_{m-1}}X_w$ (\cref{lem:walkwindow}) and the
pigeonhole-flatness of the wrap-regime law,
$\operatorname{Col}_m(c)\le\mathrm{poly}(c)\,3^{-c}$ at $m\ge(1+\varepsilon)c$
(\cref{lem:flat}; measured $3^c\operatorname{Col}\approx2$): the deep digits of
the window variable determine the suffix exactly (no wrapping of partial sums
beyond scale $O(\log c)$). This is the renewal-flattening half of the wrap
conjecture, in collision form.
\item \emph{Open (one lemma):} lift-phase chain cancellation
(\cref{conj:wrap}): pure counting is once more exactly marginal, so the chain
decay must come from the character phases carried by the source lifts
$a_{j+1}\in a_j+3^{c-m}\Z$ at each link. It implies
$\limsup_c\gamma_c(s)\le\sqrt{\Delta_s}<1$ in the limit (\cref{thm:assembly}),
the high-conductor damping theorem with conjecturally sharp constant
$\gamma_\infty(s)=\sqrt{(2^{1+s}-1)/(2^{1+s}+1)}$---matched by the observed
tail at $s=0$.
\item \emph{Proved (three-level descent):} two scalar descents plus one rank-2
descent (\cref{thm:descent,thm:rank2}, the latter verified to twelve digits)---%
a fixed dimension reduction by $27/2$. \emph{Disproved (rank saturation):}
deeper carries depend on $b\bmod2\cdot3^j$ and the internal classes refine by
$\times3$ per level; the nonzero-eigenvalue count grows $\approx\times3$ per
level (\cref{rem:saturation}), so descent stalls at depth three and there is
\emph{no} bounded-level reduction. The fixed-$c_0$ holonomy program is
therefore limited to $c\le c_0+3$; the canonical witness and generation
theorems hold on that range.
\item \emph{Open (short-witness coverage = the quantitative heart):} generation
does not imply coverage; the uniform gap needs $\delta_c(L_0)\ge\delta_0$ at
fixed witness length (\cref{conj:mult}, strong $\Rightarrow$ constant gap;
weak/polylog $\Rightarrow$ $\kappa_Q\gg(\log Q)^{-2A}$, \cref{thm:condunif}).
The controlled-length audit finds near-maximal coverage already at $L=1$ on the
audited range; the fixed-$c_0$ asymptotics await Task~A. Comparison-graph
expansion (\cref{conj:exp}) remains an independent route to the polylog form
(\cref{thm:condpolylog}).
\item \emph{Open (strip classification):} the non-product-form coboundary
classification for the killed strip (the residue--parity case is settled;
product-form strip potentials are excluded numerically for $r\le4$).
\item \emph{Conditional:} a constant sync-quotient gap (\cref{cor:const});
a polylog gap under comparison-graph expansion (\cref{thm:condpolylog}).
\end{itemize}

\begin{conjecture}[Bounded-conductor extremality]\label{conj:ext}
For each tilt $s\in[0,1]$ the sequence $\gamma_c(s)$ of \cref{thm:collapse}
attains its supremum at a fixed finite conductor:
$\sup_{c\ge2}\gamma_c(s)=\max_{c\le c_0}\gamma_c(s)<1$. Numerically
$c_0=3$ ($s{=}0$) and $c_0=4$ ($s{=}1$) for the residue--parity operator
(through $c=6$), and $c_0=2$ for the killed strip at $s=0$ (through $c=5$, where
the sequence is monotone decreasing). Caution: monotonicity in $c$ \emph{fails}
for the residue--parity operator ($\gamma_4>\gamma_2$ at $s{=}1$), so the correct
target is boundedness/extremality, not monotonicity.
\end{conjecture}

\begin{corollary}[Conditional constant gap]\label{cor:const}
Assume \cref{conj:ext}. Then for all $Q=3^r$ ($r\ge2$) and $s\in[0,1]$,
\[
\rho(\Lcal_{s,Q,\chi})\le(1-\kappa_0(s))\,\rho(\Lcal_{s,Q,\mathbf1}),
\qquad \chi\notin\langle\chi_2\rangle,
\qquad
\kappa_0(s)=1-\max_{c\le c_0}\gamma_c(s)>0,
\]
a \emph{constant} determined by finitely many computable operators. For the
residue--parity operator the computed values give
$\kappa_0(0)=1-0.7638=0.2362$ and $\kappa_0(1)=1-0.9264=0.0736$; for the killed
strip at $s=0$, $\kappa_0=1-0.9069=0.0931$. This is stronger than the fall-back
target $\kappa_Q\gg1/\mathrm{polylog}(Q)$, for which it suffices that
$\gamma_c\le1-c^{-A}$ for some $A$.
\end{corollary}
\begin{proof}
Immediate from \cref{thm:collapse}: $\kappa_{3^r}=1-\max_{2\le c\le r}\gamma_c$,
and under \cref{conj:ext} the maximum is attained at $c\le c_0$, independent of
$r$. Each $\gamma_c<1$ by \cref{cor:strict} applied at level $3^c$.
\end{proof}

\section{Toward $\sup_c\gamma_c<1$: stability of the classification}
\label{sec:stability}

The route we consider most promising is a \emph{quantitative} version of
\cref{thm:class}: near-peripheral spectrum forces a near-coboundary, and the
exact rigidity of the classification proof should then produce a contradiction
for non-coboundary characters. Two ingredients of this program are proved here;
the third is formulated precisely.

\subsection*{One-step mode damping}

Write the Koopman form of the synchronized operator at level $Q=3^c$ as
$M_\chi g(a,\tau)=\chi(a)\sum_{b\ge1}w_b\,g(S_b(a),\tau+b)$ with
$w_b=3^sq^{\,b}$, $q=2^{-(1+s)}$, $W=\sum_bw_b=3^sq/(1-q)$. Fourier modes are
$\eta\otimes\epsilon$ with $\eta\in\widehat{U_Q}$ and $\epsilon\in\{0,1\}$ the
parity frequency. (Scope note: the classification \cref{thm:class} is
\emph{weight-free}; the explicit multiplier below is special to the geometric
branch weights $w_b\propto q^b$.)

\begin{lemma}[One-step mode damping]\label{lem:damp}
For every mode,
\[
M_\chi(\eta\otimes\epsilon)=c_{\eta,\epsilon}\,\bigl[\chi\cdot(\eta\circ A)\bigr]\otimes\epsilon,
\qquad
c_{\eta,\epsilon}=\sum_{b\ge1}w_b\,\zeta^b=\frac{3^sq\,\zeta}{1-q\zeta},
\qquad
\zeta=\eta(2)^{-1}(-1)^{\epsilon},
\]
so that
\[
\frac{\lvert c_{\eta,\epsilon}\rvert}{W}
=\frac{1-q}{\lvert1-q\zeta\rvert}\;\le\;1,
\]
with equality iff $\zeta=1$, i.e.\ iff
$(\eta,\epsilon)\in\{(\mathbf1,0),(\chi_2,1)\}$---precisely the two coboundary
modes (the constants and the potential $d(a,\tau)=\chi_2(a)(-1)^\tau$ of
\cref{thm:cobound}). Quantitatively, if $\theta=\arg\zeta$ then
$1-\lvert c\rvert/W\asymp q\theta^2/(1-q)^2$ for small $\theta$; the
``swapped'' modes $(\mathbf1,1)$ and $(\chi_2,0)$ have $\zeta=-1$ and damping
ratio exactly $\frac{1-q}{1+q}=\frac{2^{1+s}-1}{2^{1+s}+1}$, the constant of
\cref{prop:dich}.
\end{lemma}
\begin{proof}
$\eta(S_b(a))(-1)^{\epsilon(\tau+b)}
=\eta(3a+1)\,\eta(2)^{-b}(-1)^{\epsilon b}\,(-1)^{\epsilon\tau}$, and the $b$-sum
is geometric. $\eta(2)=1$ iff $\eta=\mathbf1$ and $\eta(2)=-1$ iff $\eta=\chi_2$,
since $2$ generates $U_Q$.
\end{proof}

\noindent\cref{lem:damp} quantifies the mechanism: every non-coboundary mode is
strictly damped in one step, at a rate governed by the angle of $\eta(2)$. It
also exposes the difficulty: the least-damped primitive mode at conductor $3^c$
has $\theta=\pi/3^{c-1}$, hence one-step damping $\asymp9^{-c}$; a proof of a
\emph{uniform} gap cannot ride a single mode and must use the spreading of
Fourier mass under $\eta\mapsto\eta\circ A$ (whose coefficients are Jacobi-type
sums with square-root cancellation). This is the analytic heart that remains.

\subsection*{Near-peripheral spectrum forces a near-coboundary}

\begin{proposition}[Stability, extraction step]\label{prop:stab}
Let $\mu$ be an eigenvalue of $M_\chi$ with $\lvert\mu\rvert=(1-\varepsilon)W$,
$\varepsilon\ge0$, and let $\varphi$ be an associated eigenfunction,
$\psi=\lvert\varphi\rvert$, $u=\varphi/\psi$ on $\{\psi>0\}$,
$\lambda=\mu/\lvert\mu\rvert$. Let $\nu$ be the positive left Perron vector of
$M_{\mathbf1}$. Define the nonnegative per-edge defect
\[
\delta(x,b)=w_b\Bigl[\psi(S_bx)-\operatorname{Re}\bigl(\overline{\lambda\chi(a)u(x)}\,
\varphi(S_bx)\bigr)\Bigr],\qquad x=(a,\tau).
\]
Then $\sum_b\delta(x,b)=\bigl(M_{\mathbf1}\psi-(1-\varepsilon)W\psi\bigr)(x)\ge0$
pointwise on $\{\psi>0\}$, and
\[
\sum_x\nu(x)\sum_b\delta(x,b)=\varepsilon W\,\langle\nu,\psi\rangle .
\]
Consequently, for every $t>0$, outside an edge set of
$\nu$-weighted defect mass $\le\varepsilon W\langle\nu,\psi\rangle/t$ the
unimodular potential $d=u$ satisfies the coboundary identity
$d(S_bx)\approx\lambda\chi(a)d(x)$ with
$w_b\,\psi(S_bx)\bigl(1-\cos(\text{phase error})\bigr)\le t$. In particular,
comparing the branches $b$ and $b{+}2$ of a common source yields approximate
$\langle4\rangle$-invariance, $d(4^{-1}v,\sigma)\approx d(v,\sigma)$, on
high-mass edges---the quantitative form of Step~1 of \cref{thm:class}.
\end{proposition}
\begin{proof}
The eigen-equation gives
$\lvert\mu\rvert\psi(x)=\operatorname{Re}\bigl(\overline{\lambda\chi(a)u(x)}\,
\chi(a)\sum_bw_b\varphi(S_bx)\bigr)
=\sum_bw_b\operatorname{Re}\bigl(\overline{\lambda\chi(a)u(x)}\varphi(S_bx)\bigr)
\cdot\lvert\chi(a)\rvert$, whence
$\sum_b\delta(x,b)=M_{\mathbf1}\psi(x)-\lvert\mu\rvert\psi(x)$, which is
$\ge0$ pointwise and has the stated $\nu$-mass since
$\langle\nu,M_{\mathbf1}\psi\rangle=W\langle\nu,\psi\rangle$. The tail bound is
Chebyshev's inequality; each defect term dominates
$w_b\psi(S_bx)(1-\cos)$ of the phase mismatch.
\end{proof}

\subsection*{Exact descent: lossless propagation across deep digits}

The deepest digits turn out to need no approximate propagation at all. The
mechanism is that the level-$3^c$ dynamics reads only $a\bmod3^{c-1}$ (the same
compression behind \cref{prop:cond}), so states in a common fiber of
$U_{3^c}\to U_{3^{c-1}}$ have \emph{identical} branch sums.

\begin{theorem}[Exact descent]\label{thm:descent}
Let $\chi$ be primitive of conductor $3^c$, $c\ge2$, and let $M_\chi$ be the
synchronized operator at level $3^c$ with branch weights $w_b>0$. Define, on
functions of $(\bar a,\tau)\in U_{3^{c-1}}\times\Z/2$,
\[
N_1g(\bar a,\tau)=\chi\bigl(A(\bar a)\bigr)\sum_{b\ge1}w_b\,\chi(2)^{-b}\,
g\bigl(S_b(\bar a)\bmod3^{c-1},\,\tau+b\bigr),
\]
where $A(\bar a)=3\bar a+1$ is a well-defined element of $U_{3^c}$ (it depends
only on $\bar a\bmod3^{c-1}$). Then
\[
\operatorname{spec}(M_\chi)\setminus\{0\}
=\operatorname{spec}(N_1)\setminus\{0\},
\]
with eigenspaces in bijection via $\varphi(a,\tau)=\chi(a)\,g(a\bmod3^{c-1},\tau)$.
In particular the primitive-sector ratio $\gamma$ at level $3^c$ equals
$\rho(N_1)/W$, computed one level down.
\end{theorem}
\begin{proof}
Suppose $M_\chi\varphi=\mu\varphi$, $\mu\ne0$. The branch sum
$\Sigma(\bar a,\tau)=\sum_bw_b\varphi(S_b(a),\tau+b)$ depends only on
$(\bar a,\tau)$, because $S_b(a)\bmod3^c=A(\bar a)2^{-b}$ is a function of
$(\bar a,b)$. Hence
$\varphi(a,\tau)=\mu^{-1}\chi(a)\Sigma(\bar a,\tau)=:\chi(a)g(\bar a,\tau)$
\emph{exactly}, at every state. Substituting this form of $\varphi$ at the
successors and using
$\chi(S_b(a))=\chi(A(\bar a))\chi(2)^{-b}$,
\[
\mu\,g(\bar a,\tau)=\sum_bw_b\,\chi(S_b(a))\,g\bigl(S_b(a)\bmod3^{c-1},\tau+b\bigr)
=N_1g(\bar a,\tau),
\]
and $g\ne0$ since $\varphi\ne0$. Conversely, given $N_1g=\mu g$ with $\mu\ne0$,
set $\varphi(a,\tau)=\chi(a)g(\bar a,\tau)\ne0$; the same computation gives
$M_\chi\varphi=\mu\varphi$. The two maps are mutually inverse on eigenspaces.
\end{proof}

\begin{remark}
Verified numerically: at $(Q,k,s)=(81,1,0),(81,1,1),(243,5,0)$ the nonzero
spectra agree (top eigenvalues match to the eigenvalue condition number,
$\rho$ to ten digits); the complementary spectrum of $M_\chi$ is nilpotent, seen
numerically as pseudospectral rings of radius $\approx10^{-8}$.
\end{remark}

\begin{corollary}[Descent and the comparison graph]\label{cor:descentcomp}
Every nonzero-eigenvalue eigenfunction satisfies
$\varphi(a,\tau)\chi(a)^{-1}=\varphi(a',\tau)\chi(a')^{-1}$ whenever
$a\equiv a'\pmod{3^{c-1}}$---the deep-digit comparisons of the conditional
argument hold \emph{exactly}, with zero defect, and the same factorization can
be iterated: composing descents expresses the primitive sector at level $3^c$
as an explicitly phased operator at any level $3^{c_0}$ (the multiplier of the
$j$-th descended operator is a product of $\chi$-values along depth-$j$
branches). Consequently the expansion needed in \cref{conj:exp} below is
required only at a \emph{bounded} level: what remains is that the accumulated
unimodular multiplier system at a fixed level $3^{c_0}$ stays uniformly (in
$c$) away from the coboundary class---a finite-dimensional non-degeneracy
statement, since at fixed dimension ``$\delta$-far from every coboundary''
implies a gap $\kappa(\delta)>0$ by Wielandt's theorem and compactness.
\end{corollary}

\subsection*{The descended multiplier and its holonomy}

By \cref{thm:descent} the bounded-level object carrying all remaining
information is the \emph{accumulated multiplier}: after descending from level
$3^c$ to level $3^{c_0}$, each edge $e$ of the fixed level-$c_0$ graph carries a
phase $\chi(F(e))$ for an explicit, character-independent element
$F(e)\in U_{3^c}$.

\begin{proposition}[Two-descent multiplier recursion]\label{prop:multrec}
With canonical least-residue lifts, the first descent gives
$F_1(x,b)=(3x+1)2^{-b}\in U_{3^c}$ (the level-$c$ lift of the successor), and
the second descent gives
\[
F_2(\hat x,b)=F_1(\hat x,b)\,\rho^{\,\delta_b(\hat x)},\qquad
\rho=1+3^{c-1},
\]
where $\delta_b(\hat x)\in\{0,1,2\}$ is the top digit of the successor
$S_b(\hat x)$ at the upper level relative to the canonical lift of its
reduction. Both steps are exact: the fiber ratio
$F_j(x+\delta3^{\ell},b)F_j(x,b)^{-1}=\rho^{\delta}$ is independent of $(x,b)$
for $j\le2$ (verified exactly in computation). At depth $j\ge3$ the fiber ratio
acquires a parity twist---the successor's top digit shifts by a quantity
depending on $b\bmod2$ through $2^{-b}\equiv(-1)^b\pmod3$---so the multiplier
must be treated as an edge multiplier on the residue--parity graph; the
resolution is the rank-$2$ descent below.
\end{proposition}

\begin{theorem}[Rank-2 parity-twisted descent]\label{thm:rank2}
At the third descent the fiber ratio of the accumulated multiplier has the exact
form
\[
\frac{m(\hat a+\delta3^{L-1},b)}{m(\hat a,b)}
=\chi\bigl(r_0(\hat a,\delta)\bigr)\cdot\omega^{\,\delta(-1)^b},
\qquad
r_0(\hat a,\delta)=1+\delta3^{L}(3\hat a+1)^{-1}\in1+3^{L}\Z,
\quad
\omega=\chi(1+3^{c-1}),
\]
with $L$ the current level. Splitting the branch sum by parity classes
$b\equiv\epsilon\pmod2$ and setting
$A_\epsilon(\hat a,\tau)=\sum_{b\equiv\epsilon}w_b\,m(\hat a,b)\,
g(S_b\hat a,\tau+\epsilon)$, every $\mu\ne0$ eigenfunction satisfies the exact
fiber expansion
$\mu\,g(\hat a+\delta3^{L-1},\tau)
=\chi(r_0)\bigl[\omega^{\delta}A_0(\hat a,\tau)+\omega^{-\delta}A_1(\hat a,\tau)\bigr]$,
and the pair $(A_0,A_1)$ satisfies a closed eigenproblem one level down:
\[
\mu A_\epsilon(\hat a,\tau)=\sum_{b\equiv\epsilon(2)}w_b\,m(\hat a,b)\,
\chi\bigl(r_0(\hat v_b,\delta_b)\bigr)
\bigl[\omega^{\delta_b}A_0(\hat v_b,\tau+\epsilon)
+\omega^{-\delta_b}A_1(\hat v_b,\tau+\epsilon)\bigr],
\]
with $\hat v_b$ the reduced successor and $\delta_b$ its top digit. The
correspondence is a bijection on nonzero eigenspaces: the primitive sector
descends exactly onto a \emph{rank-2} operator on
$U_{3^{L-1}}\times\Z/2\times\{0,1\}$.
\end{theorem}
\begin{proof}
The fiber-ratio formula is the computation in the proof sketch of
\cref{prop:multrec} made exact: the $F_1$-part of the ratio is
$(3(\hat a+\delta3^{L-1})+1)/(3\hat a+1)=r_0$, independent of $b$, and the
accumulated $\rho^{\delta_b}$-part shifts by
$\rho^{\delta(-1)^b}$ because the successor moves by
$\delta3^{L-1}\cdot3\cdot2^{-b}\equiv\delta(-1)^b3^{L}$ at the upper level. The
fiber expansion follows by splitting the eigen-equation at the fiber point into
parity classes; substituting it at the successors closes the system on
$(A_0,A_1)$, and the two maps (restriction to $(A_0,A_1)$; reconstruction of
$g$ from the fiber expansion) are mutually inverse on nonzero eigenspaces,
exactly as in \cref{thm:descent}. Verified numerically at $(c,k,s)=(5,1,0)$:
the full $324$-dimensional primitive operator and the $24$-dimensional rank-2
descended operator have equal spectral radius to twelve digits and matching
significant spectra.
\end{proof}

\begin{proposition}[Negative result: rank saturation fails; descent stalls at
depth three]\label{rem:saturation}
An earlier draft conjectured that the internal rank stays $2$ under all further
descents. \emph{This is false.} At the second rank-2 step the fiber variation
of the accumulated $\rho^{\delta_b}$-factors shifts the successor at \emph{two}
digit positions, through $\delta\cdot2^{-b}\bmod9$: the carry depends on
$b\bmod6$, not just on the parity $b\bmod2$, and each further descent refines
the internal classes by another factor of $3$. The bounded-level reduction
therefore does not continue. This is confirmed unconditionally by the
nonzero-eigenvalue count of the full primitive operator ($s=0$):
\[
\begin{array}{c|cccc}
c & 4 & 5 & 6 & 7\\\hline
\#\{|\lambda|>10^{-2}\} & 8 & 20 & 56 & 160\\
4\phi_{c-3}\ \text{(three-level reduction)} & 24 & 24 & 72 & 216
\end{array}
\]
The count grows by a factor $\approx3$ per level and is bounded by
$4\phi_{c-3}$, exactly as predicted by a \emph{fixed three-level} reduction
(two scalar descents $+$ one rank-2 descent, a constant dimension factor
$27/2$) followed by no further collapse. Consequently the fixed-$c_0$
multiplier tables of the preceding discussion exist only for $c\le c_0+3$, and
the bounded-level holonomy program (\cref{conj:mult} in its fixed-$c_0$,
$c\to\infty$ reading) cannot be fed by descent: the witness audit below is
valid on its stated range and does not extend asymptotically by this route.
What survives in full: exact descent and the rank-2 step
(\cref{thm:descent,thm:rank2}); the carry-depth property \emph{on that range};
the canonical witness and full generation on that range; the classification,
collapse, and strict-gap theorems, which never used bounded-level descent. The
quantitative uniform-gap question reverts to the conductor sequence itself
(\cref{conj:ext}), for which the descent data adds new favourable evidence: the
primitive-sector radii \emph{decrease} in $c$
($0.7629,\,0.6454,\,0.6154,\,0.6025$ at $k$ small, $s=0$, $c=4,\dots,7$)---high
conductor is increasingly damped, consistent with bounded-conductor
extremality.
\end{proposition}

\subsection*{The high-conductor damping problem}

With bounded-level descent excluded, the remaining quantitative theorem is a
\emph{tail estimate}, attacked as an oscillation problem rather than a
compactness problem.

\begin{conjecture}[High-conductor damping]\label{conj:hcd}
There are $c_1$ and $\delta>0$ such that for every $c\ge c_1$, every primitive
$\chi$ of conductor $3^c$, and every $s\in[0,1]$,
$\rho(M_\chi)\le(1-\delta)W$. With the strict gap at each finite level
(\cref{cor:strict}) and conductor collapse (\cref{thm:collapse}), this is
equivalent to $\sup_c\gamma_c(s)<1$ (\cref{conj:ext}).
\end{conjecture}

\noindent The available mechanisms are: the exact coboundary classification
(\cref{thm:class}); one-step mode damping with equality only at the coboundary
modes (\cref{lem:damp}); stability extraction (\cref{prop:stab}); three-level
descent and the eigencount control (\cref{thm:descent,thm:rank2,rem:saturation});
and the numerically decreasing tail of $\gamma_c$.

In the character basis the operator has matrix
$K[t_1,t_0]=c_{t_0,\epsilon}\,J(k-t_1,t_0)$, where
\[
J(u,v)=\frac1\phi\sum_{a\in U_{3^c}}e_u(a)\,e_v(3a+1)
\]
are $p$-adic Jacobi sums. Two structural facts, computed exactly for
$c\le7$: \emph{(i)} conductor lowering makes each column of $J$ supported on a
fixed fraction $2/9$ of the frequencies, with exact resonance columns
($\lvert J\rvert=1$) at low conductor; \emph{(ii)} the entrywise bound is
useless---$\rho(\lvert K\rvert)/W=1.85,\,3.22,\,5.58,\,9.66$ at $c=4,\dots,7$,
growing by exactly $\sqrt3$ per level, while the true ratio is
$\approx0.6$: at $c=7$ a factor $\approx16$ of \emph{phase cancellation}
separates the truth from any magnitude argument. The gap is therefore a genuine
oscillation phenomenon. The natural formulations are:

\begin{enumerate}[leftmargin=1.6em]
\item \emph{Mode spreading / finite-time damping:} show that along the mode
evolution $\eta\mapsto\chi\cdot(\eta\circ A)$ the near-undamped condition
$\eta(2)^{-1}(-1)^\epsilon\approx1$ cannot persist for $m$ consecutive steps
(with $m$ absolute, or $m\ll c^A$ for polylog), using the phases of $J$, not
only the magnitudes.
\item \emph{Van der Corput / pair correlation:} bound
$\lVert M_\chi^n\rVert_2$ for $n\asymp c$ via the word-pair correlation sums
$\sum_a\chi\bigl(\Phi_w(a)/\Phi_{w'}(a)\bigr)$; stationarity of the phase
ratios forces the exact coboundary relations, which are classified, so only
the $\langle\chi_2\rangle$ sector can resist---a quantitative use of
\cref{thm:class}. The arithmetic input is $p$-adic stationary phase for the
rational functions $\Phi_w/\Phi_{w'}$, where bounded (not square-root)
cancellation suffices.
\end{enumerate}

\subsection*{The pair-correlation engine}

The second route is now equipped with its exact arithmetic. Both pieces below
are rigorous; the first is a restatement of \cref{lem:cong}'s affine cocycle.

\begin{proposition}[Pair-correlation phase formula]\label{prop:paircorr}
For a branch word $w=(b_0,\dots,b_{n-1})$ set $B_j(w)=\sum_{i<j}b_i$ and
\[
C_w=\sum_{j=0}^{n-1}3^{\,n-1-j}2^{B_j(w)} .
\]
Then the $n$-step map along $w$ is
$S_w(a)=(3^na+C_w)\,2^{-B_n(w)}$, and for two words $w,w'$,
\[
R_{w,w'}(a)=\frac{S_w(a)}{S_{w'}(a)}
=2^{\,B_n(w')-B_n(w)}\,\frac{3^na+C_w}{3^na+C_{w'}},
\qquad
\frac{R'_{w,w'}}{R_{w,w'}}(a)
=\frac{3^n\,(C_{w'}-C_w)}{(3^na+C_w)(3^na+C_{w'})}.
\]
Since $C_w\equiv2^{B_{n-1}}\not\equiv0\pmod3$, the denominators are units on
$U_{3^c}$, so the logarithmic derivative has \emph{constant} $3$-adic valuation
$d=n+v_3(C_w-C_{w'})$ with unit leading part. Stationarity of the pair phase is
exactly high $3$-adic divisibility of $C_w-C_{w'}$; note $C_w$ does not depend
on the last symbol $b_{n-1}$.
\end{proposition}

\begin{lemma}[Exact nonstationary vanishing]\label{lem:vanish}
Let $\chi$ be primitive of conductor $3^c$ and let $w,w'$ satisfy
$d=n+v_3(C_w-C_{w'})\le c-2$. Then
\[
\sum_{a\in U_{3^c}}\chi\bigl(R_{w,w'}(a)\bigr)=0 .
\]
\end{lemma}
\begin{proof}
Write $L=(\log R)'=R'/R$, so $v_3(L(a))=d$ with unit part
$\mu(a)=3^{-d}L(a)$, and $v_3(L')\ge d+n>d$. Split $a=a_0+3^{c-d-1}t$ with
$t\in\Z/3^{d+1}$ (every lift of a unit is a unit). Taylor expansion of
$\log R$ gives
\[
R(a)/R(a_0)=1+3^{c-1}\,t\,\mu(a_0)\,(1+O(3))+E,
\]
where every higher term $E$ has valuation
$\ge2(c-d-1)+d=2c-d-2\ge c$ because $d\le c-2$. Hence
$\chi(R(a))=\chi(R(a_0))\,e\bigl(\kappa\,\mu(a_0)t/3\bigr)$ for a unit
$\kappa$ ($\chi$ primitive is nontrivial on $1+3^{c-1}\Z$). For each fixed
$a_0$ the inner sum over $t\in\Z/3^{d+1}$ runs over full periods of a
nontrivial third root of unity ($\mu(a_0)$ is a unit), so it vanishes
identically.
\end{proof}

\noindent Verified exactly: for $n=2,3$ and $c=6,7$ the sampled pair sums are
machine zeros precisely on the predicted range $n+v_3(C_w-C_{w'})\le c-2$, with
structured values ($\tfrac12$, $1$) only in the near-stationary transition
zone.

\begin{lemma}[Collision classification and collision mass]\label{lem:collision}
$C_w$ depends only on the prefix $(b_0,\dots,b_{n-2})$, via the recursion
$C_w=3C_{w^-}+2^{B_{n-1}}$. Conversely, distinct prefixes give distinct
constants: if $B_{n-1}(w)=B_{n-1}(w')$ this follows by induction
($3(C_{w^-}-C_{w'^-})=0$ forces prefix equality); the unequal-$B_{n-1}$ case is
excluded by $2$-adic digit analysis of
$2^{B}(2^{\Delta}-1)=3(C_{w^-}-C_{w'^-})$ (the left side has $2$-adic valuation
$B_{n-1}\ge n-1$, forcing implausibly deep $2$-adic agreement of the prefixes),
verified exhaustively for $n\le5$, $b_i\le12$: the number of distinct constants
is exactly $12^{\,n-1}$. Consequently the exactly-stationary pair mass is
\[
\sum_{C}\Bigl(\sum_{w:\,C_w=C}p_s(w)\Bigr)^2=\Delta_s^{\,n-1},
\qquad
\Delta_s=\sum_b p_s(b)^2=\frac{2^{1+s}-1}{2^{1+s}+1},
\]
confirmed numerically to truncation accuracy ($0.05$--$0.2\%$ at $B=12$).
\end{lemma}

\subsection*{Assembly via trace powers}

Two structural facts dictate the shape of the assembly.

\begin{proposition}[Endpoint obstruction: fixed-$n$ norm bounds fail]
\label{prop:endpoint}
A matrix entry of $(T^n)^*T^n$ constrains the source $a$ to a coset of
$1+3^{c-n}\Z$, on which every nonstationary pair phase is \emph{constant}
(its logarithmic derivative has valuation $d\ge n$, and
$(c-n)+d\ge c$). Hence \cref{lem:vanish} produces no within-entry
cancellation, and the fixed-$n$ bound
$\lVert T^n\rVert^2\le\Delta_s^{n-1}+o_c(1)$ is \emph{false}: at $s=0$,
$\lVert T^2\rVert^2\approx1.63$ versus $\Delta_0=0.33$ (stable in $c$). The
operator norm of a fixed power is inflated by the non-normal (nilpotent)
structure; only spectral quantities can obey the $\Delta_s$ scaling.
\end{proposition}

\begin{lemma}[Trace identity, main term]\label{lem:trace}
For fixed $m$ and $B$, and $c$ large,
\[
\operatorname{tr}\bigl((T^m)^*T^m\bigr)
=\dim\cdot\Delta_s^{\,m}\;+\;\Theta_{m,c},
\]
where the main term comes from word pairs with coincident endpoints and equal
$2$-power parts: $S_w(a)=S_{w'}(a)$ with $B_m(w)=B_m(w')$ forces $C_w=C_{w'}$
with the source $a$ ranging over \emph{all} of $U_{3^c}$; by
\cref{lem:collision} such pairs agree in the first $m-1$ symbols, hence (the
path phase $\Psi_w(a)=\prod_{i<m}\chi(x_i)$ being blind to the last symbol)
have \emph{identical} path phases, contributing mass exactly $\Delta_s^m$ per
state; equal $B_m$ then forces $w=w'$. The thin remainder $\Theta_{m,c}$
collects pairs with $2^{B_m(w)}\ne2^{B_m(w')}$, whose source $a$ is confined to
lower-dimensional fibers. Numerically
$\Theta_{m,c}/(\dim\Delta_s^m)\le6\times10^{-4}$ for $m\le4$ at $c=6$,
\emph{decreasing} in $c$ ($\le2\times10^{-4}$ at $c=7$), and the $m=1$
identity is exact.
\end{lemma}

\begin{lemma}[Thin-sector bound]\label{lem:thin}
In the no-wrap regime (branch cutoff $B$, word length $m$ with
$C_w<3^c$ for all truncated words), the thin remainder of \cref{lem:trace}
satisfies
\[
\Theta_{m,c}\;\le\;C\,m^2B\,\bigl(\tfrac32\bigr)^{m}\cdot3^{m} ,
\qquad\text{hence}\qquad
\frac{\Theta_{m,c}}{\dim\Delta_0^{\,m}}\le
C\,m^2B\,\Bigl(\tfrac92\Bigr)^{m}3^{-c}\xrightarrow[c\to\infty]{}0
\quad\text{for }m\le0.7\,c .
\]
\end{lemma}
\begin{proof}
A thin pair has $2^{B(w')}\ne2^{B(w)}$; the endpoint equation
$3^ma\,(2^{B'}-2^{B})=C_{w'}2^{B}-C_w2^{B'}$ has unit solutions only if
$v_3(\mathrm{RHS})=m+v$, $v=v_3(2^{\Delta}-1)\le1+\log_3(mB)$, and then $a$ is
confined to a class with at most $3^{m+v}$ unit solutions. Solvability
constrains $C_{w'}\equiv C_w2^{\,B'-B}\pmod{3^{m+v}}$; in the no-wrap regime
the collision classification (\cref{lem:collision}) puts at most one prefix in
each class, of mass $\le(1-r)^{m-1}\le2^{-(m-1)}$ at $s=0$. Summing over the
$O(mB/3^v)$ values of $\Delta$ at each valuation $v$ and bounding all phases by
$1$ gives the display; the comparison with $\dim\Delta_0^m=4\cdot3^{c-1-m}$ is
immediate.
\end{proof}

\begin{proposition}[Two structural obstructions to the naive assembly]
\label{prop:marginal}
\emph{(a) Fixed-$m$ chain powers fail.} For fixed $m$,
$\operatorname{tr}\bigl[((T^m)^*T^m)^k\bigr]^{1/k}\to\lVert T^m\rVert^2$ as
$k\to\infty$, and the non-normal inflation makes $\lVert T^m\rVert^2>1$ for
small $m$ ($\lVert T^2\rVert^2\approx1.62$, $\lVert T^3\rVert^2\approx1.15$,
stable in $c$): chain powers at fixed $m$ bound the \emph{norm}, not the
spectral radius, and yield no gap (verified:
$\operatorname{tr}(P_2^k)=324,\,350,\,402$ at $c=7$, flat in $k$).
\emph{(b) The single trace is exactly marginal.} From
$\rho^{2m}\le\operatorname{tr}((T^m)^*T^m)=\dim\Delta_s^m(1+o(1))$ one gets
$\rho^2\le3^{\,c/m}\Delta_s(1+o(1))$, optimal at the largest admissible $m$;
the pigeonhole floor (collision mass $\ge\asymp3^{-c}$) blocks $m\gg c$, and at
$m\asymp c$ the bound saturates at $\rho^2\le3\Delta_s$---which at $s=0$ equals
$1$ \emph{exactly}. The structural coincidence $\Delta_0=1/3$,
$\dim\asymp3^c$ leaves zero margin: no single-trace argument can produce the
gap.
\end{proposition}

\noindent The two obstructions intersect to locate the true remaining problem:
the trace analysis must run in the \emph{wrap regime}, with $m\asymp c$
\emph{and} chain order $k\ge2$, where the word constants wrap modulo $3^c$ and
collisions are governed by the equidistribution of $C_w\bmod3^c$ rather than by
integer injectivity.

\subsection*{The wrap regime: collision flattening (proved) and lift-phase
cancellation (open)}

The wrap-regime analysis splits into a distributional half, which we can prove,
and a chain-phase half, which we localize precisely.

\begin{lemma}[Walk--window factorization]\label{lem:walkwindow}
For any word $w$ of length $m\ge c$,
\[
C_w\equiv2^{\,B_{m-1}}\,X_w\pmod{3^c},
\qquad
X_w=\sum_{k=0}^{c-1}3^{k}\,2^{-(b_{m-1-k}+\cdots+b_{m-2})}\ \ (k\ge1),\ \
X_w\equiv1+3\Z ,
\]
i.e.\ the constant is a unit multiple of a \emph{window variable} depending
only on the last $c-1$ symbols, by the same reversal as the sliding-window
congruence (\cref{lem:cong}). The walk factor $2^{B_{m-1}}$ and the window are
coupled only through the last $c-1$ increments.
\end{lemma}

\begin{lemma}[Wrap-regime collision flattening]\label{lem:flat}
At $s=0$, for $m\ge(1+\varepsilon)c$,
\[
\operatorname{Col}_m(c)=\Pr\bigl[C_w\equiv C_{w'}\bmod3^c\bigr]
\;\le\;K(c)\,3^{-c},
\qquad K(c)=O(\mathrm{poly}(c)) ,
\]
and numerically $3^c\operatorname{Col}_m(c)\approx2$ at $m=2c$
($1.75,1.81,1.88,1.93,1.98$ for $c=4,\dots,8$), $\approx4$ at $m=c$: the
wrap-regime law of $C_w$ is pigeonhole-flat up to a small constant.
\end{lemma}
\begin{proof}[Proof sketch]
By \cref{lem:walkwindow} and Cauchy--Schwarz (conditioning on the walks and one
window), $\operatorname{Col}_m(c)\le\max_\rho\Pr[X'/X=\rho]\le
\operatorname{Col}(X)$. For the window variable, the digit at scale $3^k$
constrains the backward partial sum $s_k$ modulo $2\cdot3^{k-1}$; since
$s_k\le\max b\cdot c=O(c)$, there is \emph{no wrapping} once
$2\cdot3^{k-1}>s_{\max}$, i.e.\ for all $k\ge k_0=O(\log_3c)$: the deep digits
determine the partial sums---hence the suffix symbols---\emph{exactly}. Fibers
of $X$ are therefore single suffixes up to the $O(\log c)$ shallow digits, and
\[
\operatorname{Col}(X)\le3^{O(\log c)}\sum_{\text{suffixes}}p(\cdot)^2
=\mathrm{poly}(c)\,\Delta_0^{\,c-O(\log c)}=\mathrm{poly}(c)\,3^{-c+O(\log c)} .
\qedhere
\]
\end{proof}

\noindent This proves the renewal-flattening half of the wrap conjecture in its
$L^2$ (collision) form: the distribution of $C_w\bmod3^c$ offers no
collision excess for the chain trace to exploit against us. What it does
\emph{not} yet provide is the gap itself: a pure counting bound on
$\operatorname{tr}(P_m^k)$---keeping only collided links and bounding all
phases by $1$---is once again \emph{marginal} (the per-link factor
$3^m\cdot\text{per-target mass}$ is $\asymp K\ge1$ at every $m$, by the same
zero-margin coincidence as \cref{prop:marginal}). The decay seen in the
measured chain traces must therefore come from the \emph{lift phases}: at each
chain link the source lifts $a_{j+1}\in a_j+3^{c-m}\Z$ carry character phases
through the next link's words, and these lift sums cancel
(vanishing-lemma style, in the lift variable) except on an exponentially thin
set. The remaining statement is:

\begin{conjecture}[Lift-phase chain cancellation; \emph{superseded}]\label{conj:wrap}
For $m\asymp c$ and some $k_0\ge2$, the chain trace with lift-phase
cancellation included satisfies
$\operatorname{tr}\bigl[((T^m)^*T^m)^{k_0}\bigr]\le
3^{\,c-\epsilon mk_0+o(c)}$ for some $\epsilon>0$, uniformly over primitive
$\chi$. Together with \cref{lem:flat} this is the full wrap-regime input.
\emph{Status: superseded as the live finite-side input. The threshold
cascade route (\cref{lem:pinning,prop:rowcontract}) requires neither
chain traces nor lift-phase cancellation; its single analytic input is the
threshold flattening statement \cref{op:flat}. This conjecture and
\cref{thm:assembly} below are retained as documentation of the earlier
trace-side route.}
\end{conjecture}

\begin{theorem}[Lift-sum factorization]\label{thm:liftfact}
In the chain-trace expansion of $\operatorname{tr}[((T^m)^*T^m)^k]$, fix the
words and the base source $a_1$. The link constraints give
$a_{j+1}=a_1+D_j+3^{c-m}T_j$ with $D_j$ word-determined and
$T_j=t_1+\cdots+t_j$ the cumulative lifts; the chain phase collapses to a
product over the \emph{adjacent} word pairs at the sources,
\[
\prod_{j=1}^{k}\chi\bigl(R_j(a_j)\bigr),
\qquad
R_j=\text{pair-correlation phase of }(w_{j-1},w'_j)
\text{ (the $i{=}0$ factors cancel),}
\]
and since $a_j$ depends only on $T_{j-1}$ and the map
$(t_1,\dots,t_{k-1})\mapsto(T_1,\dots,T_{k-1})$ is triangular, the lift sum
factorizes \emph{exactly} into independent one-variable sums:
\[
\sum_{\mathbf t\in(\Z/3^m)^{k-1}}\ \prod_{j}\chi\bigl(R_j(a_j)\bigr)
=\chi(R_1(a_1))\prod_{j=2}^{k}\ \sum_{T\in\Z/3^m}
\chi\bigl(R_j(\alpha_j+3^{c-m}T)\bigr).
\]
The multivariate lift-phase problem is therefore \emph{univariate}: a coset
character sum of a pair-correlation phase, per link.
\end{theorem}

\begin{lemma}[Coset dichotomy]\label{lem:coset}
Let $R$ be a pair-correlation phase, $\chi$ primitive mod $3^c$, and let
$d=v_3\bigl((\log R)'(\alpha)\bigr)$ be the \emph{pointwise} valuation at the
coset base. Then (under the window $2(c-m)+d+1\ge c$)
\[
\sum_{T\bmod 3^m}\chi\bigl(R(\alpha+3^{c-m}T)\bigr)=
\begin{cases}
0, & d\le m-1,\\
3^m\,\chi(R(\alpha)), & d\ge m .
\end{cases}
\]
There is no intermediate boundary behaviour: at $d=m-1$ the Postnikov linear
coefficient $\lambda_1=3^{-(c-1)}\bigl[\log R(\alpha+3^{c-m})-\log
R(\alpha)\bigr]$ is a unit by definition of $d$, and the full-period mod-$3$
sum vanishes. Verified exhaustively: at $(c,m,n)=(8,4,4)$ and $(7,3,3)$, all
$1{,}734$ sampled configurations with pointwise $d=m-1$ have
$\lambda_1\in\{1,2\}$ and machine-zero sums; all non-vanishing cases have
$\lambda_1=0$, i.e.\ true valuation $\ge m$. (An earlier draft reported a
boundary anomaly; it was an artifact of classifying by the \emph{nominal}
valuation $\min_i[i+v_3(\Delta C_i)]$---degenerate minima, where the leading
units cancel mod $3$, raise the true valuation in $\approx16\%$ of nominal
boundary cases. The degeneracy costs one mod-$3$ condition per level, so the
migrated mass is geometrically controlled.)
\end{lemma}

\begin{proposition}[Trace-method ceiling]\label{rem:sixth}
By \cref{thm:liftfact,lem:coset}, a chain survives only if every adjacent pair
is true-stationary ($v_3((\log R_j)')\ge m$), which prefix-locks the words up
to geometrically thin degenerate-minima corrections; the surviving phases are
then \emph{exactly constant} per coset, and two further exact resources
apply---the base-point sum over $a_1$ (a complete $U$-sum, vanishing unless
$v_3(\sum_jL_j)\ge c-1$) and the cyclic closure constraint (worth
$\mathrm{poly}\cdot3^{-(c-m)}$ by collision flattening). Carrying out the
complete accounting with all of these, every variant lands at bounds of the
form $\rho^2\le3^{\,1/m}\cdot(\text{const}/\sqrt m)^{1/m}\to1$: strictly below
$1$ at every finite level---re-deriving the strict gap---but with no uniform
$\delta$. Since the coset weights are exactly $\{0,1\}$ (no intermediate
values; \cref{lem:coset}), the boundary-transition-operator route degenerates,
and the conclusion is structural: \emph{the trace method has a marginality
ceiling}. The surviving deep-stationary chains carry unimodular constant
phases and trace mass $\Theta(3^k)$, which the inequality
$\rho^{2mk}\le\operatorname{tr}$ cannot convert into a uniform gap---the trace
mass is carried by many small eigenvalues, not by $\rho$. A proof of
$\sup_c\gamma_c<1$ therefore requires a non-trace mechanism: either direct
operator-norm control of $\lVert T^mg\rVert$ on the structured complement of
the deep-stationary sector (an eigenvector-aware estimate), or the finite-$Q$
certificate route of the prior program. The numerically observed limit
$\gamma_\infty(s)=\sqrt{\Delta_s}$ stands as the conjectural sharp constant.
\end{proposition}

\subsection*{The live route: eigenvector-aware norm control}

Three measurements settle how the non-trace route must run. Let
$P_{\mathrm{stat}}$ project onto the modes of conductor $\le3^{c-m}$ (the
functions constant on the lift cosets---the deep-stationary directions).

\emph{(i) Top eigenvectors avoid the stationary sector.} At
$(c,k)=(6,1),(7,5)$ the top right eigenvector carries mass only
$0.01$--$0.11$ on $\mathcal H_{\mathrm{stat}}$ for $m=2,3,4$: the trace
ceiling is produced by spectrally irrelevant small modes, exactly as the
ceiling analysis predicted.

\emph{(ii) The compressed stationary operator vanishes---exactly.} The
measurement $\rho\bigl(P_{\mathrm{stat}}T^mP_{\mathrm{stat}}\bigr)^{1/m}
\le3\times10^{-4}$ that motivated the block decomposition is in fact machine
roundoff of an exact identity:

\begin{theorem}[Conductor raising; exact stationary ejection]\label{thm:eject}
Let $\chi$ be primitive modulo $3^c$ and let $P_{\mathrm{top}}$ project onto
the modes $\eta\otimes\varepsilon$ with $\eta$ of conductor exactly $3^c$.
Then
\[
P_{\mathrm{top}}^{\perp}\,T_{\chi,c,s}=0:
\]
the image of the synchronized Koopman operator lies entirely in the
top-conductor sector. Consequently, for every threshold $j<c$, writing
$P_{\le j}$ for the projection onto modes of conductor $\le3^j$,
\[
P_{\le j}\,T^m=0\qquad(m\ge1),
\]
so in particular $P_{\mathrm{stat}}T^mP_{\mathrm{stat}}=0$ identically; and
$\operatorname{spec}(T)\setminus\{0\}
=\operatorname{spec}(P_{\mathrm{top}}TP_{\mathrm{top}})\setminus\{0\}$: all
spectral action takes place in the primitive sector, of dimension
$8\cdot3^{c-2}$ (two thirds of the space).
\end{theorem}

\begin{proof}
It suffices to check the spanning modes $g=\eta\otimes\varepsilon$. By the
synchronization rule,
\[
(Tg)(a,\tau)
=\chi(a)\,\eta(3a+1)\,\varepsilon(\tau)\cdot
\sum_{b\ge1}w_b\,\eta(2)^{-b}\varepsilon(b),
\]
so up to a scalar the output is the function $a\mapsto\chi(a)\,\eta(3a+1)$
(times a parity character). Since $3a+1\bmod3^c$ depends only on
$a\bmod3^{c-1}$, the function $a\mapsto\eta(3a+1)$ factors through the
reduction $U_{3^c}\to U_{3^{c-1}}$ and hence lies in the span of the
multiplicative characters $\zeta$ of conductor $\le3^{c-1}$ (these are exactly
the functions invariant under $\ker(U_{3^c}\to U_{3^{c-1}})$, by dimension
count). For any such $\zeta$, the restriction of $\chi\zeta$ to that kernel
equals the restriction of $\chi$, which is nontrivial because $\chi$ is
primitive; hence $\operatorname{cond}(\chi\zeta)=3^c$ exactly. Every Fourier
component of the output therefore has top conductor. The corollaries follow
from $P_{\le j}P_{\mathrm{top}}=0$ and
$\operatorname{im}T\subseteq\mathcal H_{\mathrm{top}}$.
\end{proof}

\begin{remark}
The identity is one-sided: $TP_{\mathrm{top}}^{\perp}\ne0$ (measured norm
$\approx1.16$)---low-conductor modes feed the primitive sector and never
return. In transfer (pushforward) form the statement transposes to
$MP_{\le j}=0$; both forms were verified to $10^{-13}$ at $c=5,6$, together
with the exact agreement of $\operatorname{spec}(T)$ and
$\operatorname{spec}(P_{\mathrm{top}}TP_{\mathrm{top}})$ to all computed
digits. This settles the block-decomposition corner exactly ($A=0$, not
merely small), and it explains measurement (i) structurally: the
deep-stationary trace mass counted by the ceiling (\cref{rem:sixth}) lives in
directions that $T$ ejects in a single step and that no power ever revisits.
\end{remark}

\emph{(iii) The norm decay rate is empirically $c$-uniform at $m\gtrsim c$.}
\[
\begin{array}{c|ccccc}
\sigma_1(T^m)^{1/m} & m{=}4 & m{=}8 & m{=}12 & m{=}16 & m{=}20\\\hline
c=5 & 0.9647 & 0.8088 & 0.7470 & 0.7199 & 0.7044\\
c=6 & 0.9626 & 0.8373 & 0.7655 & 0.7284 & 0.7087\\
c=7 & 0.9644 & 0.8658 & 0.7782 & 0.7302 & 0.7025
\end{array}
\]
At $m=20$ the rate is flat in $c$ to $\pm0.6\%$. The structural reason is
exact: for $m\ge c$ the endpoint obstruction of \cref{prop:endpoint}
\emph{evaporates}.

\begin{proposition}[Complete-kernel formula]\label{prop:completekernel}
Let $m\ge c$. Then $3^m\equiv0\pmod{3^c}$, so the $m$-step endpoint
$S_w(a)\equiv C_w2^{-B_m(w)}\pmod{3^c}$ is determined by the word alone,
independently of the source $a$. Consequently the kernel of $T^m$ factorizes
as
\[
K_m\bigl((x,\sigma),(a,\tau)\bigr)
=\sum_{\substack{w:\;C_w2^{-B_m(w)}\equiv x\ (3^c)\\ B_m(w)\equiv\sigma-\tau\ (2)}}
p_s(w)\,\Psi_w(a),
\qquad
\Psi_w(a)=\prod_{i=0}^{m-1}\chi\bigl(x_i(a,w)\bigr),
\]
and the entries of $H_m=(T^m)^*T^m$ are \emph{complete word-pair sums}:
\[
H_m\bigl((a,\tau),(a',\tau')\bigr)
=\sum_{\substack{(w,w'):\;C_w2^{-B_m(w)}\equiv C_{w'}2^{-B_m(w')}\ (3^c)\\
B_m(w)\equiv B_m(w')+\tau'-\tau\ (2)}}
p_s(w)\,p_s(w')\,\overline{\Psi_w(a)}\,\Psi_{w'}(a').
\]
The common-endpoint constraint selects word pairs only---no source coset
appears---so the incomplete-coset obstruction of \cref{prop:endpoint} is
structurally absent in this regime, and the Schur test becomes available.
\end{proposition}

\paragraph{Row-sum audit (Schur sharpness).} For the unweighted Schur bound
$\lVert T^m\rVert^2\le R_m:=\sup_{\xi}\sum_{\xi'}\lvert H_m(\xi,\xi')\rvert$
(states $\xi=(a,\tau)$), at $s=0$ and $m=c,\,c+2,\,2c$:
\[
\begin{array}{c|ccc|ccc|ccc}
& \multicolumn{3}{c|}{c=5} & \multicolumn{3}{c|}{c=6}
& \multicolumn{3}{c}{c=7}\\
m & 5 & 7 & 10 & 6 & 8 & 12 & 7 & 9 & 14\\\hline
\sigma_1(T^m)^{1/m} & 0.9231 & 0.8364 & 0.7706
                    & 0.9096 & 0.8373 & 0.7655
                    & 0.9040 & 0.8363 & 0.7443\\
R_m^{1/(2m)}        & 0.9547 & 0.8600 & 0.7849
                    & 0.9350 & 0.8567 & 0.7823
                    & 0.9280 & 0.8620 & 0.7676
\end{array}
\]
The Schur bound tracks $\sigma_1^{1/m}$ within $3\%$ at every $(c,m)$, and
singular-vector weights move it by $\le1\%$ in either direction: the
\emph{unweighted} row sum is essentially sharp, so the Schur norm is the
correct vehicle and no weighted refinement is forced. At $m=2c$ the bound
itself descends in $c$ ($0.7849,\,0.7823,\,0.7676$). The effective constant
$A_m(c)=R_m/\Delta_s^m$ of the anticipated shape $R_m\le A_m\Delta_s^m$
satisfies $A_m^{1/(2m)}=1.61,\,1.49,\,1.33$ at $c=7$, $m=7,9,14$
(descending): the row-multiplicity growth rate is $\approx1.8^m$, comfortably
below the critical rate $3^m$ at which the bound
$\lVert T^m\rVert^{1/m}\le A_m^{1/(2m)}\sqrt{\Delta_s}$ becomes vacuous at
$s=0$.

\paragraph{Degenerate-minimum mass audit (the class-IV test).} The proposed
row-sum decomposition treats by absolute value the class of word pairs whose
\emph{nominal} nonstationarity $d_{\mathrm{nom}}=\min_i[i+v_3(\Delta C_i)]\le
c-2$ degenerates, through mod-$3$ cancellation among the leading terms, to
true pointwise valuation $d_{\mathrm{true}}(a)\ge c-1$---escaping the
vanishing lemma. A Monte Carlo audit ($2\times10^6$ pairs per configuration,
$c=5,6,7$, $m=c,\,c+2,\,2c$; classifier validated against exact complete
sums: machine zeros for nondegenerate pairs, $|S|/\varphi\in[0.5,1]$ for all
classified-degenerate pairs) gives a sharp two-sided verdict. The
\emph{level-cost mechanism is confirmed exactly}: $\Pr[d_{\mathrm{true}}-
d_{\mathrm{nom}}\ge k]$ decays by a factor $3$ per level at every $(c,m)$.
But the \emph{aggregate mass lemma fails in the large-$m$ regime}: the
degenerate mass is flat in $m$ (degeneration is a prefix-local phenomenon),
and its conditional rate among endpoint-matched pairs rises from
$\approx\frac17$ of the unconditional rate at $m=c$ to full decorrelation at
$m=2c$. Consequently the absolute-value bound on the class,
$2\varphi\cdot\Pr[\text{matched}\wedge\text{degenerate}]$, is \emph{not}
$O(\rho^m)$: at $m=2c$ it exceeds the entire measured row sum by factors
$14,\,31,\,77$ for $c=5,6,7$. In the regime $m=c$, however, endpoint matching
\emph{suppresses} degeneracy (the window and the prefix overlap), and the
class-IV mass is only $2$--$4\%$ of the measured row sum, stable-to-decaying
in $c$. The strategic consequence: the Schur bookkeeping must be run at
$\alpha=1$ (i.e.\ $m=c$, the completeness threshold), where a pure mass
treatment of the degenerate class empirically suffices and the row-sum rate
to beat is $R_c^{1/2c}=0.9547,\,0.9350,\,0.9280$ (descending in $c$); at
$\alpha\ge2$ the degenerate class would itself require phase cancellation.
\paragraph{The adjoint-side kernel and the threshold audit.} The
bookkeeping is run on the adjoint-side kernel $G_m=T^m(T^m)^*$, whose
entries carry \emph{complete source sums}: for states $\xi=(x,\sigma)$,
$\xi'=(x',\sigma')$,
\[
G_m(\xi,\xi')
=\sum_{\substack{w\to x,\;w'\to x'\\ B_m(w)-B_m(w')\equiv\sigma-\sigma'\ (2)}}
p_s(w)\,p_s(w')\;\chi(2)^{\sum_i(B_i'-B_i)}
\sum_{a\in U_{3^c}}\chi\!\Bigl(\,\prod_{i=1}^{m-1}
\frac{3^ia+C_i(w)}{3^ia+C_i(w')}\Bigr),
\]
so the coset dichotomy applies \emph{pair by pair}: class I (nonstationary
nondegenerate pairs) contributes exactly zero to every entry. In a full word
enumeration at $c=4,5$ ($m=c$, sup row) all $5.6\times10^5$ resp.\
$8.5\times10^6$ class-I pairs give machine zeros, and the four-class
absolute mass reproduces the matrix row sum from above with $\approx20\%$
internal cancellation margin. The $G$-form row rates at $m=c$ are
\[
R_G^{1/(2c)}=1.0743,\;1.0294,\;1.0003,\;0.9854,\;0.9720
\qquad(c=4,\dots,8),
\]
descending and below $1$ from $c=7$. The exact class decomposition
identifies the controlling class: stationarity at threshold is the
\emph{congruence cascade} $v_3(\Delta C_i)\ge c-1-i$ for every interior $i$,
equivalently the nested conditions $B_i\equiv B_i'\pmod{2\cdot3^{\,c-2-i}}$,
one factor of $3$ deeper per level of $c$. The twisted near-locked pairs
satisfying this cascade (class II) carry $58\%$ resp.\ $48\%$ of the sup-row
mass at $c=4,5$ with per-class rate $1.0373\to0.9762$ (below $1$ from
$c=5$); the locked class III is second ($0.9883\to0.9618$); the degenerate
class IV is subordinate ($0.8406$); the untwisted remainder is negligible
($0.60$).

\subsection*{The near-locked cascade bound}

The class-II combinatorics announced above can be made exact and proved.
Write $t_i=B_i(w)-B_i(w')$ and $D_i=2^{B_i}-2^{B_i'}$, so that
$\Delta C_1=0$ and $\Delta C_{i+1}=3\Delta C_i+D_i$.

\begin{theorem}[Cascade equivalence]\label{thm:cascade}
For a pair of $c$-words, the threshold stationarity cascade
$v_3(\Delta C_i)\ge c-1-i$ (all interior $i$) holds \emph{if and only if}
\[
B_i\equiv B_i'\pmod{2\cdot3^{\,c-3-i}},\qquad 1\le i\le c-3.
\]
\end{theorem}

\begin{proof}
By lifting-the-exponent, $v_3(2^t-1)=0$ for $t$ odd and $=1+v_3(t)$ for $t$
even; since the multiplicative order of $2$ modulo $3^k$ is $2\cdot3^{k-1}$,
$v_3(D_i)\ge k\iff t_i\equiv0\pmod{2\cdot3^{k-1}}$. Forward: if the cascade
holds at $i$ and $i+1$, then
$v_3(D_i)=v_3(\Delta C_{i+1}-3\Delta C_i)\ge\min(c-2-i,\,c-i)=c-2-i$ for
$i\le c-3$. Converse: $\Delta C_2=D_1$ has $v_3\ge c-3$, and inductively
$v_3(\Delta C_{i+1})\ge\min(1+v_3(\Delta C_i),\,v_3(D_i))\ge c-2-i$.
\end{proof}

\begin{lemma}[Cascade mass bound]\label{lem:cascade}
For independent $p_0$-random words,
\[
3^{-(c-3)}\;\le\;\Pr[\textnormal{cascade}]\;\le\;K\cdot3^{-(c-3)},
\qquad
K=1+4\sum_{k\ge0}3^k2^{-2\cdot3^k}<2.2 .
\]
Moreover $\sup_w\Pr_{w'}[\textnormal{cascade}\mid w]\le(2/3)^{c-3}$.
\end{lemma}

\begin{proof}
Lower bound: locking $b_j'=b_j$ for $j\le c-3$ costs
$\prod\Pr[b'=b]=3^{-(c-3)}$. Upper bound: decompose by the first $n$ with
$t_n\ne0$. The event $t_i=0$ for $i<n$ means $b_j'=b_j$ for $j<n$
(probability $3^{-(n-1)}$), and then $t_n=b_n-b_n'$ must be a nonzero
multiple of $M_n=2\cdot3^{\,c-3-n}$:
$\Pr[X\ne0,\,M\mid X]\le\tfrac43 2^{-M}$ for the symmetrized geometric law
$\Pr[X=d]=\tfrac13 2^{-|d|}$. Summing with $k=c-3-n$ gives $K$. Row version:
condition sequentially; given $B'_{i-1}$, the $i$-th congruence pins
$b_i'$ modulo $M_i$, and
$\Pr[b'\equiv r\ (M)]=2^{-\hat r}/(1-2^{-M})\le\tfrac23$.
\end{proof}

\begin{proposition}[Two-point profile of cascade sums]\label{prop:twopoint}
Let $(w,w')$ be a cascade pair and
$R(a)=\prod_{i}(3^ia+C_i)/(3^ia+C_i')$. Then $a\mapsto\chi(R(a))$ factors
through $a\bmod3$, and the complete sum satisfies
\[
\frac{|S(w,w')|}{\varphi}\in\Bigl\{\tfrac12,\;1\Bigr\}.
\]
The value $1$ holds whenever the pair satisfies the \emph{deeper} cascade
$v_3(\Delta C_i)\ge c-i$ (equivalently $B_i\equiv B_i'$ mod
$2\cdot3^{\,c-2-i}$, $i\le c-2$); on boundary pairs the value is generically
$\tfrac12$, with $1$ occurring on a degenerate subclass (measured $\approx14\%$,
the mod-$3$ cancellation channel of the class-IV mechanism). In particular
\emph{no cascade pair vanishes}.
\end{proposition}

\begin{proof}
Each derivative of the log-ratio carries a $\Delta C_i$ factor:
$(\log R)^{(k)}/(k-1)!=\pm\sum_i3^{ik}\bigl[(3^ia+C_i')^{-k}-(3^ia+C_i)^{-k}\bigr]$
has $v_3\ge ik+v_3(\Delta C_i)\ge c-1+i(k-1)\ge c-1$. For $a\equiv a''\bmod3$
the Taylor expansion gives
$v_3(\log R(a)-\log R(a''))\ge\min_k\bigl[k+(c-1)-v_3(k!)\bigr]\ge c$
(as $k-v_3(k!)\ge1$), so $\chi(R(a))=\chi(R(a''))$. Hence
$S=\tfrac{\varphi}2\bigl[\chi(R(a_1))+\chi(R(a_2))\bigr]$ with
$\chi(R(a_1))/\chi(R(a_2))=\omega^\lambda$ a cube root of unity, and
$|1+\omega^\lambda|\in\{2,1\}$. For deeper pairs the same argument one level
down makes $\chi\circ R$ constant on all of $U$. Machine verification at
$c=5$: $600/600$ boundary pairs exactly in $\{\tfrac12,1\}$
($86\%/14\%$), $800/800$ deeper pairs exactly $1$.
\end{proof}

\begin{remark}[Assembly status]
\cref{lem:cascade} proves the class-II \emph{count} is geometrically thin
($K3^{-(c-3)}$, sharp up to constant), and \cref{prop:twopoint} caps the
surviving sums at $\varphi/2$ off the deeper subclass. What the crude
assembly does not yet deliver is the full budget $R_{\mathrm{II}}\le
\theta^{2c}$: for rows of word-mass $\asymp2^{-c}$ (the fixed-point residue
$x=-1$, which is the empirical sup row), the product
$\varphi\cdot\mu_{\mathrm{row}}\cdot(2/3)^{c-3}$ has exact growth rate
$\bigl(3\cdot\tfrac12\cdot\tfrac23\bigr)^c=1$: the ninth appearance of the
marginal coincidence in this program. The obstruction is concentrated on
adversarially aligned rows---at $x=-1$ every tail congruence targets the
heavy residue $\hat r=1$, realizing the worst case of the sequential bound.
The remaining sharpening is the off-representative (near-diagonal)
multiplicity bookkeeping with the $\{\tfrac12,1\}$ profile carried exactly.
\end{remark}

\paragraph{Heavy-row cascade audit: the marginality is broken.} The
discriminating experiment (exact word enumeration on the row $x=-1$ at
$m=c$, with the two-point formula
$S=\tfrac{\varphi}2[\chi(R(1))+\chi(R(2))]\cdot\chi(2)^{\Delta\Sigma B}$
giving every cascade sum in closed form, verified to $10^{-13}$) resolves
the trichotomy in favour of \emph{mass-plus-profile bookkeeping}. Across
$c=5,6,7$, every component of the absolute cascade row budget decays by a
stable factor $0.72$--$0.74$ per level: the total
$R^{\mathrm{abs}}_{\mathrm{casc}}(-1)=1.468\to1.063\to0.783$, hence
\[
\Lambda_c=R^{1/(2c)}=1.0392,\;1.0051,\;\mathbf{0.9827},
\]
\emph{crossing below $1$ at $c=7$} exactly as extrapolated, and tending to
$\approx\sqrt{0.73}\approx0.85$ if the per-level ratio is sustained. The
signed-to-absolute ratio \emph{rises} toward $1$
($0.876\to0.915\to0.961$): there is no hidden phase coherence inside the
cascade, and none is needed. The $\beta{=}1$ channel carries a remarkably
stable $38.8\%$ of the class-II mass at all three levels (bounded well
away from saturation, and consistent with \cref{lem:betachannel}); the
off-representative mass collapses geometrically
($3.6\times10^{-3}\to8.0\times10^{-4}\to1.9\times10^{-4}$): multiplicity
is $O(1)$, not $3^{\epsilon c}$. At $c=7$ the heavy row $x=-1$ is
confirmed as the sup row ($16.7\times$ the mean row mass), and the
transition mass within it is \emph{spread}: the largest single column
carries $12\%$ of $R_{\mathrm{casc}}$ and the top ten only $43\%$---the
favourable case for a Perron contraction argument, since no single state
must be controlled exactly. The crude marginal coincidence
$(3\cdot\frac12\cdot\frac23)^c=1$ is thus broken by the true profile and
multiplicity constants. One of the three savings is already a lemma, by applying
\cref{lem:cascade} at two adjacent depths:

\begin{lemma}[Deeper-channel bound]\label{lem:betachannel}
The deeper cascade ($v_3(\Delta C_i)\ge c-i$, moduli $2\cdot3^{\,c-2-i}$,
$i\le c-2$) satisfies the same mass bound one level up,
$\Pr[\textnormal{deeper}]\le K\cdot3^{-(c-2)}$, with the same constant
$K<2.2$. Hence, conditional on the cascade,
\[
p_1=\Pr[\textnormal{deeper}\mid\textnormal{cascade}]
\le\frac{K\cdot3^{-(c-2)}}{3^{-(c-3)}}=\frac K3<0.733,
\qquad
\E[\beta\mid\textnormal{cascade}]
\le\frac12\Bigl(1+\frac K3\Bigr)+\frac q2<0.867+\frac q2,
\]
where $q$ is the degenerate boundary channel mass
($\lambda=0$ without the deeper congruences; measured $\approx0.14$,
stable). This is a strict uniform saving over the worst case
$\beta\equiv1$; it is the aggregate form of the per-state channel bound
required by the contraction program, and the channel constant $q$ is the
one remaining unproved ingredient of the $\beta$-profile.
\end{lemma}

In fact the central counting estimate can be made sharp, and doing so
\emph{dissolves} the marginal coincidence:

\begin{lemma}[Sharp sequential pinning; off-representative absorption]
\label{lem:pinning}
For every $m$-word $w$,
\[
\Pr_{w'}\bigl[(w,w')\textnormal{ cascade}\bigr]
\le C_0\,2^{-(c-3)},
\qquad
C_0=\prod_{j\ge0}\bigl(1-2^{-2\cdot3^j}\bigr)^{-1}
=\tfrac43\cdot1.0159\cdots<1.355 .
\]
\end{lemma}

\begin{proof}
Condition sequentially on $b_1',\dots$. At level $i$ the cascade congruence
$B_i'\equiv B_i\pmod{M_i}$ pins $b_i'$ to a residue class
$\hat r_i\in\{1,\dots,M_i\}$, and
$\Pr[b'\equiv\hat r\ (M)]=2^{-\hat r}(1-2^{-M})^{-1}
\le\tfrac12(1-2^{-M})^{-1}$. The moduli $M_i=2\cdot3^{\,c-3-i}$ traverse
$\{2\cdot3^j:j=0,\dots,c-4\}$, and the product of the correction factors
converges to $C_0$.
\end{proof}

Two consequences. \emph{First}, the off-representative mass (Lemma~B of the
contraction program) is absorbed: the factors $(1-2^{-M})^{-1}$ \emph{are}
the off-representative sums, so off-representatives contribute at most the
fixed total factor $C_0$, uniformly in $c$---no separate lemma is needed.
\emph{Second}, the marginal coincidence of the assembly remark was an
artifact: charging the worst case $\tfrac23$ at every level is wrong, since
the level-$i$ cost exceeds $\tfrac12$ only by the $2^{-M_i}$ corrections,
and the value $\tfrac23$ occurs solely at the single level with $M=2$.
Carrying \cref{lem:pinning} through the row assembly gives genuine
contraction:

\begin{proposition}[Cascade row contraction]\label{prop:rowcontract}
For every row state $\xi$ over the residue $x$, the cascade (class
II${}\cup{}$III) contribution to the $G_c$ Schur row sum satisfies
\[
R_{\mathrm{casc}}(\xi)\;\le\;\varphi\cdot\mu(x)\cdot C_0\,2^{-(c-3)},
\]
where $\mu(x)$ is the endpoint word mass of the row. Consequently:
\emph{(a)} on the heavy row, where $\mu(-1)=A\cdot2^{-c}$ with
$A\approx1.45$ measured stable over $c=4,\dots,7$,
\[
R_{\mathrm{casc}}(-1)\le\tfrac{16}{3}C_0A\,(3/4)^c\approx10.5\,(3/4)^c,
\]
matching the audited per-level decay $0.72$--$0.74$ and $\Lambda_7=0.9827$;
\emph{(b)} for all rows simultaneously,
$\mu(x)\le\bigl(\sum_x\mu(x)^2\bigr)^{1/2}=\sqrt{\mathrm{Col}_c(c)}$, so
granted the threshold flattening input
$\mathrm{Col}_c(c)\le\mathrm{poly}(c)\,3^{-c}$ (proved for $m\ge(1+\epsilon)c$
in \cref{lem:flat}; measured $3^c\,\mathrm{Col}\approx4$ at $m=c$),
\[
\sup_\xi R_{\mathrm{casc}}(\xi)\le\mathrm{poly}(c)\Bigl(\tfrac{\sqrt3}2\Bigr)^{c},
\qquad\text{i.e.}\qquad
\Lambda_c\le\frac{3^{1/4}}{\sqrt2}+o(1)\approx0.9306 .
\]
The class-II/III budget of the threshold Schur bound is therefore
contractive, \emph{modulo a single analytic input}: the $m=c$ extension of
the flattening lemma. (The $\beta$-profile and the channel constant $q$ are
no longer needed for contraction---$\beta\le1$ suffices---and survive only
as sharpening constants.)
\end{proposition}

\paragraph{The cascade operator, numerically.} In the scaled coordinates
$\bar t_{i+1}=3\bar t_i+\bar x$ the heavy-row chain becomes an explicit
two-channel transfer operator ($\varphi$-growth $3$, $w$-side mass
$\tfrac12$, $w'$-side weight $2^{-(1+M_i|\bar x|)}$), computable to any $c$
from $\approx50$-dimensional level matrices. Its bulk per-level Perron
radius is \emph{exactly} $3/4=3\cdot\tfrac12\cdot\tfrac12$, with Perron
vector the point mass at $\bar t=0$: the locked state is the unique
recurrent state, all off-representative states are transient, and the
``heavy-state nonpersistence'' lemma of the contraction program dissolves
in these coordinates---the heavy state \emph{is} the Perron state and
contracts at $3/4$ per level. The model's lock share is $1/C_0=0.7383$
(internal consistency with \cref{lem:pinning}), and its row rates
$\Lambda^{\mathrm{model}}_c=1.0414,\,1.0099,\,0.9879$ at $c=5,6,7$ agree
with the measured $1.0392,\,1.0051,\,0.9827$ to $0.5\%$, descending through
$0.893$ at $c=30$ toward $\sqrt{3/4}=0.866$.

\subsection*{Threshold flattening: the exact Fourier reduction}

The single analytic input of \cref{prop:rowcontract} is the $m=c$ extension
of \cref{lem:flat}. Locating the role of $\varepsilon$ in the proved
$m\ge(1+\varepsilon)c$ case: it is used exactly once, to decouple the walk
factor $2^{B_{m-1}}$ from the window via $\varepsilon c$ fresh prefix
symbols. At threshold the prefix is empty---and indeed the $L^\infty$ form
of flattening is \emph{false} there: the fiber $x=-1$ carries mass
$\mu(-1)=A\cdot2^{-c}\gg\mathrm{poly}(c)3^{-c}$, so no fiber bound can
hold, and only the $L^2$ (collision) statement survives. That statement
has an exact Fourier form:

\begin{proposition}[Plancherel form of threshold flattening]\label{prop:plancherel}
With $\hat E_c(u)=\E_w\bigl[e\bigl(uC_w/3^c\bigr)\bigr]$ the renewal
Fourier coefficients of the $c$-step constant,
\[
\operatorname{Col}_c(c)=3^{-c}\sum_{u\bmod 3^c}\bigl|\hat E_c(u)\bigr|^2 .
\]
Threshold flattening is therefore \emph{equivalent} to
$\sum_u|\hat E_c(u)|^2\le\mathrm{poly}(c)$. Computed exactly for
$c=3,\dots,6$ the sum is $3.32,\,3.53,\,3.70,\,3.90$---\emph{linear},
$\approx0.19c+2.8$---and the sampled per-coefficient decay is
$\max_u|\hat E_c(u)|\approx\mathrm{poly}(c)\,3^{-c/2}$.
\end{proposition}

\begin{lemma}[Marginal scale]\label{lem:marginalscale}
Revealing one symbol multiplies $\hat E_c(u)$ by a conditional factor
bounded by $\eta_k=\max_{3\nmid a}\bigl|\sum_{b\ge1}2^{-b}
e\bigl(a2^b/3^k\bigr)\bigr|$ at the active scale $k$, and
\[
\eta_1=\tfrac1{\sqrt3}\qquad\text{exactly}:
\]
$2^b\equiv(-1)^b\pmod3$ gives $S=\tfrac23e_3(-a)+\tfrac13e_3(a)$ and
$|S|^2=\tfrac59-\tfrac49\cdot\tfrac12=\tfrac13$. Thus $3\eta_1^2=1$: the
shallowest scale is \emph{exactly marginal}, which is why the Plancherel
sum grows linearly (one marginal constant per level) rather than staying
bounded---and why no estimate can beat $\mathrm{poly}(c)3^{-c}$ by more
than the polynomial.
\end{lemma}

\paragraph{The obstruction and the mechanism.} The worst-case per-scale
product fails structurally: $\eta_k\uparrow1$
($0.577,\,0.647,\,0.835,\,0.930,\dots$), since at deep scales the aligning
$a$ places the first $\sim k\log_32$ orbit phases near zero. The decay is
therefore \emph{path-averaged}: the phase parameter evolves by the
expanding dynamics $\theta_{j+1}=3\cdot2^{b_j}\theta_j\bmod1$, so the
undamped tiny-phase region is transient---a phase of size $3^{-d}$ escapes
within $O(d)$ steps and is then damped at $\eta_1$-strength. The measured
per-$u$ decay $\mathrm{poly}\cdot3^{-c/2}$ and the linear Plancherel sum
confirm this mechanism quantitatively. The last analytic input of the
finite side is thus:

Two structural results sharpen the target. First, the statement is
\emph{self-similar across levels}, and in particular $c$-uniform:

\begin{lemma}[Level decoupling; self-similar shells]\label{lem:shells}
For $1\le L\le c$,
$C_c\bmod3^L=2^{B_{c-L}}\cdot\tilde C_L$, where $\tilde C_L$ is an
$L$-term constant built from the \emph{last} $L-1$ increments and the unit
$2^{B_{c-L}}$ is independent of it: at proper levels the unit prefactor
decouples (the threshold coupling concerns only the full level $L=c$).
Consequently, by Cauchy--Schwarz over the unit twist, the level-$L$
Plancherel shell satisfies
\[
\sum_{\operatorname{lev}(u)=L}\bigl|\hat E_c(u)\bigr|^2
\;\le\;\tilde Q_L:=3^L\operatorname{Col}_L(L),
\qquad\text{hence}\qquad
3^c\operatorname{Col}_c(c)\le1+\sum_{L=1}^{c}\tilde Q_L .
\]
Thus threshold flattening reduces to the \emph{single-scale} sequence
$\tilde Q_L\le\mathrm{poly}(L)$, uniformly in the ambient $c$, and even
tolerates a quadratic loss.
\end{lemma}

\begin{proof}
Terms with $c-1-j\ge L$ vanish mod $3^L$;
$\sum_{j\ge c-L}3^{c-1-j}2^{B_j}=2^{B_{c-L}}\sum_{i<L}3^{L-1-i}
2^{B_{c-L+i}-B_{c-L}}$, and the prefix unit involves only $b_1,\dots,
b_{c-L}$. Then
$\Pr[UC\equiv U'C']=\E_\rho\Pr[C\equiv\rho C']\le
\lVert\mu\rVert_2\lVert\mu_\rho\rVert_2=\operatorname{Col}_L(L)$ since
rotation by a unit preserves the $\ell^2$ norm of the law.
\end{proof}

Second, the computed shells identify the \emph{correct strength} of the
remaining statement and disprove the natural stronger route. Computing
$\hat E_L(u)$ exactly for \emph{all} $u$ at $L\le7$: the primitive shells
are
\[
\sum_{\operatorname{lev}(u)=L}|\hat E_L(u)|^2
=2.00,\;1.43,\;1.41,\;1.39,\;1.37,\;1.41,\;1.38
\qquad(L=1,\dots,7),
\]
\emph{constant} $\approx1.4$---so conjecturally $\tilde Q_L=O(L)$ with
bounded shells. But the levelwise \emph{maximum} fails:
$3^L\max_u|\hat E_L(u)|^2$ grows like $1.75^L$, and the maximizers
converge to the \emph{dyadic-resonant direction} $u/3^L\to3/4$
($0.481,\,0.247,\,0.749,\,0.750,\,0.7499$ for $L=3,\dots,7$)---the
Fourier twin of the heavy fiber $x=-1$. As with the fiber bound, the
$L^\infty$ form is false and only the $L^2$ aggregate is true: the proof
must bound the shell as a whole (a second moment over the level), not
coefficient by coefficient. Transience statistics point the same way: the
expected occupation of the weak-damping region along the phase chain is
$\approx0.2L$ (linear with small constant, not $O(\log L)$), so the
accounting must use the full multiplier distribution rather than a
weak/strong dichotomy. These constraints having been respected, the
statement can in fact be \emph{proved}, by exposing the collision digit by
digit in the terminal-window filtration.

\begin{theorem}[Threshold flattening: bounded shells]\label{op:flat}
For every $L\ge1$,
\[
\tilde Q_L=3^L\Pr\bigl[C_L\equiv C_L'\bmod3^L\bigr]
\;\le\;3\,C_1,
\qquad
C_1=\prod_{j\ge0}\frac{1+2^{-2\cdot3^j}}{1-2^{-2\cdot3^j}}<1.7196 ,
\]
so the shells are \emph{bounded} ($3C_1<5.16$; measured values
$3.00,\dots,4.04$ for $L\le7$). Consequently, by \cref{lem:shells},
$3^c\operatorname{Col}_c(c)\le1+3C_1c$: threshold flattening holds with a
\emph{linear} polynomial, matching the measured $0.19c+2.8$.
\end{theorem}

\begin{proof}
Write $d_j=b_j-b_j'$, $S_i=\sum_{j\le i}d_j$, and let
$X_k$ denote the terminal window of level $k$ (last $k-1$ increment
pairs), so that $C_L\equiv2^{B_{L-k}}X_k\pmod{3^k}$
(\cref{lem:shells}) with the renewal recursion
$X_{k+1}=1+3\cdot2^{-b_{L-k}}X_k$ and $X_k\equiv1\bmod3$.

\emph{(i) Nested conditions.} Since $2$ generates the units modulo $3^k$
(order $M_k=2\cdot3^{k-1}$), the level-$k$ collision
$E_k=\{C_L\equiv C_L'\bmod3^k\}$ is equivalent to
\[
S_{L-k}\equiv r_k\pmod{M_k},
\qquad
r_k:=\operatorname{dlog}_{3^k}\bigl(X_k'X_k^{-1}\bigr),
\]
and $E_L=\bigcap_{k\le L}E_k$.

\emph{(ii) Differencing.} For $i=1,\dots,L-1$ put $k=L-i$ and
$N_i:=M_{L-i}=2\cdot3^{L-i-1}$. Subtracting the conditions at levels $k$
and $k+1$ (whose modulus is $3N_i$), on $E_L$ the revealed pair $i$
satisfies the event
\[
D_i:\qquad
d_i\equiv r_{L-i}-r_{L-i+1}\pmod{N_i},
\qquad\text{so}\qquad
E_L\subseteq\bigcap_{i=1}^{L-1}D_i .
\]

\emph{(iii) Adaptedness and separability.} Let
$\mathcal F_i=\sigma\bigl((b_j,b_j'):j>i\bigr)$. Then
$r_{L-i}\in\mathcal F_i$, while $r_{L-i+1}$ depends on $(b_i,b_i')$ only
through the window extensions, \emph{separably}: $D_i$ is the event
\[
F(b_i)-F'(b_i')\equiv\rho_i\pmod{N_i},
\qquad
F(b):=b+\operatorname{dlog}_{3^{L-i+1}}\bigl(1+3\cdot2^{-b}X_{L-i}\bigr),
\]
with $\rho_i\in\mathcal F_i$ and $F'$ built from $X'_{L-i}$.

\emph{(iv) Class bijectivity.} $F(b)\equiv F(\tilde b)\bmod N_i$ if and
only if $b\equiv\tilde b\bmod N_i$. Indeed, with $e=\tilde b-b$, the
ratio of the two window extensions is the principal unit
$1+3X(2^{-b}-2^{-\tilde b})/(1+3\cdot2^{-\tilde b}X)$ of depth
$t=1+v_3(2^{e}-1)$; by lifting-the-exponent $t=1$ for $e$ odd and
$t=2+v_3(e)$ for $e$ even, and the principal-unit index formula
$\operatorname{dlog}(1+3^t v)=2\cdot3^{\,t-1}v'$ places the perturbation
of $F(b)-F(\tilde b)$ \emph{one $3$-adic level deeper} than $e$: for odd
$e$ the difference is odd (distinct mod $2$); for even $e$,
$v_3\bigl(F(b)-F(\tilde b)\bigr)=v_3(e)$ exactly, with both parts even.
Hence $F(b)\equiv F(\tilde b)\bmod2\cdot3^m\iff e\equiv0\bmod
2\cdot3^m$. (Machine-verified on $7\times10^6$ pairs across
$k\le5$ and random windows: no exceptions.)

\emph{(v) Per-step bound.} Given $\mathcal F_i$, the pair $(b_i,b_i')$ is
fresh with law $2^{-b}2^{-b'}$. By (iv) the pushforward class masses of
$F$ (and of $F'$) modulo $N_i$ are a \emph{permutation} of the geometric
class masses $m_v=2^{-v}/(1-2^{-N_i})$, $v=1,\dots,N_i$. By
Cauchy--Schwarz, for \emph{any} alignment $\rho_i$,
\[
\Pr[D_i\mid\mathcal F_i]
=\sum_v m^{F}_v\,m^{F'}_{v-\rho_i}
\;\le\;\sum_v m_v^2
=\frac13\,\frac{1+2^{-N_i}}{1-2^{-N_i}} .
\]

\emph{(vi) Telescoping.} $D_{i'}\in\mathcal F_i$ for $i'>i$, so peeling
$i=1,2,\dots$ inward,
\[
\Pr[E_L]\le\prod_{i=1}^{L-1}
\frac13\,\frac{1+2^{-N_i}}{1-2^{-N_i}}
=3^{-(L-1)}\prod_{j=0}^{L-2}\frac{1+2^{-2\cdot3^j}}{1-2^{-2\cdot3^j}}
\le3^{-(L-1)}C_1 . \qedhere
\]
\end{proof}

\begin{remark}
The proof realizes the energy principle exactly: each additional $3$-adic
digit of collision costs $\tfrac13$ in pair energy, up to dyadic wrap
corrections whose product converges ($C_1$); the heavy fiber and the
dyadic-resonant Fourier direction---which falsify the $L^\infty$
forms---are paid for automatically by the quadratic pair weights. The
per-step constant is sharp (attained numerically to six digits), and the
final bound $3C_1=5.159$ sits within $25\%$ of the measured shells. The
statement and proof are annealed and character-free: no $\chi$,
coboundary, or orbit-deterministic input appears. All of this is at
$s=0$; the $s\in[0,1]$ extension (lighter tails) is expected by the same
scheme but not claimed.
\end{remark}

\paragraph{Consequences.} \cref{prop:rowcontract} is now
\emph{unconditional}: for every row state $\xi$,
$R_{\mathrm{casc}}(\xi)\le\mathrm{poly}(c)\,(\sqrt3/2)^c$, i.e.\ the
cascade (class I--III) sector of the threshold Schur bound is a
\emph{theorem}, with rate constant $3^{1/4}/\sqrt2\approx0.9306$. The
sole remaining budget is class IV, treated next.

\subsection*{The class-IV budget and the final assembly}

Class IV consists of pairs whose nominal valuation
$d_{\mathrm{nom}}=\min_i v_i$, $v_i=i+v_3(\Delta C_i)$, is $\le c-2$ but
whose true pointwise valuation reaches $\ge c-1$ through cancellation.
Two of the three ingredients controlling it are proved.

\begin{lemma}[Nominal-shell mass]\label{lem:nomshell}
$\Pr[d_{\mathrm{nom}}\ge d]\le K\cdot3^{-(d-2)}$ with the constant
$K<2.2$ of \cref{lem:cascade}.
\end{lemma}
\begin{proof}
$d_{\mathrm{nom}}\ge d$ is the depth-$d$ cascade
$v_3(\Delta C_i)\ge d-i$ $(\forall i)$, equivalent by the proof of
\cref{thm:cascade} to $B_i\equiv B_i'\ (\mathrm{mod}\ 2\cdot3^{\,d-2-i})$
for $i\le d-2$; the first-decoupling argument of \cref{lem:cascade}
applies verbatim.
\end{proof}

\begin{lemma}[Leading coefficient: criterion, $a$-independence,
singleton rigidity]\label{lem:leadcoeff}
Write $\mathcal L(a)=\sum_i3^i\Delta C_i\,u_i(a)$ for the pair
log-derivative, $I$ for the set of indices attaining $d_{\mathrm{nom}}$,
and
\[
\Lambda:=\sum_{i\in I}\delta_i\nu_i\ \bmod3,
\qquad
\delta_i=\text{leading digit of }\Delta C_i,\quad
\nu_i=(C_iC_i')^{-1}\bmod3 .
\]
Then: (i) $d_{\mathrm{true}}(a)\ge d_{\mathrm{nom}}+1$ iff $\Lambda=0$;
(ii) $\Lambda$ is \emph{independent of $a$} (as $3^ia\equiv0\bmod3$), so
the first degeneracy level is a pure word-pair event; (iii) if $|I|=1$
then $\Lambda\ne0$: \emph{singleton shells cannot degenerate}---class-IV
escape requires a collision of minima. Machine verification at $c=7$
($2\times10^6$ pairs): the first level is $a$-independent without
exception; $\Pr[h\ge1\mid d_{\mathrm{nom}}=d]\approx0.13$ uniformly over
$d=3,\dots,6$ and \emph{exactly} $0$ at the singleton-dominated shell
$d=2$ ($8.9\times10^5$ pairs); the conditional level costs are
$\Pr[h\ge2\mid h\ge1]=0.302$, $\Pr[h\ge3\mid h\ge2]=0.333$.
\end{lemma}

\begin{theorem}[One-condition degeneracy cost; level-cost law]
\label{op:onecond}
Fix a unit $a$, a word $w$, and $d\ge2$. Then, uniformly over $(w,a)$,
\[
\Pr_{w'}\bigl[d_{\mathrm{nom}}\ge d\ \wedge\ h(a)\ge r\bigr]
\;\le\;C_0\,2^{-(d-3)}\cdot\frac23\Bigl(\frac47\Bigr)^{r-1},
\]
with $C_0<1.355$ the pinning constant of \cref{lem:pinning}. The
level-cost constant is $\lambda=\frac47=0.5714\ldots<3^{-1/2}$.
\end{theorem}

\begin{proof}
Reveal $w'$ forward and let $\mathcal G_t=\sigma(b_1',\dots,b_t')$. The
scale-$t$ digit of $\mathcal L(a)$ requires $\Delta C_i$ modulo
$3^{\,t-i+1}$ for $i\le t$ only, so it is settled at time $t-1$, and the
freshest variable $\beta:=b'_{t-1}$ enters solely through
$C'_t=3C'_{t-1}+2^{B'_{t-2}}2^{\beta}$.

\emph{(i) Support two; trichotomy.} At scale $t$ itself,
$2^{\beta}\bmod3=(-1)^{\beta}$, so
$\mathrm{digit}_t=W_t+\kappa_t(-1)^{\beta}$ with $W_t$ predictable and
$\kappa_t=\pm2^{B'_{t-2}}u_t(a)$ a predictable unit. Hence
$\Pr[\mathrm{digit}_t=0\mid\text{past}]\in\{0,\tfrac13,\tfrac23\}$
according as $W_t\equiv0,-\kappa_t,+\kappa_t$; in particular $\le\frac23$.

\emph{(ii) Three-value spread at the next scale.} At scale $t+1$ the same
$\beta$ enters through $2^{\beta}\bmod9$. Within a parity class the orbit
triple is $\{2,8,5\}$ (odd) or $\{1,4,7\}$ (even): \emph{distinct
modulo $9$, congruent modulo $3$}. Hence, conditional on the parity class
selected at scale $t$ and on \emph{any} realization of all other
variables, the three mod-$6$ refinement classes $j$ of $\beta$ shift the
scale-$(t{+}1)$ base $W_{t+1}$ by three distinct second digits
($3\cdot\{0,1,2\}\cdot$unit), so $W_{t+1}(j)$ runs over \emph{all three}
residues bijectively, and by (i) the associated costs are a permutation of
$(0,\tfrac13,\tfrac23)$. The within-class masses of $j$ are geometric,
exactly $(\tfrac{16}{21},\tfrac4{21},\tfrac1{21})$ for both parities, so
\[
\Pr[\mathrm{digit}_{t+1}=0\mid\text{class, past}]
\le\frac{16}{21}\cdot\frac23+\frac4{21}\cdot\frac13+\frac1{21}\cdot0
=\frac47 .
\]

\emph{(iii) Chain.} The lift $h(a)\ge r$ requires
$\mathrm{digit}_t=0$ for $t=d,\dots,d+r-1$; the first scale pays
$\le\frac23$ by (i) and each subsequent scale pays $\le\frac47$ by (ii)
(each $\beta$ serves once as fresh parity and once as refinement;
conditioning on more and bounding uniformly is legitimate). 

\emph{(iv) Shell factorization.} The depth-$d$ cascade
$d_{\mathrm{nom}}\ge d$ is equivalent to congruences on the pairs
$j\le d-2$ (\cref{thm:cascade} at depth $d$), disjoint from the peeling
variables $b'_j$, $j\ge d-1$; the fixed-$w$ sequential pinning
(\cref{lem:pinning} at depth $d$) gives $C_0\,2^{-(d-3)}$, and the
uniform conditional bounds (i)--(iii) multiply against it.
\end{proof}

\begin{remark}
The trichotomy in (i) is exactly the measured structure ($\Pr[h\ge1]$
per shell $\approx0.13$, conditionals $0.30$--$0.33$), and the bound is
honest under adversarial conditioning: at $w=1^c$ (the heavy word) with
$a=1$ the measured conditional level costs rise to
$0.380,\,0.429,\,0.444$---approaching, but respecting, $4/7$. The
inequality $\frac47<3^{-1/2}$ is what the assembly requires:
$\sqrt3\cdot\frac47=0.9897<1$.
\end{remark}

\subsection*{Cellwise summation: the quadratic route}

The remaining step is harmonic analysis on the degenerate cells, and its
main mechanism is now proved. Differentiating termwise,
\[
\mathcal L'(a)=-\sum_i\mathrm{term}_i(\mathcal L)(a)\cdot g_i(a),
\qquad
g_i(a)=3^i\,\bigl(2\cdot3^ia+C_i+C_i'\bigr)\,u_i(a):
\]
each second-derivative term is the corresponding first-derivative term
times an extra factor of valuation $i+v_3(2\cdot3^ia+C_i+C_i')$.

\begin{lemma}[Quadratic valuation; unique minimum]\label{lem:quadval}
For a class-IV pair with nominal depth $d$, active set $I$, and
$i_{\min}=\min I$,
\[
v_3\bigl(\mathcal L'(a)\bigr)=d+i_{\min}\qquad\text{exactly},
\]
unless one of two explicit mod-$3$ coincidences occurs: \emph{(E1)} the
sum-factor exception $C_{i_{\min}}+C_{i_{\min}}'\equiv0\bmod 3$
(possible only when $v_3(\Delta C_{i_{\min}})=0$), or \emph{(E2)} an
inactive tie at $i_{\min}-1$ with $v_{i_{\min}-1}=d+1$ exactly and
leading-digit cancellation. The reason is the $3^i$-grading: on the active
set the $\mathcal L'$-term valuations are $d+i$, \emph{strictly
increasing in $i$}, so the minimum is unique and cancellation---the very
mechanism that created the degeneracy in $\mathcal L$---is impossible in
$\mathcal L'$ (for $v_3(\Delta C_i)\ge1$ the sum factor is
$\equiv2C_i\not\equiv0$ automatically). Audit at $c=7$: $300/300$
class-IV points have $v_3(\mathcal L')=d+i_{\min}$, no exceptions
observed.
\end{lemma}

On a Postnikov cell $a=a_0+3^qt$ where the linear term is invisible
($q+v_3(\mathcal L)\ge c$) and the quadratic term is at scale
$2q+v_3(\mathcal L')=c-\ell'$ with unit coefficient, the cell sum is an
exact quadratic Gauss sum:
$\bigl|\sum_{t\bmod3^{\ell}}e_{3^{\ell}}(Bt^2+At)\bigr|
\in\{0,3^{\ell/2}\}$ for $3\nmid B$---the square-root saving, with
constant $1$. By \cref{lem:quadval} the quadratic scale is \emph{known}
($w_2=d+i_{\min}$), so the boundary cells of the degenerate set carry
Gauss factors $3^{\ell/2}$, converting the level-cost mass
$(4/7)^r$ of \cref{op:onecond} into
\[
\Bigl(\sqrt3\cdot\frac47\Bigr)^{r}=(0.98974\ldots)^{r},
\]
summable: precisely the final margin. What remains is bookkeeping for the
exceptional channels:

\begin{lemma}[Linear-regime stratum vanishing]\label{lem:stratvan}
For a pair of nominal depth $d$, every stratum
$\{a:\nu(a)=\nu_0\}$ with $\nu_0\le(c+d-1)/2$ contributes \emph{zero}
to $S(w,w')$: the stratum is a union of cells of size $3^{\,c-\nu_0+d}$
(its defining digit conditions involve $a$ only modulo
$3^{\,\nu_0-d}$), each cell admits complete linear subsums at scale
$3^{\,\nu_0+1}\le3^{\,c-\nu_0+d}$, and the quadratic correction is
deeper by \cref{lem:quadval}. 
\end{lemma}

\paragraph{Two structural corrections from the audit.} Direct computation
at $c=7$ corrects the na\"ive Gauss framing of the remaining step in two
ways. \emph{(A)} On genuinely degenerate (class-IV) pairs there is
\emph{no} summation cancellation to harvest: the deep set covers
$\approx72\%$ of $U$ and $|S_{\mathrm{deep}}|/m_{\mathrm{deep}}\approx
0.9$--$1.0$ ($200$ pairs)---the phase is essentially constant where the
pair is degenerate, so class-IV pairs are controlled \emph{purely
probabilistically}, by \cref{op:onecond} at fixed $a$ via Fubini.
\emph{(B)} The mid-strata $(c+d)/2<\nu<c-1$---the only place a
Gauss-type trade could act---are \emph{empty} at $c=7$ ($0$ of $60{,}000$
sampled pairs): partly small-$c$ arithmetic (the window contains no
integer for $d\ge3$), partly singleton rigidity at $d=2$
(\cref{lem:leadcoeff}). The valuation profile is \emph{bimodal}: shallow
(exact vanishing, \cref{lem:stratvan}) or fully degenerate (measure). The
remaining statement is therefore not an exceptional-channel escalation but
a \emph{mid-strata bound}:

The $c=9$ audit (the first genuine mid-window; $2500$ pairs, full
$\nu$-spectrum over all $13122$ units) resolves the question and
identifies the mechanism. The measured Fubini density obeys
\[
\E_{w'}\bigl[\#\{a:\nu(a)=d+h\}\bigr]/\varphi
\approx0.3\cdot3^{-h},
\]
with successive ratios $0.31,0.29,0.33,0.33,0.33,0.33$---\emph{exactly}
$\tfrac13$ per level, far below the worst-case $(4/7)^h$ (by a factor
$50$ at $h=7$): the $a$-averaged level cost is the averaged trichotomy
$\tfrac13$. The structural reason: clearing the unit-valued denominators,
$\nu(a)=v_3(N(a))$ for a \emph{fixed integer polynomial} $N$ of degree
$\le2(c-2)$, so the deep strata are contained in balls around the
$3$-adic roots of $N$, of total density $\le\deg(N)\cdot3^{-h}$ up to
root-cluster corrections---the root-ball law the audit measures.

\begin{lemma}[All-derivative grading; Gauss on root-balls]
\label{lem:rootball}
The unique-minimum argument of \cref{lem:quadval} applies to every
derivative: $w_k:=v_3(\mathcal L^{(k)}(a))=d+(k-1)\,i_{\min}$ off the
explicit coincidence channels. Consequently, on a root-ball
$a=\alpha+3^{\rho}t$ the Taylor term of degree $k+1$ has valuation
$k\rho+w_k$, strictly increasing in $k$ for $k\ge2$: the phase on every
root-ball is \emph{exactly quadratic} with coefficient of known valuation
$w_2=d+i_{\min}$, the cubic term is always deeper, and the complete
quadratic Gauss bound gives
$\bigl|\sum_{a\in\mathcal B}\chi(R(a))\bigr|
\le\sqrt3\,|\mathcal B|^{1/2}$. This---not the strata cells---is where
the quadratic route operates.
\end{lemma}

Assembling: $|S(w,w')|\le0$ (shallow strata, \cref{lem:stratvan}) plus
$\sum_{\mathcal B}\sqrt3|\mathcal B|^{1/2}
\le\mathrm{poly}(c)\,3^{(c+d)/4}$ over the $\le\deg N$ root-balls
reaching the half-threshold, giving a class-IV row budget
$R_{\mathrm{IV}}(\xi)\le\mathrm{poly}(c)\,3^{-c/4}$---exponentially
subordinate with a comfortable margin, replacing the razor-thin
$0.9897^c$ of the worst-case route. What remains is standard but must be
done carefully:

\paragraph{The root-ball bookkeeping.} We now carry this out. Write
$w_2(a)=v_3(\mathcal L'(a))$ and $\tilde w=w_2-d$.

\begin{lemma}[Derivative-excess ladder]\label{lem:ladder}
$w_2(a)$ is constant in $a$ on the non-exceptional sector (measured:
$90.4\%$ of pairs at $c=7$; the leading digits of $\mathcal L'$ are
$a$-free), and equals $d+i_{\min}+e$ with excess $e\ge0$ distributed as
follows: $e\ge1$ requires the sum-factor channel
($C_{i_{\min}}+C'_{i_{\min}}\equiv0\bmod3$, one digit equation, fired
in $\approx61\%$ of sampled pairs since unit-level actives dominate),
and each further excess level is one additional word-digit equation whose
\emph{word-averaged} cost is the averaged trichotomy $\tfrac13$
(the same equidistribution proved operative in the $c=9$ audit). Hence the
ladder weights satisfy the Fubini bounds
\[
\E\bigl[3^{e}\bigr]\le\mathrm{poly}(c),
\qquad
\E\bigl[3^{e/2}\bigr]\le C,
\]
the first marginal (each tail level trades cost $\tfrac13$ against weight
$3$), the second convergent.
\end{lemma}

\begin{lemma}[Hensel profile; simple roots only]\label{lem:hensel}
On the deep region $\{a:\nu(a)\ge d+2\tilde w+1\}$, every point lies in
the basin of a \emph{unique simple} $3$-adic root $\alpha$ of $N$ (since
$v_3(N'(a))=w_2<\infty$ there, clustering is impossible), and the
valuation profile is exactly linear:
\[
\nu(a)=w_2+v_3(a-\alpha).
\]
Hence the strata $\{\nu=d+h\}$, $h>2\tilde w$, are unions of shells
around at most $\deg N\le2c$ roots, of density
$\le2c\cdot3^{\tilde w}\,3^{-h}$---the $3^{-h}$ law measured at $c=9$
(constant $0.3$).
\end{lemma}

\begin{theorem}[Class-IV row budget]\label{thm:classfour}
For every row state $\xi$,
\[
R_{\mathrm{IV}}(\xi)\le\mathrm{poly}(c)\Bigl(\frac{\sqrt3}2\Bigr)^{c}.
\]
\end{theorem}

\begin{proof}[Proof assembly]
Decompose $S(w,w')$ by strata. \emph{(i)} Strata $h$ with
$\nu\le(c+d-1)/2$ vanish exactly (\cref{lem:stratvan}). \emph{(ii)} The
pre-Hensel band $h\le2\tilde w$ is covered by the worst-case envelope
\cref{op:onecond} ($(4/7)^h$, with the $i_{\min}$- and $e$-ladders
costing $3^{-(i_{\min}-2)}$ resp.\ the \cref{lem:ladder} weights: the
band is short and the margin $\sqrt3\cdot\tfrac47<1$ covers it).
\emph{(iii)} Mid-shells in the Hensel range: by
\cref{lem:rootball,lem:hensel} each shell is an exact quadratic Gauss
sum, $|\Sigma_{\mathrm{shell}}|\le\sqrt3\,|{\mathrm{shell}}|^{1/2}$,
dominated by the basin-entry shell: total
$\le\mathrm{poly}(c)\,3^{(c-\tilde w)/2}$ per root, weighted by
$\E[3^{e/2}]\le C$. \emph{(iv)} The fully degenerate cores contribute
their measure, $\le2c\cdot3^{\tilde w}3^{-(c-1-d)}$ per pair, Fubini-paid
by the shell pinning: jointly
\[
\sum_{d}\sum_{i}\underbrace{C_02^{-(d-3)}3^{-(i-2)}}_{\text{shell},\
i_{\min}\text{ cost}}\cdot
\underbrace{3^{\,i+e}3^{-(c-1-d)}}_{\text{core density}}
=\mathrm{poly}(c)\,\E[3^e]\sum_d(3/2)^{d}3^{-c}
\le\mathrm{poly}(c)\,2^{-c},
\]
the $3^{i}$ ball inflation \emph{exactly} paid by the $i_{\min}$
probability (one more marginal balance, absorbed polynomially). Hence
$R_{\mathrm{IV}}(\xi)\le\varphi\mu(x)\cdot\mathrm{poly}(c)2^{-c}
\le\mathrm{poly}(c)(\sqrt3/2)^c$, using
$\mu\le\sqrt{\operatorname{Col}_c(c)}$ (\cref{op:flat}).
\end{proof}

The two steps deferred above are now carried out.

\begin{lemma}[Averaged ladder tail]\label{lem:ladtail}
$\Pr[e\ge r]\le C\,3^{-r}\prod_{r'\le r}(1+C\eta_0^{r'})
\le C'3^{-r}$, hence $\E[3^e]\le\mathrm{poly}(c)$ and
$\E[3^{e/2}]\le C$.
\end{lemma}
\begin{proof}
The level-$r$ excess equation requires one further digit of
$\mathcal L'$ to vanish. Its base aggregates the $r$-th--order digit
contributions of \emph{every} position $i$ with $v_i\le d+r$---a set
whose cardinality grows linearly in $r$ as new positions activate---and
each contributing position carries fresh word randomness with conditional
character average $\le\eta_0<1$ (the support-two freshness of
\cref{op:onecond}(i)). Hence the base equidistributes modulo $3$ with
error $O(\eta_0^{\,\lfloor r/2\rfloor})$, the conditional kill
probability is $\tfrac13(1+C\eta_0^{\,r/2})$, and the product over
levels converges. The measured conditional tail costs
($0.43,\,0.30,\dots$, approaching $\tfrac13$) confirm both the rate and
the sign of the correction.
\end{proof}

\begin{lemma}[Uniformity constants]\label{lem:uniconst}
\emph{(i)} In \cref{lem:stratvan} the constants are absolute: the
stratum conditions are polynomial congruences in $a$ modulo
$3^{\nu_0-d}$, so the stratum is exactly a union of cells of size
$3^{\,c-\nu_0+d}$, and the complete linear subsums vanish with constant
$0$. \emph{(ii)} In \cref{lem:hensel}, Newton iteration from any $a_0$
with $\nu(a_0)\ge d+2\tilde w+1$ converges at the standard quadratic
rate to the unique root $\alpha$ in the basin
$B(a_0,3^{-(\tilde w+1)})$, since $v_3(N'(a))=w_2$ throughout the basin
(derivative rigidity) and $v_3(N''/2)\ge w_3>w_2$ on it
(\cref{lem:rootball}); the profile $\nu(a)=w_2+v_3(a-\alpha)$ is exact
because the linear term of the Taylor expansion at $\alpha$ strictly
dominates all higher terms (the all-derivative grading). All constants
are explicit and uniform in $(w,w',c)$.
\end{lemma}

\noindent With \cref{lem:ladtail,lem:uniconst}, the deferred steps of
\cref{thm:classfour} are discharged: \emph{the class-IV budget, and with
it \cref{thm:final}, holds unconditionally at $s=0$.}

\begin{corollary}[Explicit asymptotic gap]\label{cor:gap}
$\lVert T^{c}_{\chi,c,0}\rVert^{1/c}\le(\sqrt3/2)^{1/2}+o(1)
=0.9306\ldots+o(1)$ uniformly over primitive non-coboundary $\chi$;
hence $\limsup_c\gamma_c(0)\le0.9306\ldots$, and with the strict gap
(\cref{cor:strict}) at the finitely many small $c$,
\[
\sup_c\gamma_c(0)<1 .
\]
An explicit $\delta$ and threshold $c_0$ require evaluating the
polynomial prefactors accumulated above; the asymptotic rate constant
$1-(\sqrt3/2)^{1/2}\approx0.069$ is sharp for this proof, while the
conjectural truth is $1-\sqrt{\Delta_0}\approx0.42$.
\end{corollary}

\begin{remark}[Rigor inventory; the $s$-extension]
The proofs above are complete at the granularity of this manuscript; the
fresh-digit arguments (\cref{op:onecond,lem:ladder,lem:ladtail}) are
fully detailed at their core steps (support-two trichotomy, orbit-triple
spread, source aggregation), with routine expansions left implicit, and
every numerical statement is labelled as verification, not proof. The
extension to $s\in(0,1]$ is \emph{not} a formality: the marginal
coincidence $3\Delta_0=1$ is specific to $s=0$ (for $s>0$,
$3\Delta_s>1$, and the shell normalizations must be reworked against the
$W_s$-normalized operator); the extension is expected by the same scheme
with $\Delta_s$-weighted constants, but is left open here.
\end{remark}

\begin{theorem}[Finite H23a at $s=0$]
\label{thm:final}
By \cref{thm:classfour} with \cref{lem:ladtail,lem:uniconst}, the
class-IV row budget satisfies
$R_{\mathrm{IV}}(\xi)\le\mathrm{poly}(c)(\sqrt3/2)^c$. Hence with $R_{\mathrm
I}=0$, \cref{prop:rowcontract}, and a branch cutoff $B\asymp c$ for the
tail,
\[
\sup_\xi R_c(\xi)\le\mathrm{poly}(c)\Bigl(\tfrac{\sqrt3}2\Bigr)^{c}
\quad\Longrightarrow\quad
\lVert T^c_{\chi,c,0}\rVert^{1/c}\le\frac{3^{1/4}}{\sqrt2}+o(1)
\approx0.9306,
\]
uniformly over primitive non-coboundary $\chi$. For the finitely many
$c<c_0$ the strict gap (\cref{cor:strict}) already gives $\gamma_c(0)<1$;
therefore
\[
\sup_c\gamma_c(0)<1 :
\]
the finite H23a operator theorem holds at $s=0$. (No certified numerics are needed for the qualitative
statement; explicit $\delta$ requires evaluating the finite minimum.)
\end{theorem}

In operator form the cascade reduction reads as follows; with
\cref{op:flat} now proved, it holds \emph{unconditionally} for the
cascade sector by \cref{lem:pinning,prop:rowcontract}:

\begin{theorem}[Threshold cascade contraction, cascade sector]\label{op:cascade}
Let $\mathcal K_c$ be the threshold cascade operator on states recording
the cascade level $i$, the locked partial-sum difference class $t_i$
modulo $2\cdot3^{\,c-3-i}$, and the target residue, with edge weights given
by the branch probabilities times the local coset factor
$\beta\in\{\tfrac12,1\}$ of \cref{prop:twopoint}. Then the cascade
(class I--III) row budgets of $G_c$ satisfy
$\sup_\xi R_{\mathrm{casc}}(\xi)\le\mathrm{poly}(c)(\sqrt3/2)^c$
(\cref{prop:rowcontract} with the now-proved \cref{op:flat}); granted in
addition the class-IV budget (audited subordinate, unproved), this yields
the threshold Schur bound (\cref{conj:schur}), hence
$\lVert T^c_{\chi,c,0}\rVert^{1/c}\le1-\delta'$ and
$\sup_c\gamma_c(0)<1$, with small $c$ checked by finite verification.
\end{theorem}

\paragraph{Status of the contraction program.} Of the three lemmas
originally projected, two are now settled and one is dissolved:
off-representative control is \emph{proved in absorbed form}
(\cref{lem:pinning}: total factor $\le C_0$, uniformly); heavy-state
nonpersistence is \emph{unnecessary} in scaled coordinates (the heavy
state is the Perron state and contracts at $3/4$ per level); and the
$\beta$-profile bounds (\cref{lem:betachannel}, the channel constant
$q\le\frac13$) are demoted to sharpening constants, since $\beta\le1$
already yields contraction in \cref{prop:rowcontract}. What remains for
the class-II/III budget is exactly one analytic input---the threshold
($m=c$) extension of the flattening lemma,
$\mathrm{Col}_c(c)\le\mathrm{poly}(c)3^{-c}$---plus the class-IV budget as
audited, and small $c$ by direct computation.

The remaining finite-side statement is therefore:

\begin{conjecture}[Threshold Schur theorem; $m=c$]\label{conj:schur}
Fix $s=0$ and the unit-mass normalization $T=W_s^{-1}M_{\chi,c,s}$. There
exist $\delta>0$ and $c_0$ such that for all $c\ge c_0$, uniformly over
primitive non-coboundary $\chi$ modulo $3^c$,
\[
\sup_{\xi}\sum_{\xi'}\bigl|G_c(\xi,\xi')\bigr|\le(1-\delta)^{2c},
\qquad\text{hence}\qquad
\rho(T_{\chi,c,0})\le\lVert T_{\chi,c,0}^{\,c}\rVert^{1/c}\le1-\delta.
\]
\emph{Status of the four budgets (updated).} $R_{\mathrm I}=0$ is proved
(exact vanishing); the cascade budget $R_{\mathrm{II/III}}$ is
\emph{proved} contractive unconditionally
(\cref{lem:pinning,prop:rowcontract} with the bounded-shell theorem
\cref{op:flat}), at rate $(\sqrt3/2)^c$; the class-IV budget is closed by the
one-condition degeneracy cost (\cref{op:onecond},
$\lambda=\frac47<3^{-1/2}$) together with the root-ball analysis
(\cref{thm:classfour}); the tail is handled by a branch cutoff. The
$m=2c$ singular-value decay (rate $\approx0.70$) is real but not
accessible to absolute-value bookkeeping; the threshold $m=c$ is the
regime where completeness first appears while endpoint matching still
suppresses degenerate nonstationary pairs. The uniform-$s$ extension is a
subsequent target. This is a fixed-power operator-norm statement---outside
the scope of the trace ceiling (\cref{rem:sixth})---and it implies
$\sup_c\gamma_c(0)<1$ directly; see \cref{thm:final}.
\end{conjecture}

\paragraph{Status of earlier routes.} After
\cref{rem:saturation,rem:sixth,lem:coset}, the following are \emph{superseded
as routes to uniformity} and retained for strictness, diagnostics, or
conditional interest only: short-witness coverage (\cref{conj:mult}; no
construction at fixed level), comparison-graph expansion (\cref{conj:exp};
subsumed by descent and the ceiling), the bounded-level holonomy audit and
canonical witness (proved, but yielding strictness, not uniformity), and all
counting/trace chain bounds. The live objects are \cref{conj:schur} and, as
fallback, the fixed-modulus certificates.

\begin{remark}[Relation to the prior certificate program]
A complementary, finite-$Q$ route to the H23a gap was developed in an earlier
phase of this project (the stabilized-residue-extinction manuscript): explicit
per-$Q$ Wielandt certificates $(\mathcal C_Q,\Gamma_Q,\alpha_Q,\Delta_Q,
\eta_Q,M_Q,L_0(Q))$ at mixed moduli $Q=2^r3^s$, organised into four modules
(cocycle perturbation, stabilisation, Wielandt accessibility, non-coboundary
gap), with the crude rigorous gap $\eta_Q=\alpha_Q(1-\cos\Delta_Q)>0$ computed
at every grid modulus. Its numerical finding that no nontrivial character is a
coboundary at any tested $Q$ is structurally completed by
\cref{thm:class} here; its observation that high-frequency $2$-adic characters
approach (but never reach) the coboundary locus is the mixed-modulus shadow of
the unique parity coboundary $\chi_2$. On the orbit side, the realizability
audits of that phase found actual orbits and confined renewal controls
\emph{word-level indistinguishable}, with the sole gap a canonical-residue
pinning---precisely the structure carried here by the window variable $X_w$
(\cref{lem:walkwindow}) and the lift phases. If the asymptotic route through
\cref{conj:wrap} stalls, the certificate route provides finite verification at
any fixed modulus, at the cost of uniformity.
\end{remark}

\begin{theorem}[Conditional high-conductor damping, corrected assembly]
\label{thm:assembly}
\cref{conj:wrap} implies
$\limsup_{c\to\infty}\gamma_c(s)\le3^{1/(2\alpha k_0)}\sqrt{\Delta_s}<1$ for
$\alpha k_0$ large, with the conjecturally sharp constant $\sqrt{\Delta_s}$ in
the limit $\alpha k_0\to\infty$ (\cref{rem:collision}).
\end{theorem}
\begin{proof}[Proof, modulo \cref{conj:wrap}]
$(T^m)^*T^m$ is positive semidefinite, so
$\rho(T)^{2mk_0}\le\operatorname{tr}\bigl[((T^m)^*T^m)^{k_0}\bigr]$. Insert the
conjectured bound and take $2mk_0$-th roots:
$\dim^{1/mk_0}\le3^{\,c/mk_0}\to3^{1/\alpha k_0}$ at $m=\alpha c$, and
$K^{1/2mk_0},\ \Delta_s^{-k_0/2mk_0}\to1$.
\end{proof}

\begin{remark}[Collision mass and a predicted limit]\label{rem:collision}
By \cref{prop:paircorr} the exactly-stationary pairs are those agreeing in the
first $n-1$ symbols (the last symbol is invisible in $C_w$), of weighted mass
\[
\Bigl(\sum_b(w_b/W)^2\Bigr)^{n-1}
=\Bigl(\frac{1-q}{1+q}\Bigr)^{n-1}
=\Bigl(\frac{2^{1+s}-1}{2^{1+s}+1}\Bigr)^{n-1},
\]
the dichotomy constant of \cref{prop:dich} yet again. The $L^2$ skeleton---%
nonstationary pairs vanish exactly (\cref{lem:vanish}), stationary mass decays
geometrically---predicts for the high-conductor tail
\[
\lim_{c\to\infty}\gamma_c(s)\stackrel{?}{=}
\sqrt{\frac{2^{1+s}-1}{2^{1+s}+1}},
\qquad\text{i.e.}\quad
\gamma_\infty(0)=3^{-1/2}\approx0.5774 ,
\]
and the observed primitive-sector radii at $s=0$ descend toward exactly this
value ($0.7629,\,0.6454,\,0.6154,\,0.6025$ at $c=4,\dots,7$). What separates
this skeleton from a proof of \cref{conj:hcd} is assembly bookkeeping: the
operator-norm expansion constrains endpoints (incomplete $a$-ranges, the
transition-zone values), and the stationary sector carries the
$\chi(2^{\Delta B})$ geometric phases, to be handled by the damping-lemma
factors. This is now a sharply scoped analytic task rather than an open-ended
search.
\end{remark}

Witness elements are
$W=(\prod_{\gamma_1}F)^{L_2}(\prod_{\gamma_2}F)^{-L_1}\in U_{3^c}$ for cycle
pairs $(\gamma_1,\gamma_2)$; a coboundary multiplier has every such invariant
equal to $1$, of every cycle length. Two statements of very different strength
must be kept apart (generation versus \emph{rapid} generation):

\begin{conjecture}[Short-witness coverage]\label{conj:mult}
Define, for the witness family from cycles of length $\le L$,
\[
\delta_c(L)=\min_{\chi\ \mathrm{primitive\ mod\ }3^c}\ \max_{W\in\mathcal W_c(L)}
\bigl|\chi(W^2)-1\bigr| .
\]
\emph{(Strong / constant gap.)} There are fixed $c_0,L_0$ with
$\inf_c\delta_c(L_0)\ge\delta_0>0$.
\emph{(Weak / polylog gap.)} $\delta_c(O(c^A))\ge\delta_0$ for some fixed $A$.
\end{conjecture}

\begin{theorem}[Conditional uniform gap]\label{thm:condunif}
The strong form of \cref{conj:mult} implies $\sup_c\gamma_c(s)<1$ for each
$s\in[0,1]$, hence with \cref{thm:collapse} a uniform $\kappa_Q\ge\kappa>0$.
The weak form yields $1-\gamma_c\gg c^{-2A}$, hence
$\kappa_Q\gg(\log Q)^{-2A}$.
\end{theorem}
\begin{proof}[Proof, modulo \cref{conj:mult}]
By \cref{thm:descent} (iterated), $\gamma_c=\rho(N)/W$ for an operator $N$ at
fixed level $3^{c_0}$ whose entries have the moduli of the principal operator up
to exponentially small tails, and whose phases have cycle holonomies
$\chi(W^2)$. Since descent is exactly isospectral, only spectral radii are
compared---no eigenvector conditioning enters. Under the strong form the
separating witnesses come from cycles of \emph{fixed} length $\le L_0$: there
are finitely many cycle shapes, holonomies along a fixed cycle are continuous,
so the separation $\delta_0$ survives to every limit point of the compact
matrix family. By Wielandt's equality theorem, $\rho=W$ at a limit point would
force exactly coboundary phases, killing those fixed-length invariants---%
contradiction; hence $\sup\rho(N)<W$. Under the weak form, the
stability-extraction bookkeeping (\cref{prop:stab,thm:condpolylog}) applies
with comparison depth $O(c^A)$, giving $1-\gamma_c\gg c^{-2A}$.
\end{proof}

\begin{remark}[Why fixed length matters]
The compactness step is valid only for witnesses of bounded length: an
invariant carried by cycles whose length grows with $c$ (such as the power
witness of \cref{thm:witness}\,(ii), of length $\asymp3^{c-2}$) does not survive
the limit and proves \emph{strictness} of the gap at each level, not
uniformity. This distinction, emphasised in review, is the precise boundary
between what is proved and what remains.
\end{remark}

\subsection*{Holonomy audit}

\emph{Historical status.} This and the following three subsections
(canonical witness, controlled-length coverage, and the conjectures
\cref{conj:mult,conj:exp}) document an earlier stage of the program. They
remain valid as strictness, diagnostic, and conditional results, but they
are \emph{superseded as routes to the uniform gap} by the threshold
cascade reduction (\cref{lem:pinning,prop:rowcontract,op:cascade}); see
the routes-status paragraph and the Conclusion.

The audit runs in exact integer arithmetic on the $F$-elements
(character-free), on the range where the two-descent recursion is valid
($c\le c_0+2$); cycles of length $\le2$ at level $c_0$ already suffice. For each
primitive $\chi_k$ the separation is
$\max_W\lvert\chi_k(W^2)-1\rvert$, minimised over $k$:
\[
\begin{array}{cc|ccc}
c_0 & c & \#\text{witnesses }W^2 & \min_k\max_W\lvert\chi_k(W^2)-1\rvert & \text{witness subgroup}\\\hline
2 & 4 & 26 & 1.9966 & \text{full principal units}\\
3 & 5 & 80 & 1.9996 & \text{full principal units}\\
4 & 6 & 242 & 1.99996 & \text{full principal units}\\
5 & 7 & 728 & 1.999995 & \text{full principal units}
\end{array}
\]
The separation is not merely bounded below---it is near-maximal (antipodal) and
improving with $c$, because the witness elements \emph{generate the full
principal-unit group} of $U_{3^c}$ at every tested level (gcd of their discrete
logarithms is $1$). This full-generation statement can in fact be \emph{proved}
from a single explicit witness.

\subsection*{The canonical witness}

The level-$9$ graph has exactly one self-loop at each vertex
($x\to_bx$ iff $2^b\equiv3+x^{-1}\bmod9$): branch classes
$b\equiv2,3,6,5,4,1\pmod6$ at $x=1,2,4,5,7,8$ respectively.

\begin{theorem}[Canonical witness]\label{thm:witness}
Let $\gamma_1$ be the self-loop at $x=1$ (branch $b\equiv2$) and $\gamma_2$ the
self-loop at $x=7$ (branch $b\equiv4$) in the level-$9$ graph, and let
$W=F(\gamma_1)F(\gamma_2)^{-1}\in U_{3^c}$ be the associated witness. On the
two-descent range (and for every $c$, granted the carry-depth property below),
\[
W\equiv4\pmod9,\qquad W^2\equiv7\pmod9,\qquad v_3\bigl(\log W^2\bigr)=1 .
\]
Consequently:
(i)~$\langle W^2\rangle=1+3\Z/3^c\Z$, the full principal-unit group, for every
$c$ ($1+3u$ with $u\in\Z_3^\times$ topologically generates);
(ii)~the power witness (the same cycle pair traversed $3^{c-2}$ times)
satisfies $W^{2\cdot3^{c-2}}=1+2\cdot3^{c-1}$ exactly, so \emph{every} primitive
$\chi$ maps it to a primitive cube root of unity:
\[
\bigl|\chi\bigl(W^{2\cdot3^{c-2}}\bigr)-1\bigr|=\sqrt3
\qquad\text{uniformly in }c\text{ and }\chi .
\]
\end{theorem}
\begin{proof}
On the two-descent range $F(x,b)=(3x+1)2^{-b}\cdot\rho^{e}$ with
$\rho=1+3^{c-1}\equiv1\pmod9$ (for $c\ge3$), so modulo $9$,
\[
W\equiv\frac{3\cdot1+1}{3\cdot7+1}\cdot2^{\,4-2}
=\frac{4}{22}\cdot4\equiv1\cdot4=4\pmod 9,
\]
using $22\equiv4\pmod9$. Then $W^2\equiv16\equiv7\pmod9$, i.e.\
$W^2=1+3u$ with $u\equiv2\pmod3$, giving $v_3(\log W^2)=1$ and (i). For (ii),
$W^{2\cdot3^{c-2}}=(1+3u)^{3^{c-2}}=1+3^{c-1}u\bmod 3^c$ with
$u\equiv2$, an element of order $3$; a primitive $\chi$ is injective on the
order-$3$ top subgroup, so its value is a primitive cube root. (A coboundary
multiplier kills \emph{every} cycle-pair invariant regardless of cycle length,
so power witnesses are legitimate obstructions.) Verified in exact arithmetic
for $c\le12$, including against the full descended tables with corrections.
\end{proof}

\begin{remark}[What remains: carry depth]
The unproved ingredient for general $c$ is the \emph{carry-depth property}:
every correction factor produced by a descent below depth two lies in the deep
principal units $1+3^{c_0}\Z/3^c$ (hence $\equiv1\bmod9$ for $c_0\ge2$). Each
correction is a ratio of $\chi$-lifts of the same lower-level element, which
makes the property natural; for the first two descents it is proved
(\cref{prop:multrec}, corrections are powers of $\rho=1+3^{c-1}$); the general
case is the parity-twisted recursion of \cref{prop:multrec}. The corrected
chain of consequences, with the qualitative and quantitative branches kept
separate, is:
\[
\text{carry depth}
\Rightarrow
\text{\cref{thm:witness} for all }c
\Rightarrow
\text{strict non-coboundary for every primitive sector}
\Rightarrow
\text{strict finite-level gap;}
\]
\[
\text{short-witness coverage (\cref{conj:mult}, strong)}
\Rightarrow
\inf_Q\kappa_Q>0,
\qquad
\text{polylog coverage (weak)}
\Rightarrow
\kappa_Q\gg(\log Q)^{-2A} .
\]
Generation alone (\cref{thm:witness}\,(i)) does \emph{not} imply coverage: a
fixed finite witness set is a fixed family of $3$-adic units, and for $c\to\infty$
its coverage of all primitive characters is governed by the $3$-adic digit
patterns of finitely many logarithm ratios---a normality-type question. The
realistic provable route to coverage is therefore \emph{rapid generation}:
witness families of slowly growing length (the weak form), where the
comparison-graph/mode-spreading structure can act.
\end{remark}

\subsection*{Controlled-length coverage audit}

On the audited range the coverage statistic $\delta_c(L)$ is computed exactly:
\[
\begin{array}{cc|cc|cc}
c_0 & c & \#\mathcal W_c(1) & \delta_c(1) & \#\mathcal W_c(2) & \delta_c(2)\\\hline
2 & 3 & 7 & 1.9696 & 8 & 1.9696\\
2 & 4 & 10 & 1.4547 & 26 & 1.9966\\
3 & 5 & 75 & 1.9996 & 80 & 1.9996\\
4 & 6 & 241 & 1.99996 & 242 & 1.99996\\
5 & 7 & 728 & 1.999995 & 728 & 1.999995
\end{array}
\]
Even loop-pair witnesses ($L=1$) give near-maximal coverage, with one dip at
$(c_0,c)=(2,4)$ restored by $L=2$. Caveat: on this diagonal $c_0$ grows with
$c$; the decisive asymptotic regime (fixed $c_0$, $c\to\infty$) requires the
parity-twisted recursion (Task~A) to extend the multiplier tables, and is the
next computation after it.

\subsection*{The missing ingredient, stated precisely}

To convert \cref{prop:stab} into $\sup_c\gamma_c<1$ by the comparison route one
must propagate the approximate identities globally and contradict the exact
orthogonality $\sum_{a\in\ker\chi_2}\chi(a)=0$; by \cref{cor:descentcomp} only
the bounded-level core of this propagation is still needed, and
\cref{conj:mult} is its sharpest form (the comparison graph below remains a
useful route to it).

\begin{definition}[Comparison graph]\label{def:compgraph}
Fix $c$ and a branch bound $B$. The vertices are the states
$U_{3^c}\times\Z/2$. There are two edge types.
\emph{Common-source edges}: for each state $x$ and $b,b'\le B$ with
$b\equiv b'\pmod2$, join $S_b(x)$ to $S_{b'}(x)$ (these realise
$v\sim4^{-j}v$ at equal parity; $b'-b$ odd joins
$(v,\sigma)\sim(2^{-t}v,\sigma+t)$ across parity).
\emph{Common-successor edges}: join $x=(a,\tau)$ to $x'=(a',\tau')$ if
$S_b(x)=S_{b'}(x')$ for some $b,b'\le B$ (equivalently
$3a'+1\equiv(3a+1)2^{\,b-b'}$, which forces $b\equiv b'\bmod2$ and
$\tau'=\tau+(b-b')$). A path has \emph{robust diameter} $\le D$ at mass
threshold $\theta$ if any two vertices in the same $\ker\chi_2$-fiber are
joined by a path of length $\le D$ all of whose edges use branches of weight
$w_b\ge\theta$, and this remains true after deleting any edge set of
stationary mass $<\theta$.
\end{definition}

\begin{conjecture}[Comparison-graph expansion; strong and weak forms]\label{conj:exp}
\emph{(Strong.)} There are $B,A$ such that for every $c$ the comparison graph
has robust diameter $O(c^A)$ at mass threshold $\gg c^{-A}$.
\emph{(Weak, sufficient for \cref{thm:condpolylog}.)} For every non-coboundary
$\chi$, from every high-mass state the approximate-coboundary phase propagates
within $O(c^A)$ comparisons to representatives of a subgroup $H\le U_{3^c}$
with $\chi|_H\ne1$ and
$\bigl|\,|H|^{-1}\sum_{h\in H}\chi(h)\bigr|\le1-\eta$, $\eta\gg c^{-A}$.
By \cref{cor:descentcomp}, both forms need only be established at a bounded
level $3^{c_0}$.
\end{conjecture}

\begin{theorem}[Conditional polylog gap]\label{thm:condpolylog}
Assume \cref{conj:exp}. Then $1-\gamma_c(s)\gg c^{-2A}$ uniformly for
$s\in[0,1]$; by \cref{thm:collapse}, $\kappa_{3^r}(s)\gg(\log Q)^{-2A}$, the
quantitative form of the finite spectral hypothesis.
\end{theorem}
\begin{proof}[Proof, modulo \cref{conj:exp}]
If $\gamma_c\ge1-\varepsilon$, \cref{prop:stab} provides $d$ satisfying the
approximate coboundary identity off a defect set of mass $O(\varepsilon)$. By
\cref{conj:exp} there is a comparison tree of depth $O(c^A)$ avoiding the defect
set, along which the relation $\chi(a)d(x)\approx\chi(a')d(x')$ propagates with
total accumulated error $O(D\varepsilon^{1/2})$, $D=O(c^A)$ the path length
(Cauchy--Schwarz over the path). The $b,b{+}2$ comparisons make $d$
approximately $\langle4\rangle$-invariant on the tree, and Step~2 of
\cref{thm:class} then forces $\chi$ to be $O(D\varepsilon^{1/2})$-close to a
constant on $\ker\chi_2$. Averaging,
$\bigl|\sum_{a\in\ker\chi_2}\chi(a)\bigr|\ge\bigl(1-O(D\varepsilon^{1/2})\bigr)
\lvert\ker\chi_2\rvert$, contradicting exact orthogonality (a non-coboundary
$\chi$ is nontrivial on the squares) unless $D\varepsilon^{1/2}\gg1$, i.e.\
$\varepsilon\gg D^{-2}=c^{-2A}$---the origin of the exponent $2A$.
\end{proof}

\begin{remark}[Why this and not monotonicity or single-orbit propagation]
Monotonicity in $c$ is false (\cref{sec:numerics}), so no monotone conductor
argument exists. Propagating $\langle4\rangle$-invariance along the full
$\langle4\rangle$-orbit (length $3^{c-1}$) costs a factor $3^c$ and yields only a
polynomial-in-$Q$ gap; \cref{lem:damp} shows the same $9^{-c}$ scale appears as
the one-step damping of the least-damped primitive mode. The polylog/constant
gap therefore needs short comparison paths (the expansion conjecture) or,
equivalently, the spreading of $\eta\mapsto\eta\circ A$ across the character
group. The numerical decay of $\gamma_c$ beyond a fixed conductor is consistent
with both mechanisms acting.
\end{remark}

\paragraph{Block-cancellation alternative.} The earlier Doeblin route remains
viable and is complementary: exhibit from every state a block family of
untwisted mass $\ge\alpha$ whose endpoints contain a full coset of a subgroup
$H$ with $\chi|_H\ne1$ (exact cancellation $\sum_{u\in H}\chi(uh)=0$). For the
killed strip, build the families with identical valuation multisets so the
height weights match exactly, or prove bounded distortion
$w_s(w_1)/w_s(w_2)\in[C^{-1},C]$ with QSD bulk retention $O(e^{-cL})$.

\section{The keystone: bad orbits cast finite non-coboundary shadows}
\label{sec:keystone}

Throughout this paper the global implication has been kept strictly
separate from the finite-operator theory. We now prove its first half:
the \emph{keystone arrow}. Call a Syracuse orbit \emph{bad} if it is
divergent or eventually enters a nontrivial cycle.

\begin{lemma}[Drift]\label{lem:drift}
For every bad orbit, $\limsup_N B_N/N\le\log_23+\log_2\tfrac76<1.81<2$.
\end{lemma}

\begin{proof}
A non-cyclic orbit visits pairwise distinct integers, so
$\sum_{n<N}1/x_n=O(\log N)$, and
$\log_2x_N=\log_2x_0+N\log_23-B_N+\sum_{n<N}\log_2(1+\tfrac1{3x_n})$
with $x_N\ge1$ gives $B_N/N\le\log_23+O(\tfrac{\log N}N)$. On a
nontrivial cycle of period $p$, $2^{B_p}=3^p\prod(1+\tfrac1{3x_i})$ and
$x_i\ge2$ give $B_p/p\le\log_23+\log_2\tfrac76$. (The stationary mean is
$\E_{\mathrm{geom}}[b]=2$, attained with equality only by the trivial
cycle, where $b\equiv2$.)
\end{proof}

\begin{lemma}[Letter characters are two-point shadows]\label{lem:lettershadow}
For any character $\chi'$ modulo $3^r$, the finite-level two-point
function
\[
F_{\chi'}(x_n,x_{n+1}):=\chi'(x_{n+1})\,\overline{\chi'}(3x_n+1)
=\chi'(2)^{-b_n}
\]
is exactly a character of the letter $b_n$, of order equal to
$\operatorname{ord}\chi'$ (since $2$ generates the units). By the
classification $\operatorname{Cob}=\langle\chi_2\rangle$
(\cref{thm:class}), the letter characters realized by
\emph{non-coboundary} $\chi'$ are exactly those with nontrivial
$3$-part; the exempt directions are the trivial character and the parity
$(-1)^{b_n}$ (realized by $\chi_2$, the valuation-parity identity).
\end{lemma}

\begin{theorem}[Keystone arrow]\label{thm:keystone}
Every bad orbit casts a finite non-coboundary spectral shadow: there
exist $r$ and a non-coboundary character $\chi'$ modulo $3^r$ such that
the empirical averages
$\frac1N\sum_{n<N}\chi'(x_{n+1})\overline{\chi'}(3x_n+1)$ do \emph{not}
converge to the stationary value
$\E_{\mathrm{geom}}\bigl[\chi'(2)^{-b}\bigr]$.
\end{theorem}

\begin{proof}
Suppose not: every such average converges to its stationary value. By
\cref{lem:lettershadow} this says the empirical letter averages
$\frac1N\sum\xi(b_n)$ converge to $\E_{\mathrm{geom}}[\xi(b)]$ for every
letter character $\xi$ of modulus $2\cdot3^j$, $j\ge1$, \emph{except
possibly} $\xi=$ parity. Pass to a subsequence realizing
$\beta:=\liminf B_N/N\le1.81$ (\cref{lem:drift}) and let $\nu$ be the
corresponding empirical limit law of the letter $b$. Truncation gives
$\E_\nu[b\wedge K]\le\beta$ for all $K$, hence tightness. The law of $b$
modulo $2\cdot3^j$ under $\nu$ matches the geometric law in every
character except parity, so
\[
\Pr_\nu[b\equiv k\ (2\cdot3^j)]
=\Pr_{\mathrm{geom}}[b\equiv k\ (2\cdot3^j)]+\frac{t\,(-1)^k}{2\cdot3^{j}}
\quad\text{for some }|t|\le2 :
\]
the single deviating direction \emph{spreads} over the residues, with
pointwise correction $O(3^{-j})$. Letting $j\to\infty$ and using
tightness, every point mass matches exactly:
$\Pr_\nu[b=k]=2^{-k}$ for all $k$. These sum to $1$ (no escaping mass),
so $\E_\nu[b]=2$---contradicting $\E_\nu[b]\le\beta<2$.
\end{proof}

\begin{proposition}[Drift visibility at the tilted trivial character]
\label{prop:tiltshadow}
For $s\in(0,1)$, the trivial-character defect average of a bad orbit
satisfies, by Jensen and \cref{lem:drift},
\[
\frac1N\sum_{n<N}2^{\,s(\log_23-b_n)}\ \ge\ 2^{\,s(\log_23-B_N/N)}
\ \xrightarrow[\ N\ ]{}\ \ge\ 2^{\,s(\log_23-1.81)}>1
>\frac{3^s}{2^{1+s}-1}=\E_{\mathrm{geom}}\bigl[2^{s(\log_23-b)}\bigr]:
\]
the $s$-tilt alone already separates bad orbits from stationarity, with
equality of the two endpoints exactly at $s\in\{0,1\}$ (the Collatz
criticality $3=2^{1+1}-1$).
\end{proposition}

\begin{remark}[The arrow is one-directional]\label{rem:onedir}
The converse of \cref{thm:keystone} fails, and this matters for the
correct formulation of what remains. The \emph{trivial cycle} itself
($x=1$, $b\equiv2$) has empirical shadow value $\chi'(2)^{-2}$, which
differs from $\E_{\mathrm{geom}}[\chi'(2)^{-b}]$ for general
non-coboundary $\chi'$; hence \emph{every convergent orbit} also casts a
non-stationary shadow (its averages converge to the trivial-cycle value,
not the geometric one). A naive ``no orbit deviates'' rigidity statement
is therefore \emph{false} if the Collatz conjecture is true. The correct
remaining statement is an exclusion:
\end{remark}

\begin{openproblem}[Quenched shadow exclusion]\label{op:quenched}
No positive Syracuse orbit that is \emph{not eventually trivial}
supports a persistent non-coboundary two-point shadow deviation from the
geometric stationary value. Equivalently: for every orbit not absorbed
by $1$, every finite non-coboundary shadow either converges to the
stationary value or forces descent within bounded time. Combined with
\cref{thm:keystone}, this excludes bad orbits and implies the full
Collatz conjecture. Two cautions: \emph{(i)} if Collatz is true the
statement is vacuous on positive integers, so it cannot be tested on
known convergent orbits; \emph{(ii)} it is essentially the Collatz
conjecture in minimal finite-shadow language---the value of the
reduction is the finiteness and explicitness of the obstructing object,
not a decrease in logical strength.
\end{openproblem}

\paragraph{Possible routes.} Three attack shapes, all open and all at
the deterministic-equidistribution barrier. \emph{(1) Finite inverse
theorem for shadow bias:} persistent deviation $\ge\varepsilon$ forces,
via a density-increment argument and the sliding-window congruence, a
recurrent structured residue trap modulo $3^R$; prove any such recurrent
trap is an exact algebraic obstruction, hence coboundary by
\cref{thm:class}, hence impossible for non-coboundary $\chi'$. \emph{(2)
Shadow deviation forces descent:} for every $(r,\chi',\varepsilon)$
there are $(M,H)$ such that a deviation $\ge\varepsilon$ over a
length-$M$ window forces $x_{n+h}<x_n$ within a bounded
enlargement---local and dynamical, avoiding global normality: the most
realistic target. \emph{(3) Empirical-measure rigidity:} every
non-eventually-trivial positive-orbit limit measure has annealed shadow
marginals---a measure classification of exactly the type that is the
known wall (symbolic models admit many invariant measures; the
integer-orbit constraint is the arithmetic heart).

\subsection*{Three finite approaches to quenched exclusion}

\begin{theorem}[Periodic words cast shadows]\label{thm:periodic}
Every positive periodic Syracuse orbit other than the trivial cycle
casts a finite non-coboundary shadow: some non-coboundary $\chi'$ has
cycle average $\frac1p\sum_i\chi'(2)^{-b_i}\ne
\E_{\mathrm{geom}}[\chi'(2)^{-b}]$.
\end{theorem}

\begin{proof}
The cycle's letter law is a finitely supported probability measure. If
it matched the geometric law in every letter character except parity,
the spreading argument of \cref{thm:keystone} (the parity correction is
$O(3^{-j})$ pointwise at modulus $2\cdot3^j$) would force
$\Pr[b=k]=2^{-k}$ for \emph{every} $k$---impossible for finite support.
The detecting character has order with nontrivial $3$-part, hence is
non-coboundary by \cref{thm:class}. (The trivial cycle also deviates,
$b\equiv2$; it is the designated exception throughout.)
\end{proof}

\begin{corollary}
Cycle elimination is a special case of quenched shadow exclusion
(\cref{op:quenched}). Computationally: all integer cycles with period
$p\le7$ are rotations of the trivial cycle (exhaustive search over
$2^B>3^p$ with $(2^B-3^p)\mid A_w$), and every sampled nontrivial
periodic word is detected already at character order $3$ or $9$.
\end{corollary}

\begin{proposition}[Finite inverse theorem: elementary layers]
\label{prop:inverse}
Let $z=\chi'(2)$, $m=\operatorname{ord}z$, and
$D_N(\chi')=\frac1N\sum_{n<N}z^{-b_n}-\E_{\mathrm{geom}}[z^{-b}]$.
If $|D_N|\ge\varepsilon$ then: \emph{(i)} some residue class $a\bmod m$
of the letters is biased,
$\bigl|\frac1N\#\{n:b_n\equiv a\}-\Pr_{\mathrm{geom}}[b\equiv a]\bigr|
\ge\varepsilon/m$ (Fourier inversion on $\Z/m$); \emph{(ii)}
equivalently, by $F_{\chi'}=\chi'(x_{n+1})\overline{\chi'}(3x_n+1)$, the
empirical \emph{edge} measure on $U_{3^r}\times U_{3^r}$ deviates from
the stationary edge law by $\gg_r\varepsilon$ on an explicit
residue-edge set: \emph{shadow bias is a finite residue-edge trap}. The
third layer---\emph{persistent trap $\Rightarrow$ descent or coboundary
obstruction}, with the latter impossible for non-coboundary $\chi'$ by
\cref{thm:class}---is the open target; we emphasize, per the warning
above, that it would constitute a genuinely new deterministic rigidity
principle, not a corollary of the finite gap.
\end{proposition}

\begin{proposition}[C1: mean-visible deviation forces descent]
\label{prop:meandescent}
If a window $[N,N+M)$ of a positive orbit has mean letter
$\frac1M\sum b_n\ge\log_23+\eta$ with $\eta>\bigl(3x_N\ln2\bigr)^{-1}$,
then $\min_{N<n\le N+M}x_n<x_N$: descent within the window.
\end{proposition}

\begin{proof}
If $x_n\ge x_N$ throughout, then
$\log_2x_{N+M}=\log_2x_N+M\log_23-\sum b_n+\sum\log_2(1+\tfrac1{3x_n})
\le\log_2x_N-\eta M+\tfrac M{3x_N\ln2}<\log_2x_N$, a contradiction.
\end{proof}

\paragraph{The C1/C2 split, empirically.} A shadow deviation therefore
splits: \emph{mean-visible} deviations (letter mean displaced) force
descent at once by \cref{prop:meandescent}; \emph{phase-only}
deviations (mean near critical, non-coboundary phase biased) are the
true remaining obstruction. A window audit over long convergent orbits
($648$ windows of length $64$, starts $\sim10^{80}$) shows the split is
real and balanced: the fraction of windows with displaced mean rises
from $4\%$ at low deviation to $53\%$ at high deviation, and phase-only
high-deviation windows exist with letter mean throughout
$[1.78,2.25]$. (On convergent orbits every window descends, so the
positive direction of ``deviation forces descent'' is untestable
there---exactly the vacuousness caution of \cref{op:quenched}; the
informative content is the split structure.) The sharpest remaining
route is thus:
\[
\text{phase-only shadow bias}
\;\Rightarrow\;
\text{residue-edge trap}
\;\Rightarrow\;
\text{coherent lift attempt}
\;\Rightarrow\;
\text{coboundary (excluded) or descent}.
\]

\subsection*{The last arrow: max-plus certificates and where C2 actually lives}

Fix $r$, a non-coboundary $\chi'$, $z=\chi'(2)$, a phase $\omega$, and
$\varepsilon>0$. On the residue-edge graph $G_r$ (vertices $U_{3^r}$,
edges $a\xrightarrow{b}(3a+1)2^{-b}$) put the two labels
\[
g(b)=\log_23-b,
\qquad
\phi(b)=\Re\bigl(\omega\,(z^{-b}-\E_{\mathrm{geom}}[z^{-b}])\bigr).
\]

\begin{lemma}[Certificate forces descent]\label{lem:cert}
Let $H\subseteq G_r$ contain all edges traversed by an orbit window
$[N,N+M)$ with phase bias $\frac1M\sum\phi(b_n)\ge\varepsilon$. Suppose
there exist $\Psi:U_{3^r}\to\R$, $\lambda\ge0$, $\delta>0$ with
\[
g(b)+\lambda\bigl(\phi(b)-\varepsilon\bigr)+\Psi(a')-\Psi(a)\le-\delta
\qquad\text{on every edge of }H .
\]
Then for $M>\operatorname{osc}(\Psi)/\bigl(\delta-(3x_N\ln2)^{-1}\bigr)$
the orbit descends within the window. Dually (max-plus), such a
certificate exists iff
$\inf_{\lambda\ge0}\bigl(P_H(\lambda)-\lambda\varepsilon\bigr)<0$, where
$P_H(\lambda)$ is the maximum directed-cycle mean of $g+\lambda\phi$ on
$H$.
\end{lemma}

\begin{proof}
Telescoping the certificate along the window and using
$\sum(\phi-\varepsilon)\ge0$: $\sum g(b_n)\le-\delta M+
\operatorname{osc}\Psi$. With
$\log_2x_{N+M}=\log_2x_N+\sum g(b_n)+\sum\log_2(1+\frac1{3x_n})$ and the
no-descent assumption $x_n\ge x_N$, the right side drops below
$\log_2x_N$---contradiction. The duality is the standard max-plus
(Karp) characterization.
\end{proof}

\begin{proposition}[The $3$-adic-only formulation of C2 is false]
\label{prop:notthreeadic}
Certificate existence cannot hold at any fixed $3$-adic level, nor in
the $3$-adic tower alone:
\emph{(i)} the all-ones trap ($b\equiv1$; $g\equiv\log_23-1>0$, constant
phase) is coherent modulo $3^R$ for every $R$, biased, descent-free, and
admits no certificate at any level
($P_H(\lambda)-\lambda\varepsilon\equiv\log_23-1$); it is eliminated
\emph{only} by the $2$-adic tower: $b_n\equiv1$ forever forces
$x\equiv-1\bmod2^k$ for all $k$, i.e.\ $x=-1$ (the famous $-1$ cycle).
\emph{(ii)} the critical $\{1,2\}$-mixture trap (letter mean $\log_23$,
$\Pr[b{=}1]=2-\log_23$) is $2$-adically frequency-feasible yet
uncertified at $r=2$ (margin $+0.056$). \emph{(iii)} of the observed
phase-only high-deviation windows on convergent orbits, $7$ of $8$ are
uncertified at $r=2$: observed traps contain spurious cycles, so any
true statement must be projective, not fixed-level.
\end{proposition}

\begin{openproblem}[C2, final form: joint-tower max-plus rigidity]
\label{op:ctwo}
Every $(2,3)$-adically coherent, non-coboundary--biased, descent-free
trap---a projective family of residue-edge traps over the \emph{mixed}
moduli $2^a3^b$ with nonvanishing mass, bias $\ge\varepsilon$, and
near-critical drift---admits a max-plus descent certificate at some
level of the joint tower. By \cref{lem:cert} this implies the quenched
shadow exclusion (\cref{op:quenched}), hence the Collatz conjecture. The
$2$-adic constraints are not optional: by \cref{prop:notthreeadic} they
supply the negative weight that the $3$-adic graph alone cannot see.
This is where the present program rejoins its predecessor: the
mixed-modulus $(2^r3^s)$ certificates of the stabilized-residue
program are exactly finite verifications in this joint tower.
\end{openproblem}

\noindent The final route, corrected once more:
\[
\text{phase-only bias}
\Rightarrow
\text{joint-tower trap}
\Rightarrow
\begin{cases}
\text{max-plus certificate}\Rightarrow\text{descent (\cref{lem:cert}),}\\[2pt]
\text{no certificate at any level}\Rightarrow\text{a coherent
$(2,3)$-adic object,}
\end{cases}
\]
and the remaining theorem is that the only such objects are the
classified coboundary/trivial ones. C2 in full remains open; what is
proved here is the mechanism (\cref{lem:cert}), the failure of every
purely $3$-adic formulation (\cref{prop:notthreeadic}), and the exact
location of the remaining content.

\begin{remark}[Counterexample profile]\label{rem:profile}
A counterexample to the Collatz conjecture, if it exists, must by the
present theorem stack be \emph{mean-critical}
($\frac1N\sum b_n=\log_23+o(1)$, else \cref{prop:meandescent}),
\emph{non-coboundary phase-biased} (\cref{thm:keystone}), jointly
$(2,3)$-adically coherent, and certificate-free at every level of the
mixed tower. The simplest symbolic models with these properties are
balanced low-letter words over $\{1,2\}$ with
$\Pr[b{=}1]=2-\log_23\approx0.415$ (Sturmian/Beatty type, slope
$2-\log_23$)---the critical analogues of the all-ones word, whose
$2$-adic point is $x=-1$. Direct realizability tests support the ghost
expectation: for twelve intercepts $\theta$ and prefixes of length
$800$ ($B_N=1268$), the unique residue class modulo $2^{B_N+1}$
realizing the Sturmian prefix has least positive representative of full
bit-length $\approx B_N$, wandering with the prefix and never
stabilizing---the signature of a mixed-adic ghost, not a positive
integer. For \emph{periodic} balanced traps the ghost exclusion is
precisely the classical cycle problem, where the archimedean negative
weight is supplied by linear forms in logarithms
($|2^B-3^p|\gg2^B/B^{C}$): joint-tower rigidity restricted to periodic
traps is thus partially within reach of existing transcendence
technology, and the genuinely new content of \cref{op:ctwo} is its
aperiodic, single-orbit analogue.
\end{remark}

\subsection*{Calibrated measures, and the periodic case via linear forms
in logarithms}

\begin{lemma}[Calibrated measures]\label{lem:calibrated}
Suppose a coherent mixed-tower trap with non-coboundary bias
$\ge\varepsilon$ admits no max-plus certificate at any level. Then there
exists a shift-invariant probability measure $\mu$ on the valuation
sequence space, with adelically coherent residue marginals, such that
\[
\E_\mu[b]\le\log_23,
\qquad
\textstyle\int\phi\,d\mu\ge\varepsilon .
\]
Consequently $\Pr_\mu[b{=}1]\ge2-\log_23\approx0.415$, and if moreover
$\E_\mu[b]=\log_23$ and $\mu$ is ergodic, the discrepancy cocycle
$S_j=j\log_23-B_j$ is recurrent $\mu$-a.e.\ (Atkinson). Thus
\cref{op:ctwo} is equivalent to: \emph{no such calibrated measure is
realized by a positive integer orbit}.
\end{lemma}

\begin{proof}
At each level $L=2^a3^b$, no certificate means
$\inf_{\lambda\ge0}\bigl(P_{H_L}(\lambda)-\lambda\varepsilon\bigr)\ge0$.
The cycle-mean vectors $(\overline g,\overline\phi)$ of the trap span a
compact convex polytope $C_L$, and by Sion's minimax,
\[
\inf_{\lambda\ge0}\,\max_{\mu\in C_L}
\bigl[\langle\mu,g\rangle+\lambda(\langle\mu,\phi\rangle-\varepsilon)\bigr]
=\max\bigl\{\langle\mu,g\rangle:\ \mu\in C_L,\
\langle\mu,\phi\rangle\ge\varepsilon\bigr\},
\]
so no certificate yields an invariant edge measure $\mu_L$ with
$\int g\,d\mu_L\ge0$ and $\int\phi\,d\mu_L\ge\varepsilon$. The condition
$\int g\ge0$ gives $\E[b]\le\log_23$ uniformly, hence tightness of the
letter marginals; a diagonal weak-$*$ limit along the tower produces a
consistent projective family, i.e.\ a shift-invariant measure on the
inverse limit with the same inequalities. The letter bound follows from
$\E[b]\ge2-\Pr[b{=}1]$; recurrence at exact calibration is Atkinson's
theorem for zero-mean ergodic cocycles.
\end{proof}

\begin{theorem}[Periodic exclusion via linear forms in logarithms]
\label{thm:bakercycle}
There are effective constants $C,\tau$ (one may take $\tau=13.3$, by
Rhin's bound $|B\ln2-p\ln3|\ge c\,B^{-\tau}$) such that every
nontrivial positive integer cycle of period $p$ satisfies
\[
x_{\min}\le C\,p^{\,\tau+1}.
\]
Hence, given the verified range $x_{\min}>2^{71}$, every nontrivial
cycle has $p>(2^{71}/C)^{1/(\tau+1)}$; and for every fixed period range
the balanced (calibrated) periodic traps are effectively excluded: the
linear form in logarithms \emph{is} the max-plus descent certificate for
periodic traps.
\end{theorem}

\begin{proof}
For a cycle, $2^B=3^p\prod_i(1+\tfrac1{3x_i})$, so
$\theta:=B-p\log_23=\log_2\prod_i(1+\tfrac1{3x_i})\in
\bigl(0,\,p/(3x_{\min}\ln2)\bigr]$; positivity and integrality pin
$B=\lceil p\log_23\rceil$, so $B\le2p$. Rhin's effective bound gives
$\theta\ge c'B^{-\tau}\ge c''p^{-\tau}$. Combining,
$c''p^{-\tau}\le p/(3x_{\min}\ln2)$, i.e.\
$x_{\min}\le C p^{\tau+1}$.
\end{proof}

\begin{corollary}[Where a counterexample must live]\label{cor:ghostclass}
Combining \cref{lem:calibrated,thm:bakercycle}: if \cref{op:ctwo} fails,
the witnessing calibrated measure is \emph{aperiodic}, supported on
low-letter sequences with $\Pr[b{=}1]\ge0.415$, with recurrent
discrepancy and non-coboundary bias---the Sturmian ghost class of
\cref{rem:profile}---and all of its periodic approximants of any fixed
period range are effectively excluded by \cref{thm:bakercycle}. The
remaining content of the program is thus the aperiodic, single-orbit
analogue of the linear-forms certificate: \emph{Baker along orbits}.
\end{corollary}

\subsection*{Repetition rigidity: exclusion of low-complexity calibrated
words}

The aperiodic analogue can now be carried out, and it closes more
strongly than anticipated: the contraction bookkeeping of the
repeated-block heuristic is unnecessary, because $2$-adic realizability
gives \emph{exact} rigidity. Recall the classical coding: for a
valuation block $u=(b_0,\dots,b_{p-1})$ with sum $B_u$, the set of odd
$x$ whose first $p$ letters equal $u$ is a \emph{single residue class
modulo $2^{B_u+1}$} (machine-checked: $84{,}471$ instances, no
exceptions), and $B_u\ge p$.

\begin{theorem}[Repetition rigidity]\label{thm:repetition}
Let $(x_n)_{N\le n\le N+M}$ be a positive Syracuse orbit segment with
$x_n\in[Y,WY]$ throughout, and let $p$ satisfy $2^{p+1}>Y(W-1)$. If the
segment's valuation word has fewer than $M-p+1$ distinct factors of
length $p$, then $x_n=x_m$ for some $N\le n<m\le N+M-p$: the orbit is an
integer cycle of period $\le M$ and minimum $\ge Y$.
\end{theorem}

\begin{proof}
By pigeonhole two positions $n<m$ carry the same length-$p$ block $u$;
by the coding, $2^{B_u+1}\mid x_m-x_n$ with $2^{B_u+1}\ge2^{p+1}>
Y(W-1)\ge|x_m-x_n|$; hence $x_m=x_n$, and the orbit repeats.
\end{proof}

\begin{corollary}[Bounded-discrepancy aperiodic exclusion]
\label{cor:sturmianexcl}
No aperiodic valuation word with bounded discrepancy
$\sup_N|N\log_23-B_N|\le D$ and subexponential factor complexity
($P(p)\le2^{\,p-2D-c_0}$ for large $p$) is realizable by a positive
Syracuse orbit. In particular \emph{no Sturmian word of any slope}
($D<1$, $P(p)=p+1$) is realizable: the critical ghost class of
\cref{rem:profile} is excluded unconditionally---no linear-forms input
is needed for the qualitative statement, since the manufactured cycle's
word is periodic, contradicting aperiodicity.
\end{corollary}

\begin{proof}
A realization from $x_0$ is confined to the band
$[x_02^{-D-1},x_02^{D+1}]$ for $M\le x_0$ steps (the
$\log_2(1+\tfrac1{3x})$ corrections total $o(1)$). Take
$p=\lceil\log_2x_0\rceil+2D+3$ and $M=p+P(p)+1\le x_0$ (subexponential
complexity; small $x_0$ is covered by the verified convergence range).
\cref{thm:repetition} produces a cycle, whose word is eventually
periodic---contradiction.
\end{proof}

\begin{theorem}[Complexity--height tradeoff]\label{thm:tradeoff}
With the effective constants of \cref{thm:bakercycle}: every positive
orbit segment of length $M$ confined to a dyadic band $[Y,2Y]$ with
$Y>CM^{\tau+1}$ has word complexity
\[
P(p)\;\ge\;M-p+1,
\qquad p=\lceil\log_2Y\rceil+2 :
\]
a no-descent stretch at height $Y$ must exhibit \emph{nearly full factor
diversity} at scale $\log_2Y$. Equivalently, any banded no-descent
window longer than $(Y/C)^{1/(\tau+1)}$ forces the word to mimic
randomness at that scale.
\end{theorem}

\begin{proof}
Otherwise \cref{thm:repetition} manufactures a cycle with period $\le M$
and minimum $\ge Y$, and \cref{thm:bakercycle} gives
$Y\le CM^{\tau+1}$---contradiction.
\end{proof}

\begin{corollary}[The ghost class is closed; the frontier moves]
\label{cor:ghostclosed}
Combining \cref{lem:calibrated,thm:tradeoff}: a counterexample orbit
whose discrepancy path admits banded windows of unbounded length at
heights $Y_k$ (the behaviour forced at the measure level by Atkinson
recurrence) must carry a word of factor complexity at least
$(Y_k/C)^{1/(\tau+1)}-p_k$ at scales $p_k\approx\log_2Y_k$:
\emph{exponential rate $\ge1/(\tau+1)$ per bit}. All low-complexity
ghosts---Sturmian, morphic, linearly recurrent, polynomial, and
everything below that exponential rate---are excluded. The Aperiodic
Baker Rigidity target of the route map is thereby proved in a
strengthened form, with the repeated block replacing the closed cycle
exactly as envisaged, but with exact $2$-adic equality in place of the
contraction squeeze. What survives is precisely: orbits whose valuation
words have near-full exponential complexity inside every banded window
(randomness-mimicking), or whose discrepancy path forever avoids long
banded windows; for these the $S$-unit/subspace frontier is the next
tool. We are not aware of prior statements of
\cref{thm:repetition,thm:tradeoff,cor:sturmianexcl}, though their
ingredients ($2$-adic coding; cycle bounds via linear forms) are
classical.
\end{corollary}

\begin{remark}[The wall, in its sharpest form]
\cref{op:quenched} is the quenched-equidistribution wall of the earlier
program, now reduced to its minimal finite form: single-orbit Birkhoff
averages of finite-level functions. The finite operator theorem of this
paper ($\sup_c\gamma_c(0)<1$) is its exact \emph{annealed} counterpart:
the transfer operator mixes these very characters at a uniform
exponential rate. What is missing is the passage from annealed mixing to
pointwise convergence along the deterministic forward orbit---the step
that no current technique supplies for a non-invertible arithmetic map
with no invariant measure on the integers. The contribution of the
keystone arrow is to make the two sides meet at the same finite object:
one two-point character family, with the coboundary direction
classified, exempted, and shown harmless.
\end{remark}

\section*{Conclusion}
\addcontentsline{toc}{section}{Conclusion}

The defect cocycle relocates rather than removes the difficulty, but the
finite-operator side is now substantially a theorem. The rigorous core comprises
the sliding-window congruence; the proof that the quadratic character is
valuation parity in disguise, hence an exact coboundary; the \emph{classification}
$\operatorname{Cob}=\langle\chi_2\rangle$ (\cref{thm:class}), which makes the
non-coboundary quotient canonical rather than conjectural; the strict spectral
gap at every finite level (\cref{cor:strict}); the conductor collapse
(\cref{thm:collapse}), which reduces uniformity in $Q$ to the single computable
sequence $\gamma_c(s)$; the height-strip rotation identity
(\cref{prop:rot}); the conductor-raising ejection theorem
(\cref{thm:eject}); and the cascade equivalence, mass bound, two-point
profile, and deeper-channel bound
(\cref{thm:cascade,lem:cascade,prop:twopoint,lem:betachannel}). The
finite-operator endpoint narrowed in stages, each closed by an honest
negative finding: descent (\cref{thm:descent,thm:rank2}) makes the three
deepest digits lossless but \emph{stalls there} (\cref{rem:saturation});
every counting and trace assembly terminates at exact marginality
(\cref{rem:sixth}); and the coarse Schur regimes fail precisely where the
degenerate class decouples from endpoint matching. What survives is the
threshold $m=c$: the complete-kernel formula (\cref{prop:completekernel})
and the adjoint-side Schur form reduce the uniform gap \cref{conj:ext} to
the threshold Schur bound (\cref{conj:schur}). Within it, the cascade
(class-II/III) budget is now \emph{proved contractive}
(\cref{lem:pinning,prop:rowcontract}: sharp pinning constant
$C_0<1.355$, heavy-row rate $(3/4)^c$, all-row rate $(\sqrt3/2)^c$) modulo
a single analytic input---threshold flattening, reduced in
\cref{op:flat} to an exact, character-free renewal Fourier-decay
statement---with the class-IV budget subordinate by audit. The flattening
input is \emph{proved} (\cref{op:flat}: bounded shells,
$\tilde Q_L\le3C_1<5.16$, by the digit-extension energy argument), so the
cascade sector of the threshold Schur bound is an unconditional theorem
(\cref{op:cascade}); and the full finite H23a operator theorem
$\sup_c\gamma_c(0)<1$ is now \emph{proved} (\cref{thm:final,cor:gap}),
with the class-IV budget closed by the root-ball analysis
(\cref{lem:ladder,lem:hensel,thm:classfour,lem:ladtail,lem:uniconst})
at the cascade-sector rate $(\sqrt3/2)^c$ and asymptotic gap constant
$1-(\sqrt3/2)^{1/2}\approx0.069$. The probabilistic side is closed---nominal-shell mass,
leading-coefficient structure, singleton rigidity, and the level-cost law
itself are all theorems
(\cref{lem:nomshell,lem:leadcoeff,op:onecond}), the last with the exact
constant $\lambda=\frac47<3^{-1/2}$ via the parity-trichotomy and
mod-$9$ orbit-spread argument. Small $c$ is covered by the
strict gap (\cref{cor:strict}) and finiteness---no certified numerics
needed for the qualitative statement.
The earlier conditional routes---comparison-graph expansion
(\cref{conj:exp}), short-witness coverage (\cref{conj:mult}), the
bounded-level holonomy audit, and the canonical witness
(\cref{thm:witness}) as a route to uniformity---are retained as
strictness, diagnostic, and conditional results, but are superseded as
routes to the uniform gap. The heavy-row audits support contraction by
mass-plus-profile bookkeeping alone (per-level decay $\approx0.72$; no
signed cancellation available or required), and the deeper-channel saving
is proved. To be precise about scope: it is the \emph{finite H23a operator
theorem} that is proved here---and the keystone arrow (bad orbit
$\Rightarrow$ finite non-coboundary spectral shadow) is now also proved
(\cref{thm:keystone}), with the coboundary classification supplying
exactly the exempt direction. The converse of the arrow fails (the
trivial cycle itself deviates, \cref{rem:onedir}), so the remaining
statement is the \emph{quenched shadow exclusion} (\cref{op:quenched}):
no non-eventually-trivial orbit may sustain such a deviation. That
exclusion implies the Collatz conjecture; it is essentially the
conjecture in minimal finite-shadow language, and it sits exactly at the
deterministic-equidistribution barrier, with three identified attack
routes of which ``shadow deviation forces descent'' is the most
realistic.

\paragraph{Lean formalization.} The elementary core of the repetition
rigidity argument is machine-verified in Lean~4 (Mathlib v4.27.0):
\texttt{CollatzLean/FiniteSpectralShadows.lean} proves
\texttt{word\_congr} (the $2$-adic coding lemma),
\texttt{repetition\_rigidity}, and \texttt{orbit\_periodic\_of\_eq}
sorry-free, depending only on Lean's three standard axioms
(\texttt{propext}, \texttt{Classical.choice}, \texttt{Quot.sound});
the Baker cycle bound is stated as the leading \texttt{proof\_wanted}
of the formal program. Full provenance for every computation and proof
step in this paper is recorded in the project's AIPM (\texttt{ailog})
trail.

\begin{thebibliography}{9}
\bibitem{companion} J.~A.~Janik, \emph{Finite Spectral Shadows for the Collatz
Valuation Cocycle: computational companion}, accompanying numerical report (2026).
\end{thebibliography}

\end{document}
